% Chapter 2, Section 3 _Linear Algebra_ Jim Hefferon
%  http://joshua.smcvt.edu/linearalgebra
%  2001-Jun-10
\section{Basis dan Dimensi}
\index{basis|(}
Bagian sebelumnya berakhir dengan pengamatan bahwa
suatu himpunan perentang bersifat minimal ketika bebas linear dan
suatu himpunan bebas linear bersifat maksimal ketika merentang ruang tersebut.
Jadi, gagasan himpunan perentang minimal dan himpunan bebas maksimal
berimpit.
Dalam bagian ini, kita akan menamai gagasan tersebut dan mempelajari
sifat-sifatnya.





\subsection{Basis}
\begin{definition} \label{def:basis}
%<*df:basis>
Suatu \definend{basis}\index{basis!definisi}\index{ruang vektor!basis}
bagi ruang vektor adalah barisan vektor yang bebas linear dan merentang ruang
tersebut.
%</df:basis>
\end{definition}

%<*BasisNotation>
Karena basis merupakan suatu barisan, yang berarti bahwa dua basis berbeda
jika memuat elemen-elemen yang sama tetapi dalam urutan yang berbeda,
% the 
% order of the elements is significant,
kita menuliskannya dengan tanda kurung sudut
\( \sequence{\vec{\beta}_1,\vec{\beta}_2,\ldots} \).\appendrefs{barisan}\@ %\spacefactor=1000{}
%</BasisNotation>
(Suatu barisan bebas linear jika multihimpunan yang terdiri atas elemen-elemen
barisan tersebut bebas linear.
Demikian pula, suatu barisan merentang ruang jika himpunan elemen-elemen
barisan tersebut merentang ruang.)

\begin{example}  \label{ex:FstBasisRTwo}
Berikut adalah suatu basis bagi \( \Re^2 \).
\begin{equation*}
  \sequence{ \colvec{2 \\ 4},\colvec{1 \\ 1} }
\end{equation*}
Basis tersebut bebas linear
\begin{equation*}
  c_1\colvec{2 \\ 4}+c_2\colvec{1 \\ 1}=\colvec{0 \\ 0}
  \quad\implies\quad
  \begin{linsys}{2}
     2c_1  &+  &1c_2  &=  &0  \\
     4c_1  &+  &1c_2  &=  &0  
   \end{linsys}
  \quad\implies\quad
  c_1=c_2=0
\end{equation*}
dan merentang \( \Re^2 \).
\begin{equation*}
  \begin{linsys}{2}
    2c_1  &+  &1c_2  &=  &x  \\
    4c_1  &+  &1c_2  &=  &y  
  \end{linsys}
  \quad\implies\quad
  c_2=2x-y\text{\ dan\ } c_1=(y-x)/2
\end{equation*}
\end{example}

\begin{example}
Basis bagi \( \Re^2 \) berikut berbeda dari basis sebelumnya
\begin{equation*}
  \sequence{\colvec[r]{1 \\ 1},\colvec[r]{2 \\ 4}}
\end{equation*}
karena urutannya berbeda.
Pemeriksaan bahwa barisan tersebut merupakan basis sama seperti dalam contoh
sebelumnya.
\end{example}

\begin{example}
Ruang \( \Re^2 \) memiliki banyak basis.
Berikut adalah basis lainnya.
\begin{equation*}
  \sequence{ \colvec[r]{1 \\ 0},\colvec[r]{0 \\ 1} }
\end{equation*}
Pemeriksaannya mudah.
\end{example}

\begin{definition} \label{df:StandardBasis}
%<*df:StandardBasis>
Untuk setiap \( \Re^n \),
\begin{equation*}
   \stdbasis_n=\sequence{
     \colvec[r]{1 \\ 0 \\ \vdotswithin{0} \\ 0},
     \colvec[r]{0 \\ 1 \\ \vdotswithin{0} \\ 0},
     \dots,\,
     \colvec[r]{0 \\ 0 \\ \vdotswithin{0} \\ 1}}
\end{equation*}
merupakan basis \definend{standar}\index{basis standar}%
\index{basis!standar} (atau basis \definend{alami}).
Kita menyatakan vektor-vektor ini dengan \( \vec{e}_1,\dots,\vec{e}_n \).
%</df:StandardBasis>
\end{definition}

\noindent 
Buku-buku kalkulus menyatakan vektor basis standar $\Re^2$ dengan
\( \vec{\imath} \) dan \( \vec{\jmath} \), bukan $\vec{e}_1$
dan $\vec{e}_2$, serta menyatakan vektor basis standar $\Re^3$ dengan
\( \vec{\imath} \), \( \vec{\jmath} \), dan \( \vec{k} \),
bukan $\vec{e}_1$, $\vec{e}_2$, dan $\vec{e}_3$.
Perhatikan bahwa \( \vec{e}_1 \) memiliki arti yang berbeda dalam pembahasan
\( \Re^3 \) dan dalam pembahasan \( \Re^2 \).

\begin{example}  \label{ex:BasisForCosPlusSin}
Perhatikan ruang
\( \set{a\cdot\cos\theta+b\cdot\sin\theta\suchthat a,b\in\Re} \)
dari fungsi-fungsi dengan variabel riil $\theta$.
Berikut adalah basis alami
$ \sequence{\cos\theta, \sin\theta}=\sequence{1\cdot\cos\theta+0\cdot\sin\theta,
             0\cdot\cos\theta+1\cdot\sin\theta}
    $.
Basis yang lebih umum bagi ruang ini adalah
\( \sequence{\cos\theta-\sin\theta,
             2\cos\theta+3\sin\theta} \).
Pemeriksaan bahwa keduanya merupakan basis terdapat dalam
\nearbyexercise{exer:VerifBasesCosPlusSin}.
\end{example}

\begin{example}
Basis alami bagi ruang vektor polinomial kubik \( \polyspace_3 \) adalah
\( \sequence{1,x,x^2,x^3} \).
Dua basis lain bagi ruang ini adalah \( \sequence{x^3,3x^2,6x,6} \) dan
\( \sequence{1,1+x,1+x+x^2,1+x+x^2+x^3} \).
Pemeriksaan bahwa masing-masing bebas linear dan merentang ruang tersebut
mudah dilakukan.
\end{example}

\begin{example}
Ruang trivial\index{ruang trivial}\index{ruang vektor!trivial}
$\set{\zero}$ hanya memiliki satu basis, yaitu basis kosong
\( \sequence{} \).
\end{example}

\begin{example}
Ruang polinomial berderajat hingga memiliki basis dengan tak hingga banyak
elemen,
\( \sequence{1,x,x^2,\ldots} \).
\end{example}

\begin{example}
Kita pernah menjumpai basis sebelumnya.
Dalam bab pertama, kita mendeskripsikan himpunan solusi sistem homogen seperti
sistem berikut
\begin{equation*}
   \begin{linsys}{4}
      x  &+  &y  &   &   &-   &w   &=  &0  \\
         &   &   &   &z  &+   &w   &=  &0  
   \end{linsys}
\end{equation*}
dengan melakukan parametrisasi.
\begin{equation*}
  \set{\colvec[r]{-1 \\ 1 \\ 0 \\ 0}y
       +\colvec[r]{1 \\ 0 \\ -1 \\ 1}w
       \suchthat y,w\in\Re }
\end{equation*}
Jadi, ruang vektor solusi tersebut merupakan rentang dari suatu himpunan
berelemen dua.
Himpunan dua vektor ini juga bebas linear, yang mudah diperiksa.
Oleh karena itu, himpunan solusi tersebut merupakan subruang dari \( \Re^4 \)
dengan basis yang terdiri atas kedua vektor ini.
\end{example}

\begin{example}
Parametrisasi juga menghasilkan basis bagi ruang vektor lain, bukan hanya bagi
himpunan solusi sistem homogen.
Untuk mencari basis bagi subruang dari $\matspace_{\nbyn{2}}$ berikut,
\begin{equation*}
  \set{\begin{mat}
         a  &b  \\
         c  &0
       \end{mat} \suchthat a+b-2c=0}
\end{equation*}
kita tulis ulang syaratnya sebagai $a=-b+2c$.
\begin{equation*}
  \set{\begin{mat}
         -b+2c  &b  \\
          c     &0
       \end{mat} \suchthat b,c \in \Re}
  =\set{b\begin{mat}[r]
         -1  &1  \\
          0  &0
       \end{mat}+
       c\begin{mat}[r]
          2  &0  \\
          1  &0
       \end{mat} \suchthat b,c \in \Re}
\end{equation*}
Jadi, berikut adalah calon basis yang wajar.
\begin{equation*}
  \sequence{\begin{mat}[r]
         -1  &1  \\
          0  &0
       \end{mat},
       \begin{mat}[r]
          2  &0  \\
          1  &0
       \end{mat} }
\end{equation*}
Perhitungan di atas menunjukkan bahwa barisan tersebut merentang ruangnya.
Kebebasan linearnya juga mudah diperiksa.
\end{example}

Perhatikan kembali \nearbyexample{ex:FstBasisRTwo}.
% To verify linearly independence we looked at 
% linear combinations of the set's members that total to
% the zero vector 
% $c_1\vec{\beta}_1+c_2\vec{\beta}_2=\binom{0}{0}$.
% The resulting 
% calculation shows that such a combination is unique,
% that $c_1$ must be $0$ and $c_2$ must be $0$.
Untuk memeriksa bahwa himpunan tersebut merentang ruang, kita meninjau
kombinasi linear yang menghasilkan suatu anggota ruang,
$c_1\vec{\beta}_1+c_2\vec{\beta}_2=\binom{x}{y}$.
Dalam contoh tersebut, kita hanya mencatat bahwa kombinasi semacam itu ada,
yakni bahwa untuk setiap $x,y$ terdapat $c_1,c_2$, tetapi perhitungannya
sebenarnya juga menunjukkan bahwa kombinasi itu unik:~$c_1$ harus sama dengan
$(y-x)/2$ dan $c_2$ harus sama dengan $2x-y$.
 
\begin{theorem} \label{th:BasisIffUniqueRepWRT}
%<*th:BasisIffUniqueRepWRT>
Dalam sebarang ruang vektor, suatu subhimpunan merupakan basis\index{basis}
jika dan hanya jika setiap vektor dalam ruang dapat dinyatakan sebagai
kombinasi linear elemen-elemen subhimpunan tersebut dengan tepat satu cara.
%</th:BasisIffUniqueRepWRT>
\end{theorem}

\noindent Kita menganggap dua kombinasi linear sama jika keduanya memiliki
suku-suku yang sama dalam urutan berbeda, atau jika keduanya hanya berbeda
karena penambahan atau penghapusan suku berbentuk
`\( 0\cdot\vec{\beta} \)'.


\begin{proof}
%<*pf:BasisIffUniqueRepWRT0>
Suatu barisan merupakan basis jika dan hanya jika vektor-vektornya membentuk
himpunan yang merentang dan bebas linear.
Suatu subhimpunan merupakan himpunan perentang jika dan hanya jika setiap
vektor dalam ruang merupakan kombinasi linear elemen-elemen subhimpunan
tersebut dengan sedikitnya satu cara.
Jadi, kita hanya perlu menunjukkan bahwa suatu subhimpunan perentang bebas
linear jika dan hanya jika setiap vektor dalam ruang merupakan kombinasi
linear elemen-elemen subhimpunan tersebut dengan paling banyak satu cara.
%</pf:BasisIffUniqueRepWRT0>

%<*pf:BasisIffUniqueRepWRT1>
Perhatikan dua pernyataan suatu vektor sebagai kombinasi linear anggota-anggota
subhimpunan tersebut.
Susun ulang kedua jumlah dan, jika perlu, tambahkan beberapa suku
\( 0\cdot\vec{\beta}_i \), sehingga kedua jumlah mengombinasikan
\( \vec{\beta} \) yang sama dalam urutan yang sama:
\( \vec{v}=\lincombo{c}{\vec{\beta}} \) dan
\( \vec{v}=\lincombo{d}{\vec{\beta}} \).
Sekarang,
\begin{equation*}
   \lincombo{c}{\vec{\beta}}=\lincombo{d}{\vec{\beta}}
\end{equation*}
berlaku jika dan hanya jika
\begin{equation*}
   (c_1-d_1)\vec{\beta}_1+\dots+(c_n-d_n)\vec{\beta}_n=\zero
\end{equation*}
berlaku.
Jadi, pernyataan bahwa setiap koefisien dalam persamaan terakhir sama dengan
nol ekuivalen dengan pernyataan bahwa \( c_i=d_i \) untuk setiap \( i \),
yakni bahwa setiap vektor dapat dinyatakan secara unik sebagai kombinasi linear
vektor-vektor \( \vec{\beta} \).
%</pf:BasisIffUniqueRepWRT1>
\end{proof}

\begin{definition} \label{def:RepresentingVectors}
%<*df:RepresentingVectors>
Dalam ruang vektor dengan basis $B$, \definend{representasi \( \vec{v} \)
terhadap \( B \)}%
\index{representasi!vektor}\index{vektor!representasi} adalah vektor kolom
koefisien-koefisien yang digunakan untuk menyatakan $\vec{v}$ sebagai
kombinasi linear vektor-vektor basis:
\begin{equation*}
  \rep{\vec{v}}{B}
  =
  \colvec{c_1 \\ c_2 \\ \vdots \\ c_n}_{B}
\end{equation*}
di mana
\( B=\sequence{\vec{\beta}_1,\dots,\vec{\beta}_n} \) dan 
\( \vec{v}=\lincombo{c}{\vec{\beta}} \).
Bilangan-bilangan \( c \) adalah
\definend{koordinat \( \vec{v} \) terhadap \( B \)}%
\index{koordinat!terhadap suatu basis}.
%</df:RepresentingVectors>
\end{definition}

\begin{example}  \label{ex:RepPolyWRTA Basis}
Dalam \( \polyspace_3 \), terhadap basis
\( B=\sequence{1,2x,2x^2,2x^3} \),
representasi \( x+x^2 \) adalah
\begin{equation*}
  \rep{x+x^2}{B}=\colvec[r]{0 \\ 1/2 \\ 1/2 \\ 0}_B
\end{equation*}
karena $x+x^2=0\cdot 1+(1/2)\cdot 2x+(1/2)\cdot 2x^2+0\cdot 2x^3$.
Terhadap basis lain \( D=\sequence{1+x,1-x,x+x^2,x+x^3} \),
representasinya berbeda.
\begin{equation*}
  \rep{x+x^2}{D}=\colvec[r]{0 \\ 0 \\ 1 \\ 0}_D
\end{equation*}
\end{example}

\begin{remark}
\nearbydefinition{def:basis} mensyaratkan basis berupa barisan agar kita dapat
menuliskan koordinat-koordinat ini dalam suatu urutan.
\end{remark}

Jika hanya ada satu basis dalam pembahasan, kita sering menghilangkan subskrip
yang menamai basis tersebut.

\begin{example}
Dalam \( \Re^2 \), untuk mencari koordinat vektor
$\vec{v}=\binom{3}{2}$ terhadap basis
\begin{equation*}
  B=\sequence{
              \colvec[r]{1 \\ 1},
              \colvec[r]{0 \\ 2} }
\end{equation*}
selesaikan
\begin{equation*}
  c_1\colvec[r]{1 \\ 1}
  +c_2\colvec[r]{0 \\ 2}
  =
  \colvec[r]{3 \\ 2}
\end{equation*}
dan peroleh $c_1=3$ serta $c_2=-1/2$.
\begin{equation*}
  \rep{\vec{v}}{B}=\colvec[r]{3 \\ -1/2}
\end{equation*}  
\end{example}

Penulisan representasi sebagai kolom menggeneralisasi kasus yang sudah dikenal:
dalam \( \Re^n \) dan terhadap basis standar \( \stdbasis_n \), vektor yang
berpangkal di titik asal dan berujung di \( (v_1,\dots,v_n) \) memiliki
representasi berikut.
\begin{equation*}
  \rep{\colvec{v_1 \\ \vdots \\ v_n}}{\stdbasis_n}
    =
  \colvec{v_1 \\ \vdots \\ v_n}_{\stdbasis_n}
\end{equation*}
Berikut adalah sebuah contoh.
\begin{equation*}
  \rep{\colvec{-1 \\ 1}}{\stdbasis_{n}}
  =\colvec{-1 \\ 1}
\end{equation*}

\begin{remark}  \label{rem:RepNotation}
Notasi $\rep{\vec{v}}{B}$ bukan notasi baku.
Notasi yang paling lazim adalah $[\vec{v}]_B$, tetapi salah satu kelebihan
$\rep{\vec{v}}{B}$ ialah notasi ini lebih sulit disalahtafsirkan atau terlewat.
\end{remark}

Kolom tersebut merepresentasikan vektor dalam arti bahwa suatu relasi linear
berlaku di antara suatu himpunan vektor jika dan hanya jika relasi tersebut
berlaku di antara himpunan representasinya.

\begin{lemma}  \label{lem:LinearRelRepresented}
%<*lm:LinearRelRepresented>
Jika $B$ adalah basis dengan $n$~elemen, maka untuk sebarang himpunan vektor,
$a_1\vec{v}_1+\cdots+a_k\vec{v}_k=\vec{0}_{V}$ jika dan hanya jika
$a_1\rep{\vec{v}_1}{B}+\cdots+a_k\rep{\vec{v}_k}{B}=\vec{0}_{\Re^n}$.
%</lm:LinearRelRepresented>
\end{lemma}

\begin{proof}
Tetapkan basis \( B=\sequence{\vec{\beta}_1,\dots,\vec{\beta}_n} \)
dan misalkan
\begin{equation*}
  \rep{\vec{v}_1}{B}=\colvec{c_{1,1} \\ \vdots \\ c_{n,1}}
  \quad\ldots\quad
  \rep{\vec{v}_k}{B}=\colvec{c_{1,k} \\ \vdots \\ c_{n,k}}
\end{equation*}
sehingga $\vec{v}_1=c_{1,1}\vec{\beta}_1+\dots+c_{n,1}\vec{\beta}_n$, dan
seterusnya.
Maka $a_1\vec{v}_1+\dots+a_k\vec{v}_k=\vec{0}$ ekuivalen dengan persamaan
berikut.
\begin{align*}
  \vec{0}
  &=a_1\cdot(c_{1,1}\vec{\beta}_1+\dots+c_{n,1}\vec{\beta}_n)
    +\dots+
    a_k\cdot(c_{1,k}\vec{\beta}_1+\dots+c_{n,k}\vec{\beta}_n)  \\
  &=(a_1c_{1,1}+\dots+a_kc_{1,k})\cdot\vec{\beta}_1
    +\dots+
    (a_1c_{n,1}+\dots+a_kc_{n,k})\cdot\vec{\beta}_n 
\end{align*}
Jelas bahwa persamaan terakhir benar jika koefisien-koefisiennya nol.
Namun, karena \( B \) merupakan basis,
\nearbytheorem{th:BasisIffUniqueRepWRT} menyatakan bahwa persamaan terakhir
benar jika dan hanya jika koefisien-koefisiennya nol.
Jadi, relasi tersebut ekuivalen dengan sistem berikut.
\begin{align*}
    a_1c_{1,1}+\dots+a_kc_{1,k} &=0    \\
                               &\alignedvdots \\ 
    a_1c_{n,1}+\dots+a_kc_{n,k} &=0
\end{align*}
Berikut adalah bentuk ekuivalennya dalam vektor kolom.
\begin{equation*}
  a_1\colvec{c_{1,1} \\ \vdots \\ c_{n,1}}
  +\dots+
  a_k\colvec{c_{1,k} \\ \vdots \\ c_{n,k}}
  =\colvec{0 \\ \vdots \\ 0}
\end{equation*}
Perhatikan bahwa bukan hanya suatu relasi berlaku bagi satu himpunan jika dan
hanya jika relasi itu berlaku bagi himpunan lainnya, melainkan relasinya juga
sama\Dash nilai-nilai \( a_i \) tetap sama.
\end{proof}

\begin{example}
\nearbyexample{ex:RepPolyWRTA Basis} mencari representasi
\( x+x^2\in\polyspace_3 \) terhadap
\( B=\sequence{1,2x,2x^2,2x^3} \).
\begin{equation*}
  \rep{x+x^2}{B}=\colvec[r]{0 \\ 1/2 \\ 1/2 \\ 0}_B
\end{equation*}
Relasi berikut
\begin{equation*}
  2\cdot(x+x^2)-1\cdot(2x)-2\cdot(x^2) = 0+0x+0x^2+0x^3
\end{equation*}
direpresentasikan oleh relasi berikut.
\begin{equation*}
  2\cdot\rep{x+x^2}{B}-\rep{2x}{B}-2\cdot\rep{x^2}{B}
  =2\cdot\colvec{0 \\ 1/2 \\ 1/2 \\ 0}
   -\colvec{0 \\ 1 \\ 0 \\ 0}
   -2\cdot\colvec{0 \\ 0 \\ 1/2 \\ 0}
  =\colvec{0 \\ 0 \\ 0 \\ 0}
\end{equation*}

\end{example}

Penggunaan utama representasi akan muncul kemudian, tetapi definisinya
diberikan di sini karena fakta bahwa setiap vektor merupakan kombinasi linear
vektor-vektor basis secara unik adalah sifat penting suatu basis, dan juga
untuk membantu menegaskan satu hal.
Untuk perhitungan koordinat dan keperluan lain, kita akan membatasi perhatian
pada ruang-ruang yang memiliki basis dengan hanya berhingga banyak elemen.
Pembahasannya dimulai dalam subbagian berikutnya.

\begin{exercises}
  \recommended \item Tentukan apakah masing-masing berikut merupakan basis
    bagi $\polyspace_2$.
    \begin{exparts*}
      \partsitem $\sequence{x^2-x+1, 2x+1, 2x-1}$
      \partsitem $\sequence{x+x^2, x-x^2}$
    \end{exparts*}
    \begin{answer}
      \begin{exparts}
        \partsitem Ini merupakan basis bagi $\polyspace_2$.
           Untuk menunjukkan bahwa himpunan tersebut merentang ruang, kita
           meninjau suatu $a_2x^2+a_1x+a_0\in\polyspace_2$ yang umum dan mencari
           skalar $c_1, c_2, c_3\in\Re$ sedemikian sehingga
           $a_2x^2+a_1x+a_0=c_1\cdot(x^2-x+1) +c_2\cdot(2x+1) +c_3(2x-1)$.
           Sistem linear yang dihasilkan adalah sebagai berikut.
           \begin{equation*}
             \begin{linsys}{3}
               c_1  &\spaceforemptycolumn  &     &  &     &=  &a_2 \\
                    &  &2c_2 &+ &2c_3 &=  &a_1 \\
                    &  &c_2  &- &c_3   &=  &a_0
             \end{linsys}
             \grstep{(-1/2)\rho_2+\rho_3}
             \begin{linsys}{3}
               c_1  &\spaceforemptycolumn  &     &  &     &=  &a_2\hfill\hbox{} \\
                    &  &2c_2 &+ &2c_3    &=  &a_1\hfill\hbox{} \\
                    &  &     &  &-2c_3   &=  &a_0-(1/2)a_1
             \end{linsys}
           \end{equation*}
           Jadi, untuk sebarang rangkap tiga nilai $a_i$, kita dapat menghitung
           nilai-nilai $c_j$ sebagai $c_1=a_2$, $c_2=(1/4)a_1+(1/2)a_0$, dan
           $c_3=(1/4)a_1-(1/2)a_0$.
           Oleh karena itu, setiap elemen $\polyspace_2$ merupakan kombinasi
           dari $x^2-x+1$, $2x+1$, dan~$2x-1$ yang diberikan.

           Untuk membuktikan bahwa himpunan ketiga polinomial tersebut bebas
           linear, kita dapat menyusun persamaan
           $0x^2+0x+0=c_1\cdot(x^2-x+1) +c_2\cdot(2x+1) +c_3(2x-1)$
           lalu menyelesaikannya; hasilnya adalah $c_1=0$, $c_2=0$, dan
           $c_3=0$.
           Sebagai alternatif, kita dapat mengamati bahwa solusi dalam paragraf
           sebelumnya unik, lalu menggunakan
           \nearbytheorem{th:BasisIffUniqueRepWRT}.
        \partsitem Ini bukan basis.
           Himpunan tersebut tidak merentang ruang karena tidak ada kombinasi
           $c_1\cdot (x+x^2) +c_2\cdot (x-x^2)$ dari kedua polinomial yang
           menghasilkan polinomial $3\in\polyspace_2$.
      \end{exparts}
    \end{answer}
  \recommended \item 
    Tentukan apakah masing-masing berikut merupakan basis bagi \( \Re^3 \).
    \begin{exparts*}
      \partsitem \( \sequence{
                 \colvec[r]{1 \\ 2 \\ 3},
                 \colvec[r]{3 \\ 2 \\ 1},
                 \colvec[r]{0 \\ 0 \\ 1}}  \)
      \partsitem \( \sequence{
                 \colvec[r]{1 \\ 2 \\ 3},
                 \colvec[r]{3 \\ 2 \\ 1}}  \)
      \partsitem \( \sequence{
                 \colvec[r]{0 \\ 2 \\ -1},
                 \colvec[r]{1 \\ 1 \\ 1},
                 \colvec[r]{2 \\ 5 \\ 0}}  \)
      \partsitem \( \sequence{
                 \colvec[r]{0 \\ 2 \\ -1},
                 \colvec[r]{1 \\ 1 \\ 1},
                 \colvec[r]{1 \\ 3 \\ 0}}  \)
    \end{exparts*}
    \begin{answer}
      Menurut \nearbytheorem{th:BasisIffUniqueRepWRT}, masing-masing merupakan
      basis jika dan hanya jika setiap vektor dalam ruang dapat kita nyatakan
      secara unik sebagai kombinasi linear vektor-vektor yang diberikan.
      \begin{exparts}
        \partsitem Ya, ini merupakan basis.
          Relasi
          \begin{equation*}
            c_1\colvec[r]{1 \\ 2 \\ 3}
            +c_2\colvec[r]{3 \\ 2 \\ 1}
            +c_3\colvec[r]{0 \\ 0 \\ 1}
            =\colvec{x \\ y \\ z}
          \end{equation*}
          menghasilkan
          \begin{equation*}
            \begin{amat}{3}
              1  &3  &0  &x  \\
              2  &2  &0  &y  \\
              3  &1  &1  &z
            \end{amat}
            \grstep[-3\rho_1+\rho_3]{-2\rho_1+\rho_2}
            \repeatedgrstep{-2\rho_2+\rho_3}
            \begin{amat}{3}
              1  &3  &0  &x \hfill\hbox{}  \\
              0  &-4 &0  &-2x+y \hfill\hbox{}  \\
              0  &0  &1  &x-2y+z 
            \end{amat}
          \end{equation*}
          yang memiliki solusi unik
          \( c_3=x-2y+z \), \( c_2=x/2-y/4 \), dan
          \( c_1=-x/2+3y/4 \).
        \partsitem Ini bukan basis.
          Penyusunan relasi seperti dalam butir sebelumnya,
          \begin{equation*}
            c_1\colvec[r]{1 \\ 2 \\ 3}
            +c_2\colvec[r]{3 \\ 2 \\ 1}
            =\colvec{x \\ y \\ z}
          \end{equation*}
          menghasilkan sistem linear dengan reduksi
          \begin{equation*}
            \begin{amat}{2}
              1  &3  &x  \\ 
              2  &2  &y  \\
              3  &1  &z
            \end{amat}
            \grstep[-3\rho_1+\rho_3]{-2\rho_1+\rho_2}
            \repeatedgrstep{-2\rho_2+\rho_3}
            \begin{amat}{2}
              1  &3  &x\hfill\hbox{}  \\ 
              0  &-4 &-2x+y\hfill\hbox{}  \\
              0  &0  &x-2y+z
            \end{amat}
          \end{equation*}
          Solusi ada jika dan hanya jika komponen-komponen $x$, $y$, dan $z$
          dari vektor bertinggi tiga memenuhi $x-2y+z=0$.
          Sebagai contoh, kita dapat mencari koefisien $c_1$ dan $c_2$ yang
          sesuai ketika $x=1$, $y=1$, dan $z=1$.
          Namun, tidak ada nilai-nilai $c$ yang sesuai untuk
          $x=1$, $y=1$, dan $z=2$.
          Jadi, ini bukan basis; himpunan tersebut tidak merentang ruang.
        \partsitem Ya, ini merupakan basis.
          Penyusunan relasinya menghasilkan reduksi berikut
          \begin{equation*}
            \begin{amat}{3}
              0  &1  &2  &x \hfill\hbox{} \\
              2  &1  &5  &y \hfill\hbox{} \\
             -1  &1  &0  &z 
            \end{amat}
            \grstep{\rho_1\swap\rho_3}
            \repeatedgrstep{2\rho_1+\rho_2}
            \repeatedgrstep{-(1/3)\rho_2+\rho_3}
            \begin{amat}{3}
             -1  &1  &0   &z  \hfill\hbox{} \\
              0  &3  &5   &y+2z \hfill\hbox{} \\
              0  &0  &1/3 &x-y/3-2z/3
            \end{amat}
          \end{equation*}
          yang memiliki solusi unik untuk setiap rangkap tiga komponen
          $x$, $y$, dan $z$.
        \partsitem Tidak, ini bukan basis.
          Reduksi
          \begin{equation*}
            \begin{amat}{3}
              0  &1  &1  &x   \\
              2  &1  &3  &y   \\
             -1  &1  &0  &z  
            \end{amat}
            \grstep{\rho_1\swap\rho_3}
            \repeatedgrstep{2\rho_1+\rho_2}
            \repeatedgrstep{(-1/3)\rho_2+\rho_3}
            \begin{amat}{3}
             -1  &1  &0  &z  \hfill\hbox{} \\
              0  &3  &3  &y+2z  \hfill\hbox{} \\
              0  &0  &0  &x-y/3-2z/3
            \end{amat}
          \end{equation*}
          tidak memiliki solusi untuk setiap rangkap tiga $x$, $y$, dan $z$.
          Rentang himpunan yang diberikan hanya memuat vektor bertinggi tiga
          yang memenuhi \( x=y/3+2z/3 \).
      \end{exparts}  
     \end{answer}
  \recommended \item 
    Nyatakan representasi vektor terhadap basis yang diberikan.
    \begin{exparts}
      \partsitem \( \colvec[r]{1 \\ 2} \),
        \( B=\sequence{\colvec[r]{1 \\ 1},\colvec[r]{-1 \\ 1}}\subseteq\Re^2 \)
      \partsitem \( x^2+x^3 \),
        \( D=\sequence{1,1+x,1+x+x^2,1+x+x^2+x^3}\subseteq\polyspace_3 \)
      \partsitem \( \colvec[r]{0 \\ -1 \\ 0 \\ 1} \),
        \( \stdbasis_4\subseteq\Re^4 \)
    \end{exparts}
    \begin{answer}  
      \begin{exparts}
        \partsitem Kita menyelesaikan
          \begin{equation*}
            c_1\colvec[r]{1 \\ 1}
            +c_2\colvec[r]{-1 \\ 1}
            =\colvec[r]{1 \\ 2}
          \end{equation*}
          dengan reduksi
          \begin{equation*}
             \begin{amat}[r]{2}
               1  &-1  &1  \\
               1  &1   &2
             \end{amat}
             \grstep{-\rho_1+\rho_2}
             \begin{amat}[r]{2}
               1  &-1  &1  \\
               0  &2   &1
             \end{amat}
          \end{equation*}
          lalu menyimpulkan bahwa \( c_2=1/2 \), sehingga \( c_1=3/2 \).
          Jadi, representasinya adalah sebagai berikut.
          \begin{equation*}
            \rep{\colvec[r]{1 \\ 2}}{B}=\colvec[r]{3/2 \\ 1/2}_B
          \end{equation*}
        \partsitem Relasi
           $c_1\cdot(1)+c_2\cdot(1+x)+c_3\cdot(1+x+x^2)+c_4\cdot(1+x+x^2+x^3)
             =x^2+x^3$
           mudah diselesaikan secara langsung dan menghasilkan $c_4=1$,
           $c_3=0$, $c_2=-1$, dan $c_1=0$.
           \begin{equation*}
              \rep{x^2+x^3}{D}=\colvec[r]{0 \\ -1 \\ 0 \\ 1}_D
           \end{equation*} 
        \partsitem \( \rep{\colvec[r]{0 \\ -1 \\ 0 \\ 1}}{\stdbasis_4}
                     =\colvec[r]{0 \\ -1 \\ 0 \\ 1}_{\stdbasis_4} \)
      \end{exparts}  
    \end{answer}
  \item  Nyatakan representasi vektor terhadap masing-masing dari kedua basis.
    \begin{equation*}
      \vec{v}=\colvec{3  \\ -1}
      \quad
      B_1=\sequence{\colvec{1  \\ -1}, \colvec{1  \\ 1}},\;
      B_2=\sequence{\colvec{1 \\ 2}, \colvec{1 \\ 3}}
    \end{equation*}
    \begin{answer}
       Penyelesaian
      \begin{equation*}
        \colvec{3 \\ -1}=\colvec{1 \\ -1}\cdot c_1
                         +\colvec{1 \\ 1}\cdot c_2
      \end{equation*}
      menghasilkan $c_1=2$ dan~$c_2=1$.
      \begin{equation*}
        \rep{\colvec{3 \\ -1}}{B_1}
           =\colvec{2 \\ 1}_{B_1}
      \end{equation*}
      Demikian pula, penyelesaian
      \begin{equation*}
        \colvec{3 \\ -1}=\colvec{1 \\ 2}\cdot c_1
                         +\colvec{1 \\ 3}\cdot c_2
      \end{equation*}
      menghasilkan representasi berikut.
      \begin{equation*}
        \rep{\colvec{3 \\ -1}}{B_2}
           =\colvec{10 \\ -7}_{B_2}
      \end{equation*}     
    \end{answer}
  \item  
    Carilah suatu basis bagi \(  \polyspace_2 \), yaitu ruang semua polinomial
    kuadratik.
    Apakah setiap basis semacam itu harus memuat polinomial dari setiap
    derajat:~derajat nol, derajat satu, dan derajat dua?
    \begin{answer}
      Basis alaminya adalah \( \sequence{1,x,x^2} \).
      Ada basis bagi $\polyspace_2$ yang tidak memuat polinomial berderajat
      satu maupun berderajat nol.
      Salah satunya adalah \( \sequence{1+x+x^2,x+x^2,x^2} \).
      (Namun, setiap basis memiliki sedikitnya satu polinomial berderajat dua.)
    \end{answer}
  \item
    Carilah suatu basis bagi himpunan solusi sistem berikut.
    \begin{equation*}
      \begin{linsys}{4}
         x_1  &-  &4x_2  &+  &3x_3  &-  &x_4  &=  &0  \\
        2x_1  &-  &8x_2  &+  &6x_3  &-  &2x_4 &=  &0  
      \end{linsys}
    \end{equation*}
    \begin{answer}
      Reduksi
      \begin{equation*}
        \begin{amat}[r]{4}
          1  &-4  &3  &-1  &0  \\
          2  &-8  &6  &-2  &0
        \end{amat}
        \grstep{-2\rho_1+\rho_2}
        \begin{amat}[r]{4}
          1  &-4  &3  &-1  &0  \\
          0  &0   &0  &0   &0
        \end{amat}
      \end{equation*}
      menunjukkan bahwa satu-satunya syarat adalah
      $x_1=4x_2-3x_3+x_4$.
      Himpunan solusinya adalah
      \begin{multline*}
        \set{\colvec{4x_2-3x_3+x_4 \\ x_2 \\ x_3 \\ x_4}
               \suchthat x_2,x_3,x_4\in\Re}                          \\
        =\set{x_2\colvec[r]{4 \\ 1 \\ 0 \\ 0}
             +x_3\colvec[r]{-3 \\ 0 \\ 1 \\ 0}
             +x_4\colvec[r]{1 \\ 0 \\ 0 \\ 1} \suchthat x_2,x_3,x_4\in\Re}
      \end{multline*}
      sehingga calon basis yang langsung terlihat adalah sebagai berikut.
      \begin{equation*}
       \sequence{\colvec[r]{4 \\ 1 \\ 0 \\ 0},
                   \colvec[r]{-3 \\ 0 \\ 1 \\ 0},
                   \colvec[r]{1 \\ 0 \\ 0 \\ 1}  }
      \end{equation*}  
      Kita telah menunjukkan bahwa himpunan ini merentang ruang, dan
      pemeriksaan bahwa himpunan tersebut juga bebas linear bersifat rutin.
    \end{answer}
  \recommended \item
    Carilah suatu basis bagi \( \matspace_{\nbyn{2}} \), yaitu ruang matriks
    \( \nbyn{2} \).
    \begin{answer}
      Ada banyak basis.
      Berikut adalah salah satu basis alami.
      \begin{equation*}
        \sequence{
           \begin{mat}[r]
             1  &0  \\
             0  &0
           \end{mat},
           \begin{mat}[r]
             0  &1  \\
             0  &0
           \end{mat},
           \begin{mat}[r]
             0  &0  \\
             1  &0
           \end{mat},
           \begin{mat}[r]
             0  &0  \\
             0  &1
           \end{mat}  }
      \end{equation*} 
    \end{answer}
  \recommended \item 
      Carilah suatu basis bagi masing-masing ruang berikut.
      \begin{exparts}
        \partsitem Subruang
          $\set{a_2x^2+a_1x+a_0\suchthat a_2-2a_1=a_0}$ dari $\polyspace_2$
        \partsitem Ruang vektor baris berlebar tiga yang jumlah komponen
          pertama dan keduanya sama dengan nol
        \partsitem Subruang matriks $\nbyn{2}$ berikut
          \begin{equation*}
            \set{\begin{mat}
                   a  &b  \\
                   0  &c  
                 \end{mat} \suchthat c-2b=0}
          \end{equation*}
      \end{exparts}
      \begin{answer}
        Untuk setiap butir, ada banyak jawaban yang mungkin.
        \begin{exparts}
          \partsitem Salah satu caranya ialah membuat parametrisasi dengan
            menyatakan $a_2$ sebagai kombinasi kedua koefisien lainnya,
            $a_2=2a_1+a_0$.
            Maka $a_2x^2+a_1x+a_0$ menjadi
            $(2a_1+a_0)x^2+a_1x+a_0$, dan
            \begin{multline*} 
              \set{(2a_1+a_0)x^2+a_1x+a_0\suchthat a_1,a_0\in\Re}          \\
              =\set{a_1\cdot(2x^2+x)+a_0\cdot(x^2+1)\suchthat a_1,a_0\in\Re}
            \end{multline*}
            menyarankan basis
            $\sequence{2x^2+x,x^2+1}$.
            Perhitungan ini baru menunjukkan bahwa basis tersebut merentang,
            tetapi pemeriksaan bahwa basis itu bebas linear bersifat rutin.
          \partsitem Parametrisasikan
            $\set{\rowvec{a  &b  &c}\suchthat a+b=0}$ untuk memperoleh
            $\set{\rowvec{-b &b &c}\suchthat b,c\in\Re}$, yang menyarankan
            penggunaan barisan
            $\sequence{\rowvec{-1 &1 &0},\rowvec{0 &0 &1}}$.
            Kita telah menunjukkan bahwa barisan tersebut merentang, dan
            pemeriksaan bahwa barisan itu bebas linear mudah dilakukan.
          \partsitem Penulisan ulang
            \begin{equation*}
              \set{\begin{mat}
                     a  &b  \\
                     0  &2b
                   \end{mat}\suchthat a,b\in\Re}
              =\set{a\cdot\begin{mat}[r]
                     1  &0  \\
                     0  &0 
                   \end{mat}
                  +b\cdot\begin{mat}[r]
                           0  &1  \\
                           0  &2
                          \end{mat} \suchthat a,b\in\Re}
            \end{equation*}
            menyarankan basis berikut.
            \begin{equation*}
              \sequence{\begin{mat}[r]
                          1  &0  \\
                          0  &0
                        \end{mat},
                        \begin{mat}[r]
                          0  &1  \\
                          0  &2
                        \end{mat}}
            \end{equation*}
         \end{exparts}  
      \end{answer}
  \item Carilah suatu basis bagi setiap ruang, lalu periksalah bahwa basis
    tersebut memang suatu basis.
    \begin{exparts}
      \partsitem
        Subruang $M=\set{a+bx+cx^2+dx^3\suchthat a-2b+c-d=0}$ dari
        $\polyspace_3$.
      \partsitem  Subruang dari $\matspace_{\nbyn{2}}$ berikut.
         \begin{equation*}
           W=\set{
             \begin{mat}
               a  &b  \\
               c  &d
             \end{mat}
             \suchthat a-c=0}
         \end{equation*}

    \end{exparts}
    \begin{answer}
      \begin{exparts}
        \partsitem
          Parametrisasikan $a-2b+c-d=0$ sebagai $a=2b-c+d$, $b=b$, $c=c$,
          dan~$d=d$ untuk memperoleh deskripsi $M$ berikut sebagai rentang
          suatu himpunan tiga vektor.
          \begin{equation*}
            M=\set{
             (2+x)\cdot b+(-1+x^2)\cdot c+(1+x^3)\cdot d
             \suchthat b,c,d\in\Re
             }
          \end{equation*}
          Untuk menunjukkan bahwa himpunan tiga vektor ini merupakan basis,
          tinggal diperiksa bahwa himpunan tersebut bebas linear.
          \begin{equation*}
             0+0x+0x^2+0x^3=(2+x)\cdot c_1+(-1+x^2)\cdot c_2+(1+x^3)\cdot c_3
          \end{equation*}
          Dari suku-suku $x$, kita memperoleh $c_1=0$.
          Dari suku-suku $x^2$, kita memperoleh $c_2=0$.
          Suku-suku $x^3$ menghasilkan $c_3=0$.
       \partsitem
         Pertama, parametrisasikan deskripsinya
         (perhatikan bahwa tidak disebutkannya $b$ dan~$d$ dalam deskripsi
         $W$ bukan berarti keduanya nol atau tidak ada, melainkan bahwa
         keduanya tidak dibatasi).
         \begin{equation*}
           W=\set{
             \begin{mat}
               0 &1 \\
               0 &0
             \end{mat}\cdot b
             +
             \begin{mat}
               1 &0 \\
               1 &0
             \end{mat}\cdot c
             +
             \begin{mat}
               0 &0 \\
               0 &1
             \end{mat}\cdot d
             \suchthat b,c,d\in\Re}
         \end{equation*}
         Hal ini memberikan $W$ sebagai rentang suatu himpunan berelemen tiga.
         Pembuktian selesai setelah kita menunjukkan bahwa himpunan tersebut
         bebas linear.
         \begin{equation*}
             \begin{mat}
               0 &0 \\
               0 &0
             \end{mat}
             =
             \begin{mat}
               0 &1 \\
               0 &0
             \end{mat}\cdot c_1
             +
             \begin{mat}
               1 &0 \\
               1 &0
             \end{mat}\cdot c_2
             +
             \begin{mat}
               0 &0 \\
               0 &1
             \end{mat}\cdot c_3   
         \end{equation*}
         Dengan menggunakan entri-entri kanan atas, kita memperoleh $c_1=0$.
         Entri-entri kiri atas menghasilkan $c_2=0$, dan entri-entri kiri bawah
         menunjukkan bahwa $c_3=0$.
     \end{exparts}
   \end{answer}
  \item \label{exer:VerifBasesCosPlusSin} 
    Periksalah \nearbyexample{ex:BasisForCosPlusSin}.
    \begin{answer}
      Kita akan menunjukkan bahwa barisan kedua merupakan basis; pemeriksaan
      untuk barisan pertama serupa.
      Kita akan menunjukkannya langsung dari definisi basis karena contoh ini
      muncul sebelum
      \nearbytheorem{th:BasisIffUniqueRepWRT}.

      Untuk melihat bahwa barisan tersebut bebas linear, kita susun
      \( c_1\cdot(\cos\theta-\sin\theta)+c_2\cdot(2\cos\theta+3\sin\theta)=
        0\cos\theta+0\sin\theta \).
      Dengan mengambil \( \theta=0 \) dan \( \theta=\pi/2 \), diperoleh sistem
      berikut
      \begin{equation*}
        \begin{linsys}{2}
          c_1\cdot 1    &+  &c_2\cdot 2   &=  &0  \\
          c_1\cdot (-1) &+  &c_2\cdot 3   &=  &0  
        \end{linsys}
        \grstep{\rho_1+\rho_2}
        \begin{linsys}{2}
          c_1  &+  &2c_2   &=  &0  \\
               &+  &5c_2   &=  &0 
        \end{linsys}
      \end{equation*}
      yang menunjukkan bahwa $c_1=0$ dan $c_2=0$.

      Perhitungan untuk rentangnya juga mudah; untuk sebarang $x,y\in\Re$,
      relasi
      \( c_1\cdot(\cos\theta-\sin\theta)+c_2\cdot(2\cos\theta+3\sin\theta)=
        x\cos\theta+y\sin\theta \)
      menghasilkan \( c_2=x/5+y/5 \) dan \( c_1=3x/5-2y/5 \), sehingga
      rentangnya adalah seluruh ruang.
    \end{answer}
  \recommended \item 
    Carilah rentang setiap himpunan, lalu carilah suatu basis bagi rentang
    tersebut.
    \begin{exparts*}
       \partsitem $\set{1+x,1+2x}$ di $\polyspace_2$
       \partsitem $\set{2-2x,3+4x^2}$ di $\polyspace_2$
    \end{exparts*}
    \begin{answer}
      \begin{exparts}
         \partsitem Pertanyaan tentang $a_0+a_1x+a_2x^2$ mana yang dapat
           dinyatakan sebagai $c_1\cdot (1+x)+c_2\cdot (1+2x)$ menghasilkan
           tiga persamaan linear yang menjelaskan koefisien $x^2$, koefisien
           $x$, dan suku konstanta.
           \begin{equation*} 
             \begin{linsys}{2}
               c_1 &+ &c_2  &= &a_0 \\
               c_1 &+ &2c_2 &= &a_1 \\
                   &  &0    &= &a_2
             \end{linsys}
           \end{equation*}
           Metode Gauss dengan substitusi balik menunjukkan, asalkan $a_2=0$,
           bahwa $c_2=-a_0+a_1$ dan $c_1=2a_0-a_1$.
           Jadi, dengan $a_2=0$, kita dapat menghitung $c_1$ dan $c_2$ yang
           sesuai untuk sebarang $a_0$ dan $a_1$.
           Dengan demikian, rentangnya adalah seluruh himpunan polinomial linear
           $\set{a_0+a_1x\suchthat a_0,a_1\in\Re}$.
           Parametrisasi himpunan tersebut sebagai
           $\set{a_0\cdot 1+a_1\cdot x\suchthat a_0,a_1\in\Re}$
           menyarankan basis $\sequence{1,x}$
           (kita telah menunjukkan bahwa basis tersebut merentang;
           pemeriksaan kebebasan linearnya mudah).
        \partsitem Dari
          \begin{equation*}
            a_0+a_1x+a_2x^2
            =c_1\cdot(2-2x)+c_2\cdot(3+4x^2)
            =(2c_1+3c_2)+(-2c_1)x+(4c_2)x^2
          \end{equation*}
          kita memperoleh sistem berikut.
           \begin{equation*} 
             \begin{linsys}{2}
               2c_1  &+ &3c_2  &= &a_0 \\
               -2c_1 &  &      &= &a_1 \\
                     &  &4c_2  &= &a_2
              \end{linsys}
             \grstep{\rho_1+\rho_2}
             \repeatedgrstep{(-4/3)\rho_2+\rho_3}
             \begin{linsys}{2}
               2c_1  &+ &3c_2  &= &a_0\hfill\hbox{} \\
                     &  &3c_2  &= &a_0+a_1\hfill\hbox{} \\
                     &  &0     &= &(-4/3)a_0-(4/3)a_1+a_2
              \end{linsys}
           \end{equation*}
           Jadi, satu-satunya polinomial kuadratik $a_0+a_1x+a_2x^2$ yang
           memiliki nilai-nilai $c$ terkait adalah polinomial yang memenuhi
           $0=(-4/3)a_0-(4/3)a_1+a_2$.
           Karena itu, rentangnya adalah sebagai berikut.
           \begin{equation*}
             \set{(-a_1+(3/4)a_2)+a_1x+a_2x^2\suchthat a_1,a_2\in\Re}
           \end{equation*}
           Parametrisasi menghasilkan
           $\set{a_1\cdot (-1+x)+a_2\cdot ((3/4)+x^2)\suchthat a_1,a_2\in\Re}$,
           yang menyarankan $\sequence{-1+x,(3/4)+x^2}$
           (pemeriksaan bahwa barisan tersebut bebas linear bersifat rutin).
      \end{exparts} 
    \end{answer}
  \recommended \item 
    Carilah suatu basis bagi masing-masing subruang dari ruang polinomial kubik
    $\polyspace_3$ berikut.
    \begin{exparts}
      \partsitem  Subruang polinomial kubik $p(x)$ sedemikian sehingga
        $p(7)=0$
      \partsitem  Subruang polinomial $p(x)$ sedemikian sehingga
        $p(7)=0$ dan $p(5)=0$
      \partsitem  Subruang polinomial $p(x)$ sedemikian sehingga
        $p(7)=0$, $p(5)=0$, dan~$p(3)=0$
      \partsitem  Ruang polinomial $p(x)$ sedemikian sehingga
        $p(7)=0$, $p(5)=0$, $p(3)=0$, dan~$p(1)=0$
    \end{exparts}
    \begin{answer}
       \begin{exparts}
        \partsitem Subruangnya adalah sebagai berikut.
          \begin{equation*}
             \set{a_0+a_1x+a_2x^2+a_3x^3 \suchthat a_0+7a_1+49a_2+343a_3=0 }
          \end{equation*}
          Dengan menulis ulang $a_0=-7a_1-49a_2-343a_3$, diperoleh
          \begin{equation*}
             \set{(-7a_1-49a_2-343a_3)+a_1x+a_2x^2+a_3x^3 
               \suchthat a_1,a_2,a_3\in\Re }
          \end{equation*}
          Pemisahan parameter-parameter menyarankan
          \( \sequence{-7+x,-49+x^2,-343+x^3} \) sebagai basis
          (yang mudah diperiksa).
        \partsitem Subruang yang diberikan merupakan kumpulan polinomial kubik
          $p(x)=a_0+a_1x+a_2x^2+a_3x^3$ sedemikian sehingga
          $a_0+7a_1+49a_2+343a_3=0$ dan $a_0+5a_1+25a_2+125a_3=0$.
          Metode Gauss
          \begin{equation*}
            \begin{linsys}{4}
              a_0  &+  &7a_1  &+  &49a_2  &+  &343a_3  &=  &0  \\
              a_0  &+  &5a_1  &+  &25a_2  &+  &125a_3  &=  &0    
            \end{linsys}
            \grstep{-\rho_1+\rho_2}
            \begin{linsys}{4}
              a_0  &+  &7a_1  &+  &49a_2  &+  &343a_3  &=  &0  \\
                   &   &-2a_1 &-  &24a_2  &-  &218a_3  &=  &0    
            \end{linsys}
          \end{equation*}
          menghasilkan $a_1=-12a_2-109a_3$ dan $a_0=35a_2+420a_3$.
          Penulisan ulang
          $(35a_2+420a_3)+(-12a_2-109a_3)x+a_2x^2+a_3x^3$ sebagai
          $a_2\cdot(35-12x+x^2)+a_3\cdot(420-109x+x^3)$ menyarankan basis
          $\sequence{35-12x+x^2,420-109x+x^3}$.
          Perhitungan di atas menunjukkan bahwa basis tersebut merentang ruang.
          Pemeriksaan kebebasan linearnya bersifat rutin.
          (\textit{Catatan.}
          Pemeriksaan yang berguna ialah memastikan bahwa kedua polinomial
          dalam basis memiliki tujuh dan lima sebagai akar.)
        \partsitem Di sini terdapat tiga syarat pada polinomial kubik, yaitu
          $a_0+7a_1+49a_2+343a_3=0$, $a_0+5a_1+25a_2+125a_3=0$, dan
          $a_0+3a_1+9a_2+27a_3=0$.
          Metode Gauss
          \begin{equation*}
            \begin{linsys}{4}
              a_0  &+  &7a_1  &+  &49a_2  &+  &343a_3  &=  &0  \\
              a_0  &+  &5a_1  &+  &25a_2  &+  &125a_3  &=  &0  \\    
              a_0  &+  &3a_1  &+  &9a_2   &+  &27a_3   &=  &0    
            \end{linsys}
            \grstep[-\rho_1+\rho_3]{-\rho_1+\rho_2}
            \repeatedgrstep{-2\rho_2+\rho_3}
            \begin{linsys}{4}
              a_0  &+  &7a_1  &+  &49a_2  &+  &343a_3  &=  &0  \\
                   &   &-2a_1 &-  &24a_2  &-  &218a_3  &=  &0  \\    
                   &   &      &   &8a_2   &+  &120a_3  &=  &0    
            \end{linsys}
          \end{equation*}
          menghasilkan satu variabel bebas $a_3$, dengan
          $a_2=-15a_3$, $a_1=71a_3$, dan $a_0=-105a_3$.
          Parametrisasi yang diperoleh adalah sebagai berikut.
          \begin{multline*} 
            \set{(-105a_3)+(71a_3)x+(-15a_3)x^2+(a_3)x^3\suchthat a_3\in\Re} \\
            =
            \set{a_3\cdot(-105+71x-15x^2+x^3)\suchthat a_3\in\Re}
          \end{multline*}
          Oleh karena itu, calon basis yang wajar adalah
          $\sequence{-105+71x-15x^2+x^3}$.
          Berdasarkan perhitungan di atas, basis tersebut merentang ruang.
          Basis itu jelas bebas linear karena merupakan himpunan berelemen satu
          (dan elemen tunggal tersebut bukan objek nol ruang).
          Jadi, setiap polinomial kubik yang melalui ketiga titik $(7,0)$,
          $(5,0)$, dan $(3,0)$ merupakan kelipatan polinomial ini.
          (\textit{Catatan.}
          Seperti dalam pertanyaan sebelumnya, pemeriksaan yang berguna ialah
          memastikan bahwa substitusi tujuh, lima, dan tiga ke dalam polinomial
          ini selalu menghasilkan nol.)
        \partsitem Ini merupakan subruang trivial dari $\polyspace_3$.
          Jadi, basisnya kosong, yaitu $\sequence{}$.
      \end{exparts}
      \noindent\textit{Catatan.}
      Sebagai alternatif, polinomial dalam butir ketiga dapat kita peroleh
      dengan mengalikan $(x-7)(x-5)(x-3)$.
     \end{answer}
  \item 
     Kita telah melihat bahwa hasil pengurutan ulang suatu basis dapat menjadi
     basis lain.
     Apakah hasilnya selalu merupakan basis?
     \begin{answer}
       Ya.
       Kebebasan linear dan rentang tidak berubah karena pengurutan ulang.
     \end{answer}
  \item  
    Dapatkah suatu basis memuat vektor nol?
    \begin{answer}
      Tidak ada himpunan bebas linear yang memuat vektor nol.
    \end{answer}
  \recommended \item
    Misalkan \( \sequence{\vec{\beta}_1,\vec{\beta}_2,\vec{\beta}_3} \)
    merupakan basis bagi suatu ruang vektor.
    \begin{exparts}
      \partsitem Tunjukkan bahwa
        \( \sequence{c_1\vec{\beta}_1,c_2\vec{\beta}_2,c_3\vec{\beta}_3} \)
        merupakan basis ketika \( c_1, c_2, c_3\neq 0 \).
        Apa yang terjadi ketika sedikitnya satu \( c_i \) bernilai $0$?
      \partsitem Buktikan bahwa
        \( \sequence{\vec{\alpha}_1,\vec{\alpha}_2,\vec{\alpha}_3} \)
        merupakan basis, dengan \( \vec{\alpha}_i=\vec{\beta}_1+\vec{\beta}_i \).
    \end{exparts}
    \begin{answer}
      \begin{exparts}
         \partsitem Untuk menunjukkan bahwa barisan tersebut bebas linear,
           perhatikan bahwa jika
           \( d_1(c_1\vec{\beta}_1)+d_2(c_2\vec{\beta}_2)+d_3(c_3\vec{\beta}_3)
               =\zero \)
           maka
           \( (d_1c_1)\vec{\beta}_1+(d_2c_2)\vec{\beta}_2+(d_3c_3)\vec{\beta}_3
               =\zero \), yang selanjutnya menyiratkan bahwa setiap
           \( d_ic_i \) bernilai nol.
           Namun, dengan \( c_i\neq 0 \), hal itu berarti setiap \( d_i \)
           bernilai nol.
           Pembuktian bahwa barisan tersebut merentang ruang hampir sama;
           karena $\sequence{\vec{\beta}_1,\vec{\beta}_2,\vec{\beta}_3}$
           merupakan basis sehingga merentang ruang, untuk sebarang $\vec{v}$
           kita dapat menulis
           \( \vec{v}=d_1\vec{\beta}_1+d_2\vec{\beta}_2+d_3\vec{\beta}_3 \),
           lalu
           \( \vec{v}=(d_1/c_1)(c_1\vec{\beta}_1)
               +(d_2/c_2)(c_2\vec{\beta}_2)+(d_3/c_3)(c_3\vec{\beta}_3) \).

           Jika salah satu skalarnya nol, hasilnya bukan basis karena tidak
           bebas linear.
         \partsitem Menunjukkan bahwa
           $\sequence{2\vec{\beta}_1,\vec{\beta}_1+\vec{\beta}_2,
             \vec{\beta}_1+\vec{\beta}_3}$ bebas linear mudah dilakukan.
           Untuk menunjukkan bahwa barisan tersebut merentang ruang, misalkan
           \( \vec{v}=d_1\vec{\beta}_1+d_2\vec{\beta}_2+d_3\vec{\beta}_3 \).
           Maka, kita dapat merepresentasikan \( \vec{v} \) yang sama terhadap
           \( \sequence{2\vec{\beta}_1,\vec{\beta}_1+\vec{\beta}_2,
                         \vec{\beta}_1+\vec{\beta}_3} \)
           dengan cara berikut
           $\vec{v}=(1/2)(d_1-d_2-d_3)(2\vec{\beta}_1)
           +d_2(\vec{\beta}_1+\vec{\beta}_2)+d_3(\vec{\beta}_1+\vec{\beta}_3)$.
      \end{exparts}   
    \end{answer}
  \item 
    Carilah satu vektor $\vec{v}$ yang membuat masing-masing barisan berikut
    menjadi basis bagi ruangnya.
    \begin{exparts*}
      \partsitem $\sequence{\colvec[r]{1 \\ 1},\vec{v}}$ di $\Re^2$
      \partsitem $\sequence{\colvec[r]{1 \\ 1 \\ 0},
                            \colvec[r]{0 \\ 1 \\ 0},\vec{v}}$ di $\Re^3$
      \partsitem $\sequence{x,1+x^2,\vec{v}}$ di $\polyspace_2$
    \end{exparts*} 
    \begin{answer}
      Jika $\vec{v}$ dihilangkan, masing-masing membentuk himpunan bebas linear.
      Untuk mempertahankan kebebasan linear, kita harus memperluas rentang
      masing-masing himpunan.
      Artinya, kita harus menentukan rentang masing-masing himpunan (tanpa
      $\vec{v}$), lalu memilih $\vec{v}$ yang berada di luar rentang tersebut.
      Terakhir, kita harus memeriksa bahwa hasilnya merentang seluruh ruang yang
      diberikan.
      Pemeriksaan-pemeriksaan itu bersifat rutin.
      \begin{exparts}
        \partsitem Sebarang vektor yang bukan kelipatan vektor yang diberikan,
          yakni sebarang vektor yang tidak terletak pada garis $y=x$, dapat
          digunakan.
          Salah satunya adalah $\vec{v}=\vec{e}_1$.
        \partsitem Dari pemeriksaan langsung, kita melihat bahwa vektor
          $\vec{e}_3$ tidak berada dalam rentang himpunan kedua vektor yang
          diberikan.
          Pemeriksaan bahwa himpunan yang dihasilkan merupakan basis bagi
          $\Re^3$ bersifat rutin.
        \partsitem Bagi setiap anggota rentang
          $\set{c_1\cdot(x)+c_2\cdot(1+x^2)\suchthat c_1,c_2\in\Re}$,
          koefisien $x^2$ sama dengan suku konstantanya.
          Jadi, kita memperluas rentangnya jika menambahkan polinomial kuadratik
          yang tidak memiliki sifat ini, misalnya $\vec{v}=1-x^2$.
          Pemeriksaan bahwa hasilnya merupakan basis bagi $\polyspace_2$ mudah.
      \end{exparts}
    \end{answer}
  \recommended \item Perhatikan
    \( 2+4x^2, 1+3x^2, 1+5x^2 \in\polyspace_2\).
    \begin{exparts}
      \partsitem Carilah suatu relasi linear di antara ketiganya.
      \partsitem Representasikan ketiganya terhadap
        $B=\sequence{1-x,1+x,x^2}$.
      \partsitem Periksalah bahwa relasi linear yang sama berlaku di antara
        representasi-representasinya, seperti dalam
        \nearbylemma{lem:LinearRelRepresented}.
    \end{exparts}
    \begin{answer}
    \begin{exparts}
      \partsitem $1\cdot(2+4x^2)-3\cdot(1+3x^2)+1\cdot(1+5x^2)=0+0x+0x^2$
      \partsitem Perhitungan sederhana menghasilkan
        \begin{equation*}
          \rep{2+4x^2}{B}=\colvec{1 \\ 1 \\ 4}
          \quad
          \rep{1+3x^2}{B}=\colvec{1/2 \\ 1/2 \\ 3}
          \quad
          \rep{1+5x^2}{B}=\colvec{1/2 \\ 1/2 \\ 5}
        \end{equation*}
      \partsitem Pemeriksaannya langsung.
        \begin{equation*}
          1\cdot\colvec{1 \\ 1 \\ 4}
          -3\cdot\colvec{1/2 \\ 1/2 \\ 3}
          +1\cdot\colvec{1/2 \\ 1/2 \\ 5}
          =\colvec{0 \\ 0 \\ 0}
        \end{equation*}
    \end{exparts}
    \end{answer}
  \recommended \item  
    Jika
    \( \sequence{\vec{\beta}_1,\dots,\vec{\beta}_n } \)
    merupakan basis, tunjukkan bahwa dalam persamaan berikut
    \begin{equation*}
       c_1\vec{\beta}_1+\dots+c_k\vec{\beta}_k
       =
       c_{k+1}\vec{\beta}_{k+1}+\dots+c_n\vec{\beta}_n
    \end{equation*}
    setiap \( c_i \) bernilai nol.
    Generalisasikan hasilnya.
    \begin{answer}
      Untuk menunjukkan bahwa setiap skalar bernilai nol, cukup kurangkan
      \( c_1\vec{\beta}_1+\dots+c_k\vec{\beta}_k
          -c_{k+1}\vec{\beta}_{k+1}-\dots-c_n\vec{\beta}_n=\zero \).
      Generalisasi yang jelas adalah bahwa dalam sebarang persamaan yang hanya
      melibatkan vektor-vektor \( \vec{\beta} \), dan setiap
      \( \vec{\beta} \) hanya muncul satu kali, setiap skalar bernilai nol.
      Sebagai contoh, persamaan dengan kombinasi vektor-vektor basis berindeks
      genap (yakni $\vec{\beta}_2$, $\vec{\beta}_4$, dan seterusnya) di ruas
      kanan dan vektor-vektor basis berindeks ganjil di ruas kiri juga
      menghasilkan kesimpulan bahwa semua koefisien bernilai nol.
    \end{answer}
  \item  
    Suatu basis memuat sebagian vektor dari sebuah ruang vektor; dapatkah basis
    tersebut memuat semuanya?
    \begin{answer}
      Tidak; tidak ada himpunan bebas linear yang memuat vektor nol.
    \end{answer}
  \item  
    \nearbytheorem{th:BasisIffUniqueRepWRT} menunjukkan bahwa, terhadap suatu
    basis, setiap kombinasi linear bersifat unik.
    Jika suatu subhimpunan bukan basis, dapatkah kombinasi linearnya tidak unik?
    Jika dapat, apakah kombinasi linearnya pasti tidak unik?
    \begin{answer}
      Berikut adalah subhimpunan dari $\Re^2$ yang bukan basis, beserta dua
      kombinasi linear berbeda dari elemen-elemennya yang menghasilkan vektor
      yang sama.
      \begin{equation*}
        \set{\colvec[r]{1 \\ 2},\colvec[r]{2 \\ 4}}
        \qquad
        2\cdot\colvec[r]{1 \\ 2}+0\cdot\colvec[r]{2 \\ 4}
        =0\cdot\colvec[r]{1 \\ 2}+1\cdot\colvec[r]{2 \\ 4}
      \end{equation*}
      Jadi, ketika suatu subhimpunan bukan basis, kombinasi linearnya dapat
      bersifat tidak unik.

      Namun, fakta bahwa suatu subhimpunan bukan basis tidak menyiratkan bahwa
      kombinasi linearnya pasti tidak unik.
      Sebagai contoh, himpunan berikut
      \begin{equation*}
        \set{\colvec[r]{1 \\ 2}}
      \end{equation*}
      memiliki sifat bahwa
      \begin{equation*}
        c_1\cdot\colvec[r]{1 \\ 2}
        =
        c_2\cdot\colvec[r]{1 \\ 2}
      \end{equation*}
      menyiratkan $c_1=c_2$.
      Intinya, subhimpunan ini gagal menjadi basis karena tidak merentang ruang;
      pembuktian teorema tersebut menetapkan bahwa kombinasi linear bersifat
      unik jika dan hanya jika subhimpunannya bebas linear.
     \end{answer}
  \item
    Suatu matriks persegi disebut \definend{simetris}\index{matriks!simetris}%
    \index{matriks simetris}
    jika untuk semua indeks \( i \) dan \( j \), entri \( i,j \) sama dengan
    entri \( j,i \).
    \begin{exparts}
      \partsitem Carilah suatu basis bagi ruang vektor matriks simetris
        \( \nbyn{2} \).
      \partsitem Carilah suatu basis bagi ruang matriks simetris \( \nbyn{3} \).
      \partsitem Carilah suatu basis bagi ruang matriks simetris \( \nbyn{n} \).
    \end{exparts}
    \begin{answer}
      \begin{exparts}
        \partsitem Mendeskripsikan ruang vektor tersebut sebagai
          \begin{equation*}
             \set{\begin{mat}
                     a  &b  \\
                     b  &c
                  \end{mat}  \suchthat a,b,c\in\Re}
          \end{equation*}
          menyarankan basis berikut.
          \begin{equation*}
            \sequence{
              \begin{mat}[r]
                1  &0  \\
                0  &0
              \end{mat},
              \begin{mat}[r]
                0  &0  \\
                0  &1
              \end{mat},
              \begin{mat}[r]
                0  &1  \\
                1  &0
              \end{mat}  }
          \end{equation*}
          Pemeriksaannya mudah.
        \partsitem Berikut adalah salah satu basis yang mungkin.
          \begin{equation*}
            \sequence{
              \begin{mat}[r]
                1  &0  &0  \\
                0  &0  &0  \\
                0  &0  &0
              \end{mat},
              \begin{mat}[r]
                0  &0  &0  \\
                0  &1  &0  \\
                0  &0  &0
              \end{mat},
              \begin{mat}[r]
                0  &0  &0  \\
                0  &0  &0  \\
                0  &0  &1
              \end{mat},
              \begin{mat}[r]
                0  &1  &0  \\
                1  &0  &0  \\
                0  &0  &0
              \end{mat},
              \begin{mat}[r]
                0  &0  &1  \\
                0  &0  &0  \\
                1  &0  &0
              \end{mat},
              \begin{mat}[r]
                0  &0  &0  \\
                0  &0  &1  \\
                0  &1  &0
              \end{mat}  }
          \end{equation*}
         \partsitem Seperti dalam dua pertanyaan sebelumnya, kita dapat
           membentuk basis dari dua jenis matriks.
           Jenis pertama adalah matriks dengan satu entri bernilai satu pada
           diagonal dan semua entri lainnya nol (terdapat \( n \) matriks
           semacam itu).
           Jenis kedua adalah matriks dengan dua entri luar diagonal yang
           saling berhadapan bernilai satu dan semua entri lainnya nol.
           (Artinya, semua entri dalam $M$ bernilai nol, kecuali $m_{i,j}$ dan
           $m_{j,i}$ yang bernilai satu.)
      \end{exparts}  
    \end{answer}
  \item  
     Kita dapat menunjukkan bahwa setiap basis bagi $\Re^3$ memuat jumlah
     vektor yang sama.
     \begin{exparts}
       \partsitem Tunjukkan bahwa tidak ada subhimpunan bebas linear dari
          $\Re^3$ yang memuat lebih dari tiga vektor.
       \partsitem Tunjukkan bahwa tidak ada subhimpunan perentang dari $\Re^3$
          yang memuat kurang dari tiga vektor.
          \textit{Petunjuk:}
          ingatlah cara menghitung rentang suatu himpunan dan tunjukkan bahwa
          metode tersebut tidak dapat menghasilkan seluruh $\Re^3$ ketika
          diterapkan pada kurang dari tiga vektor.
     \end{exparts}
     \begin{answer}
       \begin{exparts}
         \partsitem Sebarang empat vektor dari $\Re^3$ berelasi linear karena
           persamaan vektor
           \begin{equation*}
             c_1\colvec{x_1 \\ y_1 \\ z_1}
             +c_2\colvec{x_2 \\ y_2 \\ z_2}
             +c_3\colvec{x_3 \\ y_3 \\ z_3}
             +c_4\colvec{x_4 \\ y_4 \\ z_4}
             =\colvec[r]{0 \\ 0 \\ 0}
           \end{equation*}
           menghasilkan sistem linear
           \begin{equation*}
             \begin{linsys}{4}
                x_1c_1  &+  &x_2c_2  &+  &x_3c_3  &+  &x_4c_4  &=  &0 \\
                y_1c_1  &+  &y_2c_2  &+  &y_3c_3  &+  &y_4c_4  &=  &0 \\
                z_1c_1  &+  &z_2c_2  &+  &z_3c_3  &+  &z_4c_4  &=  &0
             \end{linsys}
           \end{equation*}
           yang homogen (sehingga memiliki solusi) dan memiliki empat variabel
           tak diketahui tetapi hanya tiga persamaan, sehingga memiliki solusi
           tak trivial.
           (Tentu saja, argumen ini berlaku bagi sebarang subhimpunan dari
           $\Re^3$ dengan empat vektor atau lebih.)
         \partsitem Kita cukup membahas kasus dua vektor.
           Untuk nilai-nilai $x_1$, \ldots, $z_2$ yang diberikan, misalkan
           \begin{equation*}
             S=\set{\colvec{x_1 \\ y_1 \\ z_1},
             \colvec{x_2 \\ y_2 \\ z_2}} 
           \end{equation*}
           Untuk menentukan vektor-vektor mana
           \begin{equation*}
             \colvec{x \\ y \\ z}
           \end{equation*}
           yang berada dalam rentang $S$, susun
           \begin{equation*}
             c_1\colvec{x_1 \\ y_1 \\ z_1}
             +c_2\colvec{x_2 \\ y_2 \\ z_2}
             =\colvec{x \\ y \\ z}
           \end{equation*}
           lalu reduksi baris sistem yang dihasilkan.
           \begin{equation*}
             \begin{linsys}{2}
                x_1c_1  &+  &x_2c_2   &=  &x \\
                y_1c_1  &+  &y_2c_2   &=  &y \\
                z_1c_1  &+  &z_2c_2   &=  &z
             \end{linsys}
           \end{equation*}
           Terdapat dua variabel $c_1$ dan $c_2$, tetapi tiga persamaan, sehingga
           setelah Metode Gauss selesai, baris terbawah akan memberikan suatu
           relasi berbentuk $0=m_1x+m_2y+m_3z$.
           Dengan demikian, vektor-vektor dalam rentang himpunan berelemen dua
           $S$ harus memenuhi suatu pembatasan.
           Karena itu, rentangnya bukan seluruh $\Re^3$.
       \end{exparts}  
    \end{answer}
  \item 
    Salah satu latihan dalam subbagian Subruang menunjukkan bahwa himpunan
    \begin{equation*}
      \set{\colvec{x \\ y \\ z}\suchthat x+y+z=1}
    \end{equation*}
    merupakan ruang vektor dengan operasi-operasi berikut.
    \begin{equation*}
      \colvec{x_1 \\ y_1 \\ z_1}+\colvec{x_2 \\ y_2 \\ z_2}
      =\colvec{x_1+x_2-1 \\ y_1+y_2 \\ z_1+z_2}
      \qquad
      r\colvec{x \\ y \\ z}=\colvec{rx-r+1 \\ ry \\ rz}
    \end{equation*}
    Carilah suatu basis.
    \begin{answer}
      Dengan menggunakan operasi-operasi tak lazim ini secara cermat, kita
      memperoleh
      \begin{multline*}
        \set{\colvec{1-y-z \\ y \\ z}\suchthat y,z\in\Re}
         =\set{\colvec{-y+1 \\ y \\ 0}
               +\colvec{-z+1 \\ 0 \\ z}
              \suchthat y,z\in\Re}                           \\
         =\set{y\cdot\colvec{0 \\ 1 \\ 0}
               +z\cdot\colvec{0 \\ 0 \\ 1}
              \suchthat y,z\in\Re}
      \end{multline*}
      sehingga calon basis yang wajar adalah sebagai berikut.
      \begin{equation*}
        \sequence{\colvec[r]{0 \\ 1 \\ 0},\colvec[r]{0 \\ 0 \\ 1}}
      \end{equation*}
      Untuk memeriksa kebebasan linear, kita susun
      \begin{equation*}
         c_1\colvec[r]{0 \\ 1 \\ 0}+c_2\colvec[r]{0 \\ 0 \\ 1}
         =\colvec[r]{1 \\ 0 \\ 0}
      \end{equation*}
      (vektor di ruas kanan adalah objek nol dalam ruang ini).
      Hal itu menghasilkan sistem linear
      \begin{equation*}
        \begin{linsys}{3}
          (-c_1+1)  &+  &(-c_2+1)  &-  &1  &=  &1  \\
             c_1    &   &          &   &   &=  &0  \\
                    &   &c_2       &   &   &=  &0
        \end{linsys}
      \end{equation*}
      yang hanya memiliki solusi $c_1=0$ dan $c_2=0$.
      Pemeriksaan rentangnya serupa.
   \end{answer}
\end{exercises}














\subsection{Dimensi}
Subbagian sebelumnya mendefinisikan basis suatu ruang vektor dan menunjukkan
bahwa suatu ruang dapat memiliki banyak basis yang berbeda.
% For example, following the definition of a basis, we saw three
% different bases for $\Re^2$.
Karena itu, kita tidak dapat berbicara tentang ``basis tunggal'' bagi suatu
ruang vektor.
Memang, beberapa ruang vektor memiliki basis yang terasa lebih alami daripada
basis lainnya, misalnya basis $\stdbasis_2$ bagi $\Re^2$
% or $\Re^3$'s basis $\stdbasis_3$
atau basis $\sequence{1,x,x^2}$ bagi $\polyspace_2$.
Namun, bagi ruang vektor
$\set{a_2x^2+a_1x+a_0\suchthat 2a_2-a_0=a_1}$, tidak ada basis tertentu yang
langsung tampak sebagai basis alaminya.
Secara umum, kita tidak dapat mengaitkan satu basis tunggal dengan suatu ruang
sebagai basis yang paling baik mendeskripsikannya.

Namun, kita dapat menemukan suatu sifat basis yang terkait secara unik dengan
ruangnya.
Subbagian ini menunjukkan bahwa sebarang dua basis bagi suatu ruang memiliki
jumlah elemen yang sama.
Jadi, dengan setiap ruang kita dapat mengaitkan suatu bilangan, yaitu jumlah
vektor dalam sebarang basisnya.

Sebelum memulai, kita batasi perhatian pada ruang yang sedikitnya memiliki satu
basis dengan hanya berhingga banyak anggota.

\begin{definition} \label{df:FiniteDimensional}
%<*df:FiniteDimensional>
Suatu ruang vektor disebut \definend{berdimensi hingga\/}%
\index{ruang vektor!berdimensi hingga}\index{ruang vektor berdimensi hingga}
jika memiliki basis dengan hanya berhingga banyak vektor.
%</df:FiniteDimensional>
\end{definition}

\noindent Salah satu ruang yang tidak berdimensi hingga adalah himpunan
polinomial dengan koefisien riil,
\nearbyexample{ex:PolysOfAllFiniteDegrees}. 
Ruang ini tidak direntang oleh sebarang subhimpunan hingga karena subhimpunan
semacam itu akan memiliki polinomial berderajat terbesar, sedangkan ruang ini
memiliki polinomial dari setiap derajat.
Ruang semacam itu menarik dan penting, tetapi kita akan berfokus ke arah lain.
Mulai sekarang, kita hanya mempelajari ruang vektor berdimensi hingga.
Dalam sisa buku ini, `ruang vektor' akan berarti `ruang vektor berdimensi
hingga'.

% \begin{remark}
% One reason for sticking to finite-dimensional spaces is so that 
% the representation of a vector with respect to a 
% basis is a finitely-tall vector and we can easily write it.
% Another reason is that 
% the statement `any infinite-dimensional vector space has a basis'
% is equivalent to a statement called the Axiom of Choice
% \cite{Blass84} and so covering this would 
% move us far past this book's scope.
% (A discussion of the Axiom of Choice is in the
% Frequently Asked Questions list for \texttt{sci.math}, and
% another accessible one is \cite{Rucker}.)  
% \end{remark}

% \begin{theorem} \label{th:AllBasesSameSize}
% %<*th:AllBasesSameSizeCoord>
% In any finite-dimensional vector space, all 
% bases have the same number of elements.\index{basis}
% %</th:AllBasesSameSizeCoord>
% \end{theorem}

% \begin{proof}
% Fix a vector space with at least one finite basis.
% From among all of this space's bases, choose
% one \( B=\sequence{\vec{\beta}_1,\dots,\vec{\beta}_n} \) of minimal size.
% We will show that any other basis
% \( D=\smash{\sequence{\vec{\delta}_1,\vec{\delta}_2,\ldots}} \)
% has the same number of members, $n$.
% Because \( B \) has minimal size, \( D \) has no fewer than \( n \) vectors.
% We will argue that it cannot have more than \( n \) vectors.

% So suppose that
% $\vec{\delta}_1$, \ldots{}, $\vec{\delta}_{n+1}$ are distinct.
% The assumption that $D$ is linearly independent gives that
% the only relationship
% $a_1\vec{\delta}+1+\dots+a_{n+1}\vec{\delta}_{n+1}=\vec{0}$
% is the one where each $a_i$ is zero.
% We shall get a contradiction by finding that 
% there are nontrivial relationships among these~$n+1$ vectors.

% Represent them with respect to~$B$.
% \begin{equation*}
%   \rep{\vec{\delta}_1}{B}=\colvec{c_{1,1} \\ \vdots \\ c_{n,1}}
%   \quad\cdots\quad
%   \rep{\vec{\delta}_{n+1}}{B}=\colvec{c_{1,n+1} \\ \vdots \\ c_{n,n+1}}  
% \end{equation*}
% \nearbylemma{lem:LinearRelRepresented} says that a relationship
% holds among the vectors
% $a_1\vec{\delta}_1+\dots+a_{n+1}\vec{\delta}_{n+1}=\vec{0}$ 
% if and only if the same relationship holds among the representations.
% \begin{equation*}
%   a_1\colvec{c_{1,1} \\ \vdots \\ c_{n,1}}
%   +\cdots+
%   a_{n+1}\colvec{c_{1,n+1} \\ \vdots \\ c_{n,n+1}}
%   =\colvec{0 \\ \vdots \\ 0}  
% \end{equation*}
% Here the $c_{i,j}$ are fixed, since they are the coefficients in the 
% representations of the given vectors, while we are looking for the~$a_i$.
% Thus the above is a homogeneous linear system
% \begin{equation*}
%   \begin{linsys}{3}
%     c_{1,1}a_1 &+ &\cdots &+ &c_{1,n+1}a_{n+1} &= &0 \\ 
%               &  &        &  &               &\vdots \\ 
%     c_{n,1}a_1 &+ &\cdots &+ &c_{n,n+1}a_{n+1} &= &0 
%   \end{linsys}
% \end{equation*}
% with more unknowns, $n+1$, than equations.
% Such a system does not have only one solution, it has infinitely many solutions.
% \end{proof}


Untuk membuktikan teorema utama, kita akan menggunakan suatu hasil teknis,
yaitu Lema Pertukaran.
Pertama, kita mengilustrasikannya dengan sebuah contoh.

\begin{example}
Berikut adalah suatu basis bagi $\Re^3$ dan suatu vektor yang diberikan sebagai
kombinasi linear anggota-anggota basis tersebut.
\begin{equation*}
  B=\sequence{\colvec[r]{1 \\ 0  \\ 0},
              \colvec[r]{1 \\ 1 \\ 0},
              \colvec[r]{0 \\ 0 \\ 2}}
  \qquad
  \colvec[r]{1 \\ 2 \\ 0}
  =(-1)\cdot\colvec[r]{1 \\ 0  \\ 0}
   +2\colvec[r]{1 \\ 1 \\ 0}
   +0\cdot\colvec[r]{0 \\ 0 \\ 2}
\end{equation*}
Dua vektor basis memiliki koefisien tak nol.
Pilih salah satunya, misalnya vektor pertama.
Gantilah vektor tersebut dengan vektor yang telah kita nyatakan sebagai
kombinasi,
\begin{equation*}
  \hat{B}=\sequence{\colvec[r]{1 \\ 2  \\ 0},
              \colvec[r]{1 \\ 1 \\ 0},
              \colvec[r]{0 \\ 0 \\ 2}}
\end{equation*}
dan hasilnya merupakan basis lain bagi \( \Re^3 \).
\end{example}

\begin{lemma}[Lema Pertukaran] \label{lm:ExchangeLemma}
%<*lm:ExchangeLemma>
Misalkan \( B=\sequence{\vec{\beta}_1,\dots,\vec{\beta}_n} \) merupakan basis
bagi suatu ruang vektor, dan bagi vektor \( \vec{v} \), relasi
\( \vec{v}=\lincombo{c}{\vec{\beta}} \) memiliki \( c_i\neq 0 \).
Maka, mengganti \( \vec{\beta}_i \) dengan \( \vec{v} \) menghasilkan basis
lain bagi ruang tersebut.
%</lm:ExchangeLemma>
\end{lemma}

\begin{proof}
%<*pf:ExchangeLemma0>
Nyatakan hasil pertukaran tersebut dengan
\( \hat{B}=\sequence{\vec{\beta}_1,\dots,\vec{\beta}_{i-1},\vec{v},
                       \vec{\beta}_{i+1},\dots,\vec{\beta}_n}  \).

Pertama, kita tunjukkan bahwa $\hat{B}$ bebas linear.
Sebarang relasi
\( d_1\vec{\beta}_1+\dots+d_i\vec{v}+\dots+d_n\vec{\beta}_n=\zero \)
di antara anggota-anggota $\hat{B}$, setelah menyubstitusikan $\vec{v}$,
\begin{equation*}
  d_1\vec{\beta}_1+\dots
    +d_i\cdot(c_1\vec{\beta}_1+\dots+c_i\vec{\beta}_i+\dots+c_n\vec{\beta}_n)
    +\dots+d_n\vec{\beta}_n
  =\zero
\tag*{($*$)}\end{equation*}
menghasilkan relasi linear di antara anggota-anggota $B$.
Basis $B$ bebas linear, sehingga koefisien $d_ic_i$ dari
$\vec{\beta}_i$ bernilai nol.
Karena kita mengasumsikan $c_i$ tak nol, maka $d_i=0$.
Menggunakan hasil ini dalam persamaan~$(*)$ menunjukkan bahwa semua nilai $d$
lainnya juga nol.
Oleh karena itu, $\hat{B}$ bebas linear.
%</pf:ExchangeLemma0>

%<*pf:ExchangeLemma1>
Kita selesaikan pembuktian dengan menunjukkan bahwa $\hat{B}$ memiliki rentang
yang sama dengan $B$.
Separuh argumen, yaitu
$\spanof{\smash{\hat{B}}}\subseteq\spanof{B}$, mudah; kita dapat menulis
sebarang anggota
$d_1\vec{\beta}_1+\dots+d_i\vec{v}+\dots+d_n\vec{\beta}_n$
dari $\spanof{\smash{\hat{B}}}$ sebagai
$
  d_1\vec{\beta}_1+\dots
    +d_i\cdot(c_1\vec{\beta}_1+\dots+c_n\vec{\beta}_n)
    +\dots+d_n\vec{\beta}_n
$,
yang merupakan kombinasi linear dari kombinasi-kombinasi linear anggota $B$,
sehingga berada dalam $\spanof{B}$.
Untuk separuh argumen
$\spanof{B}\subseteq\spanof{\smash{\hat{B}}}$, ingat bahwa jika
$\vec{v}=c_1\vec{\beta}_1+\dots+c_n\vec{\beta}_n$ dengan $c_i\neq 0$
maka kita dapat menyusun ulang persamaan menjadi
$\vec{\beta}_i=(-c_1/c_i)\vec{\beta}_1+\dots+(1/c_i)\vec{v}+\dots
   +(-c_n/c_i)\vec{\beta}_n$.
Sekarang, perhatikan sebarang anggota
$d_1\vec{\beta}_1+\dots+d_i\vec{\beta}_i+\dots+d_n\vec{\beta}_n$
dari $\spanof{B}$, substitusikan pernyataan $\vec{\beta}_i$ sebagai kombinasi
linear anggota-anggota $\hat{B}$, lalu perhatikan, seperti dalam separuh pertama
argumen, bahwa hasilnya merupakan kombinasi linear dari kombinasi-kombinasi
linear anggota $\hat{B}$, sehingga berada dalam
$\spanof{\smash{\hat{B}}}$.
%</pf:ExchangeLemma1>
\end{proof}

\begin{theorem} \label{th:AllBasesSameSize}
%<*th:AllBasesSameSize>
Dalam sebarang ruang vektor berdimensi hingga, semua basis memiliki jumlah
elemen yang sama.
%</th:AllBasesSameSize>
\end{theorem}\index{basis}

\begin{proof}
%<*pf:AllBasesSameSize0>
Tetapkan suatu ruang vektor dengan sedikitnya satu basis hingga.
Dari semua basis ruang tersebut, pilih satu basis
\( B=\sequence{\vec{\beta}_1,\dots,\vec{\beta}_n} \) yang berukuran minimal.
Kita akan menunjukkan bahwa sebarang basis lain
\( D=\smash{\sequence{\vec{\delta}_1,\vec{\delta}_2,\ldots}} \) juga memiliki
jumlah anggota yang sama, yaitu $n$.
Karena \( B \) berukuran minimal, \( D \) tidak memiliki kurang dari
\( n \) vektor.
Kita akan berargumen bahwa \( D \) tidak mungkin memiliki lebih dari
\( n \) vektor.
%</pf:AllBasesSameSize0>

%<*pf:AllBasesSameSize1>
Basis \( B \) merentang ruang dan \( \vec{\delta}_1 \) berada dalam ruang,
sehingga \( \vec{\delta}_1 \) merupakan kombinasi linear tak trivial
elemen-elemen \( B \).
Menurut Lema Pertukaran, kita dapat menukar \( \vec{\delta}_1 \) dengan suatu
vektor dari \( B \), sehingga menghasilkan basis \( B_1 \), dengan satu elemen
berupa \( \vec{\delta}_1 \) dan semua \( n-1 \) elemen lainnya berupa
vektor-vektor \( \vec{\beta} \).
%</pf:AllBasesSameSize1>

%<*pf:AllBasesSameSize2>
Paragraf sebelumnya membentuk kasus dasar bagi argumen induksi.
Langkah induksi dimulai dengan basis \( B_k \) (untuk \( 1\leq k<n \)) yang
memuat \( k \) anggota \( D \) dan \( n-k \) anggota \( B \).
Kita mengetahui bahwa \( D \) memiliki sedikitnya \( n \) anggota, sehingga
terdapat \( \vec{\delta}_{k+1} \).
Representasikan vektor tersebut sebagai kombinasi linear elemen-elemen
\( B_k \).
Pokok pentingnya:~dalam representasi itu, sedikitnya satu skalar tak nol harus
terkait dengan suatu \( \vec{\beta}_i \); jika tidak, representasi tersebut
akan menjadi relasi linear tak trivial di antara elemen-elemen himpunan bebas
linear \( D \).
Tukarkan \( \vec{\delta}_{k+1} \) dengan \( \vec{\beta}_i \) untuk memperoleh
basis baru \( B_{k+1} \) dengan satu \( \vec{\delta} \) lebih banyak dan satu
\( \vec{\beta} \) lebih sedikit daripada basis sebelumnya \( B_k \).
%</pf:AllBasesSameSize2>

%<*pf:AllBasesSameSize3>
Ulangi proses tersebut hingga tidak ada \( \vec{\beta} \) yang tersisa,
sehingga \( B_n \) memuat
$\vec{\delta}_1,\dots,\vec{\delta}_n$.
Sekarang, \( D \) tidak mungkin memiliki lebih dari \( n \) vektor tersebut
karena sebarang \( \vec{\delta}_{n+1} \) yang tersisa akan berada dalam rentang
\( B_n \) (karena \( B_n \) merupakan basis), sehingga menjadi kombinasi
linear dari vektor-vektor $\vec{\delta}$ lainnya; hal ini bertentangan dengan
kebebasan linear $D$.
%</pf:AllBasesSameSize3>
\end{proof}

\begin{definition}  \label{df:Dimension}
%<*df:Dimension>
\definend{Dimensi}\index{dimensi}\index{ruang vektor!dimensi}
suatu ruang vektor adalah jumlah vektor dalam sebarang basisnya.
%</df:Dimension>
\end{definition}

\begin{example}
Sebarang basis bagi \( \Re^n \) memiliki \( n \) vektor karena basis standar
\( \stdbasis_n \) memiliki \( n \) vektor.
Jadi, definisi `dimensi' ini menggeneralisasi penggunaan istilah yang paling
dikenal, yaitu bahwa $\Re^n$ berdimensi $n$.
\end{example}

\begin{example}
Ruang \( \polyspace_n \) dari polinomial berderajat paling tinggi $n$
memiliki dimensi \( n+1 \).
Kita dapat menunjukkannya dengan memberikan sebarang basis\Dash misalnya
$\sequence{1,x,\dots,x^n}$\Dash lalu menghitung anggotanya.
\end{example}

\begin{example}
Ruang fungsi
$\set{a\cdot\cos\theta+b\cdot\sin\theta\suchthat a,b\in\Re}$  
dengan variabel riil $\theta$ memiliki dimensi~$2$ karena ruang ini memiliki
basis $\sequence{\cos\theta,\sin\theta}$.
\end{example}

\begin{example}
Ruang trivial berdimensi nol karena basisnya kosong.
\end{example}

Sekali lagi, walaupun kadang-kadang kita mengatakan `berdimensi hingga' untuk
memberi penekanan, mulai sekarang kita menganggap semua ruang vektor
berdimensi hingga.
Jadi, dalam hasil berikutnya, kata `ruang' berarti `ruang vektor berdimensi
hingga'.

\begin{corollary} \label{cor:NoLiSetGreatDim}
%<*co:NoLiSetGreatDim>
Tidak ada himpunan bebas linear yang dapat berukuran lebih besar daripada
dimensi ruang yang memuatnya.
%</co:NoLiSetGreatDim>
\end{corollary}

\begin{proof}
%<*pf:NoLiSetGreatDim>
Pembuktian \nearbytheorem{th:AllBasesSameSize} tidak pernah menggunakan fakta
bahwa \( D \) merentang ruang, melainkan hanya bahwa \( D \) bebas linear.
%</pf:NoLiSetGreatDim>
\end{proof}


\begin{example} \label{ex:RefSubSpDiagram}
Ingat kembali diagram dari Contoh I.\ref{ex:SubspRThree} yang menunjukkan
subruang-subruang dari \( \Re^3 \).
Setiap subruang dideskripsikan dengan himpunan perentang minimal, yaitu suatu
basis.
Seluruh ruang memiliki basis dengan tiga anggota, subruang bidang memiliki
basis dengan dua anggota, subruang garis memiliki basis dengan satu anggota,
dan subruang trivial memiliki basis tanpa anggota.

Dalam bagian tersebut, kita belum dapat menunjukkan bahwa ini adalah semua
subruang \( \Re^3 \).
Sekarang kita dapat menunjukkannya.
Korolari sebelumnya membuktikan bahwa tidak ada, misalnya, subruang berdimensi
lima dari ruang berdimensi tiga.
Selain itu, menurut \nearbydefinition{df:Dimension}, dimensi setiap ruang
merupakan bilangan bulat, sehingga tidak ada subruang $\Re^3$ yang entah
bagaimana berdimensi $1.5$, di antara garis dan bidang.
Jadi, daftar subruang yang kita berikan bersifat lengkap; satu-satunya
subruang dari \( \Re^3 \) berdimensi tiga, dua, satu, atau nol.
\end{example}


\begin{corollary} \label{cor:LIExpBas}
%<*co:LIExpBas>
Sebarang himpunan bebas linear dapat diperluas menjadi suatu basis.
%</co:LIExpBas>
\end{corollary}

\begin{proof}
%<*pf:LIExpBas>
Jika suatu himpunan bebas linear belum merupakan basis, himpunan itu tidak
merentang ruang.
Menurut Lema II.\ref{lm:AddVecLiSetIsLiIffVecNotInSpan}, menambahkan suatu
vektor yang tidak berada dalam rentang himpunan akan mempertahankan kebebasan
linear.
Teruslah menambahkan vektor sampai himpunan yang dihasilkan merentang ruang;
korolari sebelumnya menunjukkan bahwa hal ini akan terjadi setelah hanya
berhingga banyak langkah.
%</pf:LIExpBas>
\end{proof}

\begin{corollary}  \label{co:SpanningSetShrinksToABasis}
%<*co:SpanningSetShrinksToABasis>
Sebarang himpunan perentang dapat diperkecil menjadi suatu basis.
%</co:SpanningSetShrinksToABasis>
\end{corollary}

\begin{proof}
%<*pf:SpanningSetShrinksToABasis>
Nyatakan himpunan perentang tersebut dengan \( S \).
Jika \( S \) kosong, maka \( S \) sudah merupakan basis (ruangnya harus berupa
ruang trivial).
Jika \( S=\set{\zero} \), himpunan tersebut dapat diperkecil menjadi basis
kosong, sehingga menjadi bebas linear tanpa mengubah rentangnya.

Jika tidak, $S$ memuat vektor $\vec{s}_1$ dengan $\vec{s}_1\neq\zero$ dan kita
dapat membentuk basis \( B_1=\sequence{\vec{s}_1} \).
Jika \( \spanof{B_1}=\spanof{S} \), pembuktian selesai.
Jika tidak, terdapat \( \vec{s}_2\in\spanof{S} \) sedemikian sehingga
\( \vec{s}_2\not\in\spanof{B_1} \).
Misalkan \( B_2=\sequence{\vec{s}_1,\vec{s_2}} \); menurut
Lema~II.\ref{lm:AddVecLiSetIsLiIffVecNotInSpan}, barisan ini bebas linear,
sehingga jika \( \spanof{B_2}=\spanof{S} \), pembuktian selesai.

Kita dapat mengulangi proses ini sampai rentangnya sama, yang pasti terjadi
dalam paling banyak berhingga banyak langkah.
%</pf:SpanningSetShrinksToABasis>
\end{proof}


\begin{corollary} \label{cor:NVecsRNSpanIffLI}
%<*co:NVecsRNSpanIffLI>
Dalam ruang berdimensi \( n \), suatu himpunan yang terdiri atas \( n \) vektor
bebas linear jika dan hanya jika merentang ruang tersebut.
%</co:NVecsRNSpanIffLI>
\end{corollary}

\begin{proof}
%<*pf:NVecsRNSpanIffLI>
Pertama, kita akan menunjukkan bahwa subhimpunan dengan \( n \) vektor bebas
linear jika dan hanya jika merupakan basis.
Arah `jika' jelas benar\Dash basis bersifat bebas linear.
Arah `hanya jika' berlaku karena suatu himpunan bebas linear dapat diperluas
menjadi basis, tetapi basis memiliki \( n \) elemen, sehingga hasil perluasan
ini sebenarnya adalah himpunan semula.

Untuk menyelesaikan pembuktian, kita akan menunjukkan bahwa sebarang
subhimpunan dengan \( n \) vektor merentang ruang jika dan hanya jika merupakan
basis.
Sekali lagi, arah `jika' jelas.
Arah `hanya jika' berlaku karena sebarang himpunan perentang dapat diperkecil
menjadi basis, tetapi basis memiliki \( n \) elemen, sehingga himpunan yang
diperkecil itu sama dengan himpunan semula.
%</pf:NVecsRNSpanIffLI>
\end{proof}

Hasil utama subbagian ini, yaitu bahwa semua basis dalam suatu ruang vektor
berdimensi hingga memiliki jumlah elemen yang sama, merupakan hasil terpenting
dalam buku ini.
Seperti ditunjukkan \nearbyexample{ex:RefSubSpDiagram}, hasil tersebut
menjelaskan ruang vektor dan subruang apa saja yang mungkin ada.

Salah satu konsekuensi langsung membawa kita kembali ke saat kita meninjau dua
makna yang mungkin bagi istilah `himpunan perentang minimal'.
Pada saat itu, kita mendefinisikan `minimal' sebagai bebas linear, tetapi kita
mencatat bahwa tafsiran lain yang juga masuk akal ialah bahwa suatu himpunan
perentang bersifat `minimal' ketika memiliki jumlah elemen paling sedikit di
antara semua himpunan dengan rentang yang sama.
Setelah menunjukkan bahwa semua basis memiliki jumlah elemen yang sama, kita
mengetahui bahwa kedua makna `minimal' tersebut ekuivalen.


\begin{exercises}
  \item[{\em Anggap bahwa semua ruang berdimensi hingga kecuali dinyatakan
    lain.}]
  \recommended \item  
    Carilah suatu basis bagi \(  \polyspace_2 \) dan dimensinya.
    \begin{answer}
      Salah satu basisnya adalah \( \sequence{1,x,x^2} \), sehingga dimensinya
      tiga.
    \end{answer}
  \item 
    Carilah suatu basis bagi himpunan solusi sistem berikut beserta dimensinya.
    \begin{equation*}
      \begin{linsys}{4}
         x_1  &-  &4x_2  &+  &3x_3  &-  &x_4  &=  &0  \\
        2x_1  &-  &8x_2  &+  &6x_3  &-  &2x_4 &=  &0  
      \end{linsys}
    \end{equation*}
    \begin{answer}
      Himpunan solusinya adalah
      \begin{equation*}
        \set{\colvec{4x_2-3x_3+x_4 \\ x_2 \\ x_3 \\ x_4}
               \suchthat x_2,x_3,x_4\in\Re}
      \end{equation*}
      sehingga basis alaminya adalah sebagai berikut
      \begin{equation*}
       \sequence{\colvec[r]{4 \\ 1 \\ 0 \\ 0},
                   \colvec[r]{-3 \\ 0 \\ 1 \\ 0},
                   \colvec[r]{1 \\ 0 \\ 0 \\ 1}  }
      \end{equation*}  
      (pemeriksaan kebebasan linearnya mudah).
      Jadi, dimensinya tiga.
    \end{answer}
  \recommended \item Carilah suatu basis bagi setiap ruang beserta dimensinya.
    \begin{exparts}
      \partsitem
        $
          \set{\colvec{x \\ y \\ z \\ w}\in\Re^4
               \suchthat x-w+z=0}
        $
      \partsitem himpunan matriks $\nbym{5}{5}$ yang semua entri tak nolnya
        berada pada diagonal (misalnya pada entri $1,1$ dan $2,2$, dan
        seterusnya)
      \partsitem 
        $\set{a_0+a_1x+a_2x^2+a_3x^3\suchthat 
            \text{$a_0+a_1=0$ dan $a_2-2a_3=0$}}\subseteq\polyspace_3$
    \end{exparts}
    \begin{answer}
      \begin{exparts}
        \partsitem
          Parametrisasikan untuk memperoleh deskripsi ruang berikut.
          \begin{equation*}
            \set{\colvec{w-z \\ y \\ z \\ w}
                 =\colvec{0 \\ 1 \\ 0 \\ 0}y+
                 \colvec{-1 \\ 0 \\ 1 \\ 0}z+
                 \colvec{1 \\ 0 \\ 0 \\ 1}w
                 \suchthat y,z,w\in\Re}
          \end{equation*}
          Hal ini memberikan ruang tersebut sebagai rentang himpunan tiga
          vektor.
          Untuk menunjukkan bahwa himpunan tiga vektor itu membentuk basis,
          kita periksa bahwa himpunan tersebut bebas linear.
          \begin{equation*}
             \colvec{0 \\ 0 \\ 0 \\ 0}
                 =\colvec{0 \\ 1 \\ 0 \\ 0}c_1+
                 \colvec{-1 \\ 0 \\ 1 \\ 0}c_2+
                 \colvec{1 \\ 0 \\ 0 \\ 1}c_3
          \end{equation*}
          Komponen kedua menghasilkan $c_1=0$, sedangkan komponen ketiga dan
          keempat menghasilkan $c_2=0$ dan~$c_3=0$.
          Jadi, salah satu basisnya adalah sebagai berikut.
          \begin{equation*}
              \sequence{
                 \colvec{0 \\ 1 \\ 0 \\ 0},
                 \colvec{-1 \\ 0 \\ 1 \\ 0},
                 \colvec{1 \\ 0 \\ 0 \\ 1} 
              } 
          \end{equation*}
          Dimensinya adalah jumlah vektor dalam basis, yaitu $3$.
        \partsitem
          Parametrisasi alaminya adalah sebagai berikut.
          \begin{multline*}
            \set{
              \begin{mat}
                a &0 &0 &0 &0 \\
                0 &b &0 &0 &0 \\
                0 &0 &c &0 &0 \\
                0 &0 &0 &d &0 \\
                0 &0 &0 &0 &e 
              \end{mat}
              \suchthat a,\ldots,e\in\Re}   \\
            \set{
              \begin{mat}
                1 &0 &0 &0 &0 \\
                0 &0 &0 &0 &0 \\
                0 &0 &0 &0 &0 \\
                0 &0 &0 &0 &0 \\
                0 &0 &0 &0 &0 
              \end{mat}\cdot a
              +\begin{mat}
                0 &0 &0 &0 &0 \\
                0 &1 &0 &0 &0 \\
                0 &0 &0 &0 &0 \\
                0 &0 &0 &0 &0 \\
                0 &0 &0 &0 &0 
              \end{mat}\cdot b
              +\cdots
              \suchthat a,\ldots,e\in\Re}
          \end{multline*}
          Pemeriksaan bahwa himpunan berelemen lima tersebut bebas linear
          bersifat langsung.
          Jadi, ini merupakan basis; dimensinya $5$.
          \begin{equation*}
            \sequence{
              \begin{mat}
                1 &0 &0 &0 &0 \\
                0 &0 &0 &0 &0 \\
                0 &0 &0 &0 &0 \\
                0 &0 &0 &0 &0 \\
                0 &0 &0 &0 &0 
              \end{mat},
              \begin{mat}
                0 &0 &0 &0 &0 \\
                0 &1 &0 &0 &0 \\
                0 &0 &0 &0 &0 \\
                0 &0 &0 &0 &0 \\
                0 &0 &0 &0 &0 
              \end{mat},
              \,\ldots\,,
              \begin{mat}
                0 &0 &0 &0 &0 \\
                0 &0 &0 &0 &0 \\
                0 &0 &0 &0 &0 \\
                0 &0 &0 &0 &0 \\
                0 &0 &0 &0 &1 
              \end{mat}
          }
          \end{equation*}
        \partsitem
          Pembatasan-pembatasan tersebut membentuk sistem linear dengan dua
          persamaan dan empat variabel tak diketahui.
          Parametrisasi sistem tersebut untuk menyatakan variabel utama dalam
          bentuk variabel bebas menghasilkan $a_0=-a_1$, $a_1=a_1$,
          $a_2=2a_3$, dan~$a_3=a_3$.
          \begin{multline*}
            \set{-a_1+a_1x+2a_3x^2+a_3x^3\suchthat a_1,a_3\in\Re}         \\
            =\set{(-1+x)\cdot a_1+(2x^2+x^3)\cdot a_3\suchthat a_1,a_3\in\Re}
          \end{multline*}
          Deskripsi tersebut menunjukkan bahwa ruang itu adalah rentang
          himpunan berelemen dua $\set{-1+x,2x^2+x^3}$.
          Pembuktian selesai setelah kita menunjukkan bahwa himpunan tersebut
          bebas linear.
          Relasi
          \begin{equation*}
            0+0x+0x^2+0x^3=(-1+x)\cdot c_1+(2x^2+x^3)\cdot c_2
          \end{equation*}
          menghasilkan $c_1=0$ dari suku-suku konstanta dan $c_2=0$ dari
          suku-suku kubik.
          Salah satu basis bagi ruang tersebut adalah
          $\sequence{-1+x,2x^2+x^3}$.
          Ruang ini berdimensi dua.

      \end{exparts}
    \end{answer}
  \item
    Carilah suatu basis bagi \( \matspace_{\nbyn{2}} \), yaitu ruang vektor
    matriks \( \nbyn{2} \), beserta dimensinya.
    \begin{answer}
      Untuk ruang ini,
      \begin{multline*}
        \set{\begin{mat}
               a  &b  \\
               c  &d
             \end{mat} \suchthat a,b,c,d\in\Re}            \\
        =\set{a\cdot\begin{mat}[r]
             1  &0  \\
             0  &0
           \end{mat}
           +\dots+
           d\cdot\begin{mat}[r]
             0  &0  \\
             0  &1
           \end{mat} \suchthat a,b,c,d\in\Re}
      \end{multline*}
      berikut adalah basis alami.
      \begin{equation*}
        \sequence{
           \begin{mat}[r]
             1  &0  \\
             0  &0
           \end{mat},
           \begin{mat}[r]
             0  &1  \\
             0  &0
           \end{mat},
           \begin{mat}[r]
             0  &0  \\
             1  &0
           \end{mat},
           \begin{mat}[r]
             0  &0  \\
             0  &1
           \end{mat}  }
      \end{equation*}
      Dimensinya empat.
    \end{answer}
  \item 
    Carilah dimensi ruang vektor matriks
    \begin{equation*}
      \begin{mat}
        a  &b  \\
        c  &d
      \end{mat}
    \end{equation*}
    yang memenuhi masing-masing syarat berikut.
    \begin{exparts}
      \item $a, b, c, d\in\Re$ 
      \item $a-b+2c=0$ dan~$d\in\Re$
      \item $a+b+c=0$, $a+b-c=0$, dan~$d\in\Re$
    \end{exparts}
    \begin{answer}
     \begin{exparts}
      \partsitem Seperti dalam latihan sebelumnya, ruang
        $\matspace_{\nbyn{2}}$ dari matriks tanpa pembatasan memiliki basis
        berikut
        \begin{equation*}
         \sequence{
           \begin{mat}[r]
             1  &0  \\
             0  &0
           \end{mat},
           \begin{mat}[r]
             0  &1  \\
             0  &0
           \end{mat},
           \begin{mat}[r]
             0  &0  \\
             1  &0
           \end{mat},
           \begin{mat}[r]
             0  &0  \\
             0  &1
           \end{mat}  }
        \end{equation*}
        sehingga dimensinya empat.
      \partsitem Untuk ruang ini,
        \begin{multline*}
         \set{\begin{mat}
               a  &b  \\
               c  &d
             \end{mat} \suchthat \text{$a=b-2c$ dan $d\in\Re$}}     \\
         =\set{b\cdot\begin{mat}[r]
             1  &1  \\
             0  &0
           \end{mat}
           +c\cdot\begin{mat}[r]
             -2  &0  \\
              1  &0
           \end{mat}
           +d\cdot\begin{mat}[r]
             0  &0  \\
             0  &1
           \end{mat} \suchthat b,c,d\in\Re}
        \end{multline*}
        berikut adalah basis alami.
        \begin{equation*}
          \sequence{
            \begin{mat}[r]
              1  &1  \\
              0  &0
            \end{mat},
            \begin{mat}[r]
              -2  &0  \\
               1  &0
            \end{mat},
            \begin{mat}[r]
              0  &0  \\
              0  &1
            \end{mat}  }
        \end{equation*}
        Dimensinya tiga.
      \partsitem Penerapan Metode Gauss pada sistem linear dua persamaan
        menghasilkan $c=0$ dan $a=-b$.
        Jadi, kita memiliki deskripsi berikut
        \begin{equation*}
         \set{\begin{mat}
               -b  &b  \\
                0  &d
             \end{mat} \suchthat b,d\in\Re}
         =\set{b\cdot\begin{mat}[r]
             -1  &1  \\
             0  &0
           \end{mat}
           +d\cdot\begin{mat}[r]
             0  &0  \\
             0  &1
           \end{mat} \suchthat b,d\in\Re}
        \end{equation*}
        sehingga berikut adalah basis alami.
        \begin{equation*}
          \sequence{
            \begin{mat}[r]
              -1  &1  \\
               0  &0
            \end{mat},
            \begin{mat}[r]
              0  &0  \\
              0  &1
            \end{mat}  }
        \end{equation*}
        Dimensinya dua.
     \end{exparts} 
    \end{answer}
  \recommended \item
    Carilah dimensi subruang dari $\Re^2$ berikut.
    \begin{equation*}
      S=\set{\colvec{a+b \\ a+c}\suchthat a,b,c\in\Re}
    \end{equation*}
    \begin{answer}
      Kita tidak dapat sekadar menghitung parameternya.
      Artinya, jawabannya bukan~$3$.
      Sebaliknya, perhatikan bahwa kita dapat menyatakan setiap anggota
      $\binom{x}{y}\in\Re^2$ dalam bentuk
      \begin{equation*}
        \colvec{x \\ y}=\colvec{a+b \\ a+c}  
      \end{equation*}
      dengan pilihan $a=0$, $b=x$, dan $c=y$
      (pilihan lain juga mungkin).
      Jadi, $S$ adalah himpunan $S=\Re^2$.
      Dimensinya~$2$.
    \end{answer}
  \recommended \item 
    Carilah dimensi masing-masing ruang berikut.
    \begin{exparts}
      \partsitem  Ruang polinomial kubik $p(x)$ sedemikian sehingga $p(7)=0$
      \partsitem  Ruang polinomial kubik $p(x)$ sedemikian sehingga
        $p(7)=0$ dan~$p(5)=0$
      \partsitem  Ruang polinomial kubik $p(x)$ sedemikian sehingga
        $p(7)=0$, $p(5)=0$, dan~$p(3)=0$
      \partsitem  Ruang polinomial kubik $p(x)$ sedemikian sehingga
        $p(7)=0$, $p(5)=0$, $p(3)=0$, dan~$p(1)=0$
    \end{exparts}
    \begin{answer}
      Basis bagi ruang-ruang ini dikembangkan dalam kumpulan jawaban subbagian
      sebelumnya.
      \begin{exparts}
        \partsitem Salah satu basisnya adalah
          \( \sequence{-7+x,-49+x^2,-343+x^3} \).
          Dimensinya tiga.
        \partsitem Salah satu basisnya adalah
          $\sequence{35-12x+x^2,420-109x+x^3}$, sehingga dimensinya dua.
        \partsitem Suatu basisnya adalah $\set{-105+71x-15x^2+x^3}$.
          Dimensinya satu.
        \partsitem Ini merupakan subruang trivial dari $\polyspace_3$, sehingga
          basisnya kosong.
          Dimensinya nol.
      \end{exparts}  
    \end{answer}
  \item 
     Berapakah dimensi rentang himpunan
     $\set{\cos^2\theta,\sin^2\theta,\cos2\theta,\sin2\theta}$?
     Rentang ini merupakan subruang dari ruang semua fungsi bernilai riil
     dengan satu variabel riil.
     \begin{answer}
       Pertama, ingat bahwa $\cos2\theta=\cos^2\theta-\sin^2\theta$, sehingga
       penghapusan $\cos2\theta$ dari himpunan ini tidak mengubah rentangnya.
       Himpunan yang tersisa,
       $\set{\cos^2\theta,\sin^2\theta,\sin2\theta}$, bebas linear
       (perhatikan relasi
       $c_1\cos^2\theta+c_2\sin^2\theta+c_3\sin2\theta=Z(\theta)$
       dengan $Z$ sebagai fungsi nol, lalu ambil $\theta=0$, $\theta=\pi/4$,
       dan $\theta=\pi/2$ untuk menyimpulkan bahwa setiap $c$ bernilai nol).
       Oleh karena itu, himpunan tersebut merupakan basis bagi rentangnya.
       Hal ini menunjukkan bahwa rentang tersebut adalah ruang vektor
       berdimensi tiga.
     \end{answer}
  \item  
    Carilah dimensi \( \C^{47} \), yaitu ruang vektor rangkap-$47$ bilangan
    kompleks.
    \begin{answer}
      Berikut adalah suatu basis
      \begin{equation*}
        \sequence{(1+0i,0+0i,\dots,0+0i),\,
                  (0+1i,0+0i,\dots,0+0i),(0+0i,1+0i,\dots,0+0i),\ldots }
      \end{equation*}   
      sehingga dimensinya \( 2\cdot 47=94 \).
    \end{answer}
  \item  
    Berapakah dimensi ruang vektor $\matspace_{\nbym{3}{5}}$ dari matriks
    \( \nbym{3}{5} \)?
    \begin{answer}
      Suatu basisnya adalah
      \begin{equation*}
        \sequence{
          \begin{mat}[r]
            1  &0  &0  &0  &0  \\
            0  &0  &0  &0  &0  \\
            0  &0  &0  &0  &0
          \end{mat},
          \begin{mat}[r]
            0  &1  &0  &0  &0  \\
            0  &0  &0  &0  &0  \\
            0  &0  &0  &0  &0
          \end{mat},
          \ldots,
          \begin{mat}[r]
            0  &0  &0  &0  &0  \\
            0  &0  &0  &0  &0  \\
            0  &0  &0  &0  &1
          \end{mat}  }
      \end{equation*}
      sehingga dimensinya \( 3\cdot 5=15 \).
    \end{answer}
  \recommended \item 
    Tunjukkan bahwa barisan berikut merupakan basis bagi $\Re^4$.
    \begin{equation*}
      \sequence{\colvec[r]{1 \\ 0 \\ 0 \\ 0},
        \colvec[r]{1 \\ 1 \\ 0 \\ 0},
        \colvec[r]{1 \\ 1 \\ 1 \\ 0},
        \colvec[r]{1 \\ 1 \\ 1 \\ 1} }
    \end{equation*}
    (Kita dapat menggunakan hasil-hasil subbagian ini untuk menyederhanakan
    pekerjaan tersebut.)
    \begin{answer}
       Dalam ruang berdimensi empat, suatu himpunan empat vektor bebas linear
       jika dan hanya jika merentang ruang tersebut.
       Bentuk vektor-vektor ini memudahkan pembuktian kebebasan linear
       (perhatikan persamaan komponen keempat, lalu persamaan komponen ketiga,
       dan seterusnya).
    \end{answer}
  \item Tentukan apakah masing-masing himpunan berikut merupakan basis bagi
    $\polyspace_{2}$.
    \begin{exparts*}
      \partsitem $\set{1,x^2,x^2-x}$
      \partsitem $\set{x^2+x,x^2-x}$
      \partsitem $\set{2x^2+x+1,2x+1,2}$
      \partsitem $\set{3x^2,-1,3x,x^2-x}$
    \end{exparts*}
    \begin{answer}
      Menurut hasil-hasil bagian ini, karena $\polyspace_2$ berdimensi~$3$,
      untuk menunjukkan bahwa suatu himpunan bebas linear merupakan basis kita
      hanya perlu mengamati bahwa himpunan tersebut memiliki tiga anggota.
      Untuk menunjukkan bahwa suatu himpunan bukan basis, kita hanya perlu
      mengamati bahwa himpunan tersebut tidak memiliki tiga anggota (dalam
      kasus ini, kita tidak perlu memeriksa kebebasan linear).
      \begin{exparts}
        \partsitem  Ini merupakan basis; dari pemeriksaan langsung, himpunan
           tersebut bebas linear (elemen pertama tidak memiliki suku kuadratik
           maupun linear, elemen kedua memiliki suku kuadratik tetapi tidak
           memiliki suku linear, dan elemen ketiga memiliki suku linear), dan
           memiliki tiga elemen.
        \partsitem Ini bukan basis karena hanya memiliki dua elemen.
        \partsitem Himpunan berelemen tiga ini merupakan basis.
        \partsitem Ini bukan basis karena memiliki empat elemen.
      \end{exparts}
    \end{answer}
  \item 
     Lihat \nearbyexample{ex:RefSubSpDiagram}.
     \begin{exparts}
       \partsitem Buatlah sketsa diagram subruang serupa bagi $\polyspace_2$.
       \partsitem Buatlah sketsa diagram serupa bagi $\matspace_{\nbyn{2}}$.
     \end{exparts}
     \begin{answer}
      \begin{exparts}
       \partsitem Diagram bagi $\polyspace_2$ memiliki empat tingkat.
          Tingkat teratas memuat satu-satunya subruang berdimensi tiga, yaitu
          $\polyspace_2$ sendiri.
          Tingkat berikutnya memuat subruang-subruang berdimensi dua
          (\emph{bukan} hanya polinomial linear; melainkan sebarang subruang
          berdimensi dua, seperti polinomial berbentuk $ax^2+b$).
          Di bawahnya terdapat subruang-subruang berdimensi satu.
          Terakhir, tentu saja, terdapat satu-satunya subruang berdimensi nol,
          yaitu subruang trivial.
        \partsitem Bagi $\matspace_{\nbyn{2}}$, diagram tersebut memiliki lima
          tingkat, yang mencakup subruang berdimensi empat hingga nol.
     \end{exparts}
    \end{answer}
  \recommended \item \label{exer:DimDomToR}
    Jika \( S \) adalah suatu himpunan, fungsi-fungsi
    \( \map{f}{S}{\Re} \) membentuk ruang vektor dengan operasi alami:
    jumlah $f+g$ adalah fungsi yang diberikan oleh
    \( f+g\,(s)=f(s)+g(s) \), dan hasil kali skalarnya adalah
    \( r\cdot f \, (s)=r\cdot f(s) \).
    Berapakah dimensi ruang yang dihasilkan bagi setiap domain?
    \begin{exparts*}
      \partsitem \( S=\set{1} \)
      \partsitem \( S=\set{1,2} \)
      \partsitem \( S=\set{1,\ldots,n} \)
    \end{exparts*}
    \begin{answer}
      \begin{exparts*}
        \partsitem Satu
        \partsitem Dua
        \partsitem \( n \)
      \end{exparts*}  
    \end{answer}
  \item  
    (Lihat \nearbyexercise{exer:DimDomToR}.)
    Buktikan bahwa himpunan semua fungsi \( \map{f}{\Re}{\Re} \) dengan operasi
    alami merupakan ruang berdimensi tak hingga.
    \begin{answer}
      Kita hanya perlu memberikan suatu himpunan bebas linear tak hingga.
      Salah satu barisan semacam itu adalah \( \sequence{f_1,f_2,\ldots} \),
      dengan \( \map{f_i}{\Re}{\Re} \) diberikan oleh
      \begin{equation*}
         f_i(x)=\begin{cases}
                   1  &\text{jika \( x=i \)}  \\
                   0  &\text{selainnya}
                \end{cases}
      \end{equation*}  
      yakni fungsi yang hanya bernilai $1$ pada $x=i$.
    \end{answer}
  \item  
    (Lihat \nearbyexercise{exer:DimDomToR}.)
    Berapakah dimensi ruang vektor fungsi $\map{f}{S}{\Re}$ dengan operasi
    alami, ketika domain $S$ adalah himpunan kosong?
    \begin{answer}
      Suatu fungsi adalah himpunan pasangan berurutan $(x,f(x))$.
      Jadi, hanya ada satu fungsi dengan domain kosong, yaitu himpunan kosong.
      Ruang vektor dengan hanya satu elemen merupakan ruang vektor trivial dan
      berdimensi nol.
    \end{answer}
  \item  
    Tunjukkan bahwa sebarang himpunan empat vektor dalam \( \Re^2 \) bergantung
    linear.
    \begin{answer}
      Terapkan \nearbycorollary{cor:NoLiSetGreatDim}.
    \end{answer}
  \item  
    Tunjukkan bahwa
    \( \sequence{\vec{\alpha}_1,\vec{\alpha}_2,\vec{\alpha}_3}\subset\Re^3 \)
    merupakan basis jika dan hanya jika tidak ada bidang melalui titik asal
    yang memuat ketiga vektor tersebut.
    \begin{answer}
      Suatu bidang berbentuk
      $\set{\vec{p}+t_1\vec{v}_1+t_2\vec{v}_2\suchthat t_1,t_2\in\Re}$.
      (Bab pertama juga menyebutnya `ruang datar berdimensi $2$', dan membahas
      mengapa deskripsi ini ekuivalen dengan deskripsi yang lazim dipakai dalam
      Kalkulus, yaitu himpunan titik $(x,y,z)$ yang memenuhi syarat berbentuk
      $ax+by+cz=d$.)
      Ketika bidang tersebut melalui titik asal, kita dapat mengambil vektor
      khusus $\vec{p}$ sama dengan $\zero$.
      Jadi, dalam bahasa yang telah kita kembangkan dalam bab ini, suatu bidang
      melalui titik asal adalah rentang suatu himpunan dua vektor.

      Sekarang, perhatikan pernyataannya.
      Menyatakan bahwa ketiga vektor tidak sebidang sama dengan menyatakan bahwa
      tidak satu pun vektor berada dalam rentang kedua vektor lainnya\Dash tidak
      satu pun merupakan kombinasi linear kedua vektor lainnya.
      Ini tepat berarti bahwa himpunan berelemen tiga tersebut bebas linear.
      Menurut \nearbycorollary{cor:NVecsRNSpanIffLI}, hal itu ekuivalen dengan
      pernyataan bahwa himpunan tersebut merupakan basis bagi $\Re^3$
      (lebih tepatnya, sebarang barisan yang dibentuk dari elemen-elemen
      himpunan tersebut merupakan basis).
    \end{answer}
  \item 
    % \begin{exparts}
      % \partsitem Prove that any subspace of a finite dimensional space
      %   has a basis. 
        Buktikan bahwa sebarang subruang dari suatu ruang berdimensi hingga juga
        berdimensi hingga.
    % \end{exparts}
    \begin{answer}
      Misalkan ruang $V$ berdimensi hingga dan $S$ merupakan subruang dari $V$.

      Jika $S$ tidak berdimensi hingga, maka $S$ memiliki himpunan bebas linear
      yang tak hingga (mulailah dengan himpunan kosong dan tambahkan berulang
      kali vektor yang tidak bergantung linear pada himpunan; proses ini dapat
      berlanjut dalam tak hingga banyak langkah, sebab jika berhenti maka $S$
      berdimensi hingga).
      Namun, sebarang subhimpunan bebas linear dari~$S$ juga merupakan
      subhimpunan bebas linear dari~$V$, yang bertentangan dengan
      \nearbycorollary{cor:NoLiSetGreatDim}.
      % \begin{exparts}
      %   \partsitem The empty set is a linearly independent subset of $S$.
      %     By \nearbycorollary{cor:LIExpBas}, it can be expanded to a basis
      %     for the vector space $S$. 
      %   \partsitem Any basis for the subspace $S$ is a linearly independent 
      %     set in the superspace $V$.
      %     Hence it can be expanded to a basis for the superspace, which is
      %     finite dimensional.
      %     Therefore it has only finitely many members.  
      % \end{exparts}
    \end{answer}
  \item  
    Di mana sifat hingga dari \( B \) digunakan dalam
    \nearbytheorem{th:AllBasesSameSize}?
    \begin{answer}
      Sifat itu memastikan bahwa semua \( \vec{\beta} \) akhirnya habis.
      Artinya, sifat tersebut membenarkan kalimat pertama paragraf terakhir.
    \end{answer}
  \item
    Buktikan bahwa jika \( U \) dan \( W \) keduanya merupakan subruang
    berdimensi tiga dari \( \Re^5 \), maka \( U\intersection W \) tak trivial.
    Generalisasikan hasilnya.
    \begin{answer}
       Misalkan \( B_U \) merupakan basis bagi \( U \) dan \( B_W \) merupakan
       basis bagi \( W \).
       Perhatikan konkatenasi kedua barisan basis tersebut.
       Keenam vektornya bergantung linear karena berada dalam ruang
       berdimensi lima \( \Re^5 \).
       Pilih suatu relasi linear tak trivial dan kelompokkan suku-sukunya
       menjadi \( \vec{u}+\vec{w}=\zero \), dengan \( \vec{u}\in U \) dan
       \( \vec{w}\in W \).
       Kebebasan linear masing-masing basis menyiratkan bahwa
       \( \vec{u}\neq\zero \) dan \( \vec{w}\neq\zero \).
       Dengan demikian, \( \vec{u}=-\vec{w} \) merupakan unsur tak nol dari
       \( U\intersection W \), sehingga irisannya tak trivial.

       Generalisasi:
       jika \( U,W \) merupakan subruang dari suatu ruang vektor berdimensi
       \( n \) dan jika \( \dim(U)+\dim(W)>n \), maka keduanya memiliki irisan
       tak trivial.
    \end{answer}
  \item 
    Suatu basis bagi sebuah ruang terdiri atas unsur-unsur ruang tersebut.
    Karena itu, wajar jika kita menyelidiki bagaimana sifat `merupakan basis'
    berinteraksi dengan operasi $\subseteq$, $\cap$, dan $\cup$.
    (Tentu saja, basis sebenarnya merupakan barisan, sehingga memiliki urutan,
    tetapi operasi-operasi ini dapat diperluas secara alami.)
    \begin{exparts}
      \partsitem Pertama, perhatikan bagaimana basis-basis dapat berhubungan
        melalui $\subseteq$.
        Andaikan \( U,W \) merupakan subruang dari suatu ruang vektor dan
        \( U\subseteq W \).
        Dapatkah ada basis \( B_U \) bagi \( U \) dan basis \( B_W \) bagi
        \( W \) sedemikian sehingga \( B_U\subseteq B_W \)?
        Haruskah basis-basis demikian ada?

        Untuk setiap basis \( B_U \) bagi \( U \), haruskah ada basis
        \( B_W \) bagi \( W \) sedemikian sehingga
        \( B_U\subseteq B_W \)?

        Untuk setiap basis \( B_W \) bagi \( W \), haruskah ada basis
        \( B_U \) bagi \( U \) sedemikian sehingga
        \( B_U\subseteq B_W \)?

        Untuk setiap basis \( B_U, B_W \) bagi \( U \) dan \( W \),
        haruskah \( B_U \) merupakan subhimpunan dari \( B_W \)?
      \partsitem Apakah $\cap$ dari basis-basis merupakan basis?
        Bagi ruang apa?
      \partsitem Apakah $\cup$ dari basis-basis merupakan basis?
        Bagi ruang apa?
      \partsitem Bagaimana dengan operasi komplemen?
     \end{exparts}
     (\textit{Petunjuk.}
     Ujilah setiap konjektur pada beberapa subruang dari \( \Re^3 \).)
     \begin{answer}
       Pertama, perhatikan bahwa suatu himpunan merupakan basis bagi suatu
       ruang jika dan hanya jika himpunan itu bebas linear, sebab dalam hal
       tersebut himpunan itu merupakan basis bagi rentangnya sendiri.
       \begin{exparts}
         \partsitem Jawaban atas pertanyaan pada paragraf kedua adalah
           ``ya'' (yang juga menghasilkan jawaban ``ya'' atas kedua
           pertanyaan pada paragraf pertama).
           Jika \( B_U \) merupakan basis bagi \( U \), maka
           \( B_U \) merupakan subhimpunan bebas linear dari \( W \).
           Terapkan \nearbycorollary{cor:LIExpBas} untuk memperluasnya menjadi
           basis bagi \( W \).
           Itulah \( B_W \) yang dicari.

           Jawaban atas pertanyaan pada paragraf ketiga adalah ``tidak'', yang
           juga menghasilkan jawaban ``tidak'' atas pertanyaan pada paragraf
           keempat.
           Berikut contoh basis bagi suatu superruang yang tidak memiliki
           subbasis yang membentuk basis bagi sebuah subruang: dalam
           \( W=\Re^2 \), perhatikan basis standar \( \stdbasis_2 \).
           Tidak ada subbasis dari $\stdbasis_2$ yang membentuk basis bagi
           subruang \( U \) dari $\Re^2$ berupa garis \( y=x \).
         \partsitem Irisan tersebut merupakan basis (bagi rentangnya) karena
           irisan himpunan-himpunan bebas linear juga bebas linear
           (irisan itu merupakan subhimpunan dari masing-masing himpunan
           bebas linear).

           Namun, irisan itu bukan basis bagi irisan ruang-ruangnya.
           Sebagai contoh, berikut dua basis bagi \( \Re^2 \):
           \begin{equation*}
             B_1=\sequence{\colvec[r]{1 \\ 0},\colvec[r]{0 \\ 1}}
             \quad\text{dan}\quad
             B_2=\sequence[r]{\colvec{2 \\ 0},\colvec[r]{0 \\ 2}}
           \end{equation*}
           dan \( \Re^2\intersection\Re^2=\Re^2 \), tetapi
           \( B_1\intersection B_2 \) kosong.
           Yang dapat kita katakan hanyalah bahwa $\cap$ dari basis-basis
           merupakan basis bagi rentangnya sendiri, yaitu subruang
           \( \spanof{B_1\intersection B_2} \) yang terkandung dalam irisan
           kedua ruang.
         \partsitem $\cup$ dari basis-basis tidak harus merupakan basis:
           dalam \( \Re^2 \)
           \begin{equation*}
             B_1=\sequence{\colvec[r]{1 \\ 0},\colvec[r]{1 \\ 1}}
             \quad\text{dan}\quad
             B_2=\sequence{\colvec[r]{1 \\ 0},\colvec[r]{0 \\ 2}}
           \end{equation*}
           \( B_1\union B_2 \) tidak bebas linear.
           Syarat perlu dan cukup agar $\cup$ dari dua basis merupakan basis
           \begin{equation*}
             B_1\union B_2 \text{ bebas linear }
             \quad\iff\quad
             \spanof{B_1\intersection B_2}=\spanof{B_1}\intersection
                                            \spanof{B_2}
           \end{equation*}
           cukup mudah dibuktikan (meskipun mungkin sukar diterapkan).
         \partsitem Komplemen suatu basis tidak dapat menjadi basis karena
           memuat vektor nol.
       \end{exparts}  
     \end{answer}
  \recommended \item 
    Selidiki bagaimana `dimensi' berinteraksi dengan `subhimpunan'.
    Andaikan \( U \) dan \( W \) keduanya merupakan subruang dari suatu ruang
    vektor, dan \( U\subseteq W \).
    \begin{exparts}
      \partsitem Buktikan bahwa \( \dim (U)\leq\dim (W) \).
      \partsitem Buktikan bahwa dimensinya sama jika dan hanya jika
        \( U=W \).
      \partsitem Tunjukkan bahwa butir sebelumnya tidak berlaku jika kedua
        ruang berdimensi tak hingga.
    \end{exparts}
    \begin{answer}
       \begin{exparts}
          \partsitem Suatu basis bagi \( U \) merupakan himpunan bebas linear
            dalam \( W \), sehingga melalui \nearbycorollary{cor:LIExpBas}
            basis itu dapat diperluas menjadi basis bagi \( W \).
            Basis kedua memiliki sekurang-kurangnya sebanyak anggota basis
            pertama.
          \partsitem Satu arah sudah jelas:~jika \( U=W \), keduanya memiliki
            dimensi yang sama.
            Untuk arah sebaliknya, misalkan \( B_U \) merupakan basis bagi
            \( U \).
            Basis itu merupakan subhimpunan bebas linear dari \( W \), sehingga
            dapat diperluas menjadi basis bagi \( W \).
            Jika \( \dim(U)=\dim(W) \), basis bagi \( W \) ini tidak memiliki
            lebih banyak anggota daripada \( B_U \), sehingga sama dengan
            \( B_U \).
            Karena \( U \) dan \( W \) memiliki basis yang sama, keduanya sama.
          \partsitem Misalkan \( W \) merupakan ruang polinomial berderajat
            hingga dan \( U \) merupakan subruang polinomial yang hanya
            memiliki suku-suku berpangkat genap.
            \begin{equation*}
               U=\set{a_0+a_1x^2+a_2x^4+\dots+a_nx^{2n}\suchthat
                          a_0,\ldots,a_n\in\Re} 
            \end{equation*}
            Kedua ruang berdimensi tak hingga, tetapi \( U \) merupakan
            subruang sejati.
       \end{exparts}  
     \end{answer}
    \item Berikut bukti alternatif bagi hasil utama bagian ini,
      \nearbytheorem{th:AllBasesSameSize}.
      Pertama diberikan sebuah contoh, kemudian sebuah lema, lalu teoremanya.
      \begin{exparts}
        \partsitem Nyatakan vektor dari $\Re^3$ ini sebagai kombinasi linear
          dari anggota-anggota basis yang diberikan.
         \begin{equation*}
           B=\sequence{\colvec[r]{1 \\ 0  \\ 0},
             \colvec[r]{1 \\ 1 \\ 0},
             \colvec[r]{0 \\ 0 \\ 2}}
           \qquad
           \vec{v}=\colvec[r]{1 \\ 2 \\ 0}
         \end{equation*}
        \partsitem Dalam kombinasi tersebut, pilih sebuah vektor basis yang
          koefisiennya tak nol.
          Ubah~$B$ dengan mengganti vektor basis itu dengan~$\vec{v}$ untuk
          memperoleh barisan baru~$\hat{B}$.
          Periksa bahwa $\hat{B}$ juga merupakan basis bagi~$\Re^3$.
        \partsitem (Lema Pertukaran)
          Andaikan
          \( B=\sequence{\vec{\beta}_1,\dots,\vec{\beta}_n} \) merupakan basis
          bagi suatu ruang vektor, dan untuk vektor \( \vec{v} \), hubungan
          \( \vec{v}=\lincombo{c}{\vec{\beta}} \) memiliki \( c_i\neq 0 \).
          Buktikan bahwa mengganti \( \vec{\beta}_i \) dengan \( \vec{v} \)
          menghasilkan basis lain bagi ruang tersebut.
        \partsitem Gunakan hasil itu bersama induksi untuk membuktikan
          \nearbytheorem{th:AllBasesSameSize}.
      \end{exparts}
      \begin{answer}
      \begin{exparts}
        \partsitem
          $\colvec[r]{1 \\ 2 \\ 0}
             =(-1)\cdot\colvec[r]{1 \\ 0  \\ 0}
              +2\colvec[r]{1 \\ 1 \\ 0}
              +0\cdot\colvec[r]{0 \\ 0 \\ 2}$
        \partsitem Dua di antara vektor-vektor basis memiliki koefisien tak nol.
           Misalnya, kita dapat memilih yang pertama.
           \begin{equation*}
             \hat{B}=\sequence{\colvec[r]{1 \\ 2  \\ 0},
               \colvec[r]{1 \\ 1 \\ 0},
               \colvec[r]{0 \\ 0 \\ 2}}
           \end{equation*}
           Pemeriksaan bahwa ini merupakan basis lain bagi \( \Re^3 \)
           bersifat rutin.
        \partsitem  
          Sebut hasil pertukaran itu
          \( \hat{B}=\sequence{\vec{\beta}_1,\dots,\vec{\beta}_{i-1},\vec{v},
             \vec{\beta}_{i+1},\dots,\vec{\beta}_n}  \).
          Pertama, kita tunjukkan bahwa $\hat{B}$ bebas linear.
          Setiap hubungan
          \( d_1\vec{\beta}_1+\dots+d_i\vec{v}+\dots+d_n\vec{\beta}_n=\zero \)
          di antara anggota-anggota $\hat{B}$, setelah substitusi untuk
          $\vec{v}$,
          \begin{equation*}
          d_1\vec{\beta}_1+\dots
          +d_i\cdot(c_1\vec{\beta}_1+\dots+c_i\vec{\beta}_i+\dots+c_n\vec{\beta}_n)
          +\dots+d_n\vec{\beta}_n
          =\zero
          \tag*{($*$)}\end{equation*}
          memberikan hubungan linear di antara anggota-anggota $B$.
          Basis $B$ bebas linear, sehingga koefisien $d_ic_i$ dari
          $\vec{\beta}_i$ bernilai nol.
          Karena kita mengandaikan $c_i$ tak nol, $d_i=0$.
          Dengan menggunakan hal ini dalam persamaan~$(*)$, diperoleh bahwa
          semua $d$ lainnya juga nol.
          Jadi, $\hat{B}$ bebas linear.

          Kita selesaikan bukti dengan menunjukkan bahwa $\hat{B}$ memiliki
          rentang yang sama dengan $B$.
          Separuh dari argumen ini, yakni
          $\spanof{\smash{\hat{B}}}\subseteq\spanof{B}$, 
          mudah; setiap anggota
          $d_1\vec{\beta}_1+\dots+d_i\vec{v}+\dots+d_n\vec{\beta}_n$
          dari $\spanof{\smash{\hat{B}}}$ dapat ditulis sebagai
          $
          d_1\vec{\beta}_1+\dots
          +d_i\cdot(c_1\vec{\beta}_1+\dots+c_n\vec{\beta}_n)
          +\dots+d_n\vec{\beta}_n
          $,
          yang merupakan kombinasi linear dari kombinasi-kombinasi linear
          anggota $B$, sehingga berada dalam $\spanof{B}$.
          Untuk separuh argumen
          $\spanof{B}\subseteq\spanof{\smash{\hat{B}}}$, ingat bahwa jika
          $\vec{v}=c_1\vec{\beta}_1+\dots+c_n\vec{\beta}_n$ dengan $c_i\neq 0$
          maka persamaan tersebut dapat disusun ulang menjadi
          $\vec{\beta}_i=(-c_1/c_i)\vec{\beta}_1+\dots+(1/c_i)\vec{v}+\dots
          +(-c_n/c_i)\vec{\beta}_n$.
          Sekarang, ambil sebarang anggota
          $d_1\vec{\beta}_1+\dots+d_i\vec{\beta}_i+\dots+d_n\vec{\beta}_n$
          dari $\spanof{B}$, substitusikan bentuk $\vec{\beta}_i$ sebagai
          kombinasi linear anggota-anggota $\hat{B}$, lalu perhatikan, seperti
          pada separuh pertama argumen ini, bahwa hasilnya merupakan kombinasi
          linear dari kombinasi-kombinasi linear anggota $\hat{B}$, sehingga
          berada dalam $\spanof{\smash{\hat{B}}}$.
        \partsitem Tetapkan suatu ruang vektor yang memiliki sekurang-kurangnya
          satu basis hingga.
          Dari semua basis ruang ini, pilih satu basis
          \( B=\sequence{\vec{\beta}_1,\dots,\vec{\beta}_n} \) yang ukurannya
          minimum.
          Kita akan menunjukkan bahwa setiap basis lain
          \( D=\smash{\sequence{\vec{\delta}_1,\vec{\delta}_2,\ldots}} \)
          juga memiliki jumlah anggota yang sama, yaitu $n$.
          Karena \( B \) berukuran minimum, \( D \) memiliki tidak kurang dari
          \( n \)~vektor.
          Kita akan membuktikan bahwa \( D \) tidak mungkin memiliki lebih dari
          \( n \) vektor.

          Basis \( B \) merentang ruang tersebut dan \( \vec{\delta}_1 \)
          berada dalam ruang itu, sehingga \( \vec{\delta}_1 \) merupakan
          kombinasi linear tak trivial dari anggota-anggota \( B \).
          Menurut Lema Pertukaran, kita dapat mengganti sebuah vektor dari
          \( B \) dengan \( \vec{\delta}_1 \), sehingga menghasilkan basis
          \( B_1 \), yang salah satu anggotanya adalah \( \vec{\delta}_1 \)
          dan semua \( n-1 \) anggota lainnya adalah vektor
          \( \vec{\beta} \).

          Paragraf sebelumnya merupakan langkah basis bagi argumen induksi.
          Langkah induksi dimulai dengan basis \( B_k \)
          (untuk \( 1\leq k<n \)) yang memuat \( k \) anggota \( D \) dan
          \( n-k \) anggota \( B \).
          Kita tahu bahwa \( D \) memiliki sekurang-kurangnya \( n \) anggota,
          sehingga \( \vec{\delta}_{k+1} \) ada.
          Nyatakan vektor itu sebagai kombinasi linear anggota-anggota
          \( B_k \).
          Pokok pentingnya:~dalam representasi tersebut, sekurang-kurangnya
          satu skalar tak nol harus berkaitan dengan suatu
          \( \vec{\beta}_i \); jika tidak, representasi itu akan menjadi
          hubungan linear tak trivial di antara anggota-anggota himpunan bebas
          linear \( D \).
          Ganti \( \vec{\beta}_i \) dengan \( \vec{\delta}_{k+1} \) untuk
          memperoleh basis baru \( B_{k+1} \), dengan satu vektor
          \( \vec{\delta} \) lebih banyak dan satu vektor \( \vec{\beta} \)
          lebih sedikit daripada basis sebelumnya \( B_k \).

          Ulangi proses itu sampai tidak ada vektor \( \vec{\beta} \) yang
          tersisa, sehingga \( B_n \) memuat
          $\vec{\delta}_1,\dots,\vec{\delta}_n$.
          Sekarang, \( D \) tidak dapat memiliki lebih dari \( n \) vektor ini,
          sebab setiap \( \vec{\delta}_{n+1} \) yang masih tersisa akan berada
          dalam rentang \( B_n \) (karena \( B_n \) merupakan basis), sehingga
          menjadi kombinasi linear dari vektor-vektor $\vec{\delta}$ lainnya;
          ini bertentangan dengan kebebasan linear $D$.
      \end{exparts}
      \end{answer}
  \puzzle \item 
    \cite{Sheffer}
    Sebuah perpustakaan memiliki $n$ buku dan $n+1$ pelanggan.
    Setiap pelanggan telah membaca sekurang-kurangnya satu buku dari
    perpustakaan tersebut.
    Buktikan bahwa pasti ada dua himpunan pelanggan yang saling lepas dan telah
    membaca buku-buku yang persis sama (yakni, gabungan buku yang dibaca oleh
    para pelanggan dalam masing-masing himpunan sama).
    \begin{answer}
      (Jawaban ini berasal dari komentar situs oleh Yuzhou Gu.)
      Untuk setiap orang, tetapkan sebuah vektor $\vec{v}_i\in\R^n$, dengan
      entri bernilai $1$ jika orang tersebut membaca buku yang bersangkutan,
      dan $0$ jika tidak.
      Karena vektor-vektor itu bergantung linear, terdapat persamaan berbentuk
      $\sum a_i v_i = 0$.
      Maka himpunan semua $i$ dengan $a_i>0$ dan himpunan semua $i$ dengan
      $a_i<0$ merupakan dua himpunan saling lepas dengan gabungan buku yang
      dibaca sama.
    \end{answer}
  \puzzle \item 
    \cite{Wohascum47}
    Untuk sebarang vektor $\vec{v}$ dalam $\Re^n$ dan sebarang permutasi
    $\sigma$ dari bilangan $1$, $2$, \ldots, $n$
    (yakni, $\sigma$ merupakan penyusunan kembali bilangan-bilangan tersebut
    dalam urutan baru), definisikan $\sigma(\vec{v})$ sebagai vektor yang
    komponen-komponennya adalah $v_{\sigma(1)}$, $v_{\sigma(2)}$, \ldots, dan
    $v_{\sigma(n)}$ (dengan $\sigma(1)$ sebagai bilangan pertama dalam susunan
    baru itu, dan seterusnya).
    Sekarang tetapkan $\vec{v}$ dan misalkan $V$ merupakan rentang dari
    $\set{\sigma(\vec{v})\suchthat \mbox{$\sigma$ mempermutasikan $1$, \ldots, $n$}}$.
    Apa saja kemungkinan dimensi $V$?
    \begin{answer}
      Kemungkinan dimensi $V$ adalah $0$, $1$, $n-1$, dan $n$.

      Untuk melihatnya, pertama perhatikan kasus ketika semua koordinat
      $\vec{v}$ sama.
      \begin{equation*}
         \vec{v}=\colvec{z \\ z \\ \vdots \\ z}
      \end{equation*}
      Maka $\sigma(\vec{v})=\vec{v}$ untuk setiap permutasi $\sigma$, sehingga
      $V$ hanyalah rentang dari $\vec{v}$, yang berdimensi $0$ atau $1$
      menurut apakah $\vec{v}$ sama dengan $\zero$ atau tidak.

      Sekarang, andaikan tidak semua koordinat $\vec{v}$ sama; misalkan $x$
      dan $y$, dengan $x\neq y$, merupakan dua di antara koordinat
      $\vec{v}$.
      Maka kita dapat menemukan permutasi $\sigma_1$ dan $\sigma_2$ sedemikian
      sehingga
      \begin{equation*}
        \sigma_1(\vec{v})=\colvec{x \\ y \\ a_3 \\ \vdots \\ a_n}
        \quad\text{dan}\quad
        \sigma_2(\vec{v})=\colvec{y \\ x \\ a_3 \\ \vdots \\ a_n}
      \end{equation*}
      untuk suatu $a_3,\ldots,a_n\in\Re$.
      Oleh karena itu,
      \begin{equation*}
        \frac{1}{y-x}\bigl(\sigma_1(\vec{v})-\sigma_2(\vec{v})\bigr)=
        \colvec[r]{-1 \\ 1 \\ 0 \\ \vdotswithin{-1} \\ 0}
      \end{equation*}
      berada dalam $V$.
      Dengan kata lain, $\vec{e}_2-\vec{e}_1\in V$, dengan $\vec{e}_1$,
      $\vec{e}_2$, \ldots, $\vec{e}_n$ sebagai basis standar bagi $\Re^n$.
      Demikian pula, $\vec{e}_3-\vec{e}_2$, \ldots,
      $\vec{e}_n-\vec{e}_1$ semuanya berada dalam $V$.
      Mudah dilihat bahwa vektor-vektor
      $\vec{e}_2-\vec{e}_1$, $\vec{e}_3-\vec{e}_2$, \ldots, 
      $\vec{e}_n-\vec{e}_1$ bebas linear (yakni, membentuk himpunan bebas
      linear), sehingga $\dim V\geq n-1$.

      Terakhir, kita dapat menulis
      \begin{align*}
        \vec{v} &=x_1\vec{e}_1+x_2\vec{e}_2+\dots+x_n\vec{e}_n  \\
                &=(x_1+x_2+\dots+x_n)\vec{e}_1+x_2(\vec{e}_2-\vec{e}_1)+\dots
                   +x_n(\vec{e}_n-\vec{e}_1)
      \end{align*}
      Ini menunjukkan bahwa jika $x_1+x_2+\dots+x_n=0$, maka $\vec{v}$ berada
      dalam rentang dari $\vec{e}_2-\vec{e}_1$, \ldots,
      $\vec{e}_n-\vec{e}_1$ (yakni, dalam rentang himpunan vektor-vektor
      tersebut); demikian pula, setiap $\sigma(\vec{v})$ berada dalam rentang
      ini, sehingga $V$ sama dengan rentang tersebut dan $\dim V=n-1$.
      Di sisi lain, jika $x_1+x_2+\cdots+x_n\neq 0$, persamaan di atas
      menunjukkan bahwa $\vec{e}_1\in V$, sehingga
      $\vec{e}_1,\dots,\vec{e}_n\in V$; jadi $V=\Re^n$ dan $\dim V=n$.
    \end{answer}
\end{exercises}
\index{basis|)}



























\subsection{Ruang Vektor dan Sistem Linear}
Sekarang kita akan meninjau kembali sistem linear dan Metode Gauss dengan
bantuan perangkat dan istilah dari bab ini.
Kita akan membahas tiga pokok.

Untuk pokok pertama, ingat kembali wawasan dari Bab Satu bahwa Metode Gauss
bekerja dengan mengambil kombinasi linear baris-baris\Dash
jika dua matriks dihubungkan oleh operasi baris
$A\longrightarrow\cdots\longrightarrow B$, maka setiap baris $B$ merupakan
kombinasi linear dari baris-baris $A$.
Karena itu, ruang vektor berikut merupakan lingkungan yang tepat untuk
mempelajari operasi baris secara umum dan Metode Gauss secara khusus.

\begin{definition} \label{df:RowSpace}
%<*df:RowSpace>
\definend{Ruang baris\/}\index{matriks!ruang baris}\index{ruang baris}
suatu matriks adalah rentang dari himpunan baris-barisnya.
\definend{Rank baris\/}\index{baris!rank}\index{matriks!rank baris}\index{rank baris}
adalah dimensi ruang tersebut, yakni jumlah baris yang bebas linear.
%</df:RowSpace>
\end{definition}

\begin{example}
Jika
\begin{equation*}
  A=\begin{mat}[r]
          2  &3  \\
          4  &6
       \end{mat}
\end{equation*}
maka \( \rowspace{A} \) merupakan subruang berikut dari ruang vektor baris
berkomponen dua.
\begin{equation*}
   \set{c_1\cdot\rowvec{2 &3}+c_2\cdot\rowvec{4  &6}
          \suchthat c_1,c_2\in\Re }
\end{equation*}
Vektor baris kedua bergantung linear pada yang pertama, sehingga deskripsi di
atas dapat disederhanakan menjadi
$\set{c\cdot\rowvec{2 &3}\suchthat c\in\Re }$.
\end{example}

\begin{lemma}     \label{le:RowSpUnchByGR}
%<*lm:RowSpUnchByGR>
Jika dua matriks \( A \) dan \( B \) dihubungkan oleh suatu operasi baris
\begin{equation*}
  A\grstep{\rho_i\leftrightarrow\rho_j}B 
  \quad\text{atau}\quad
  A\grstep{k\rho_i}B 
  \quad\text{atau}\quad
  A\grstep{k\rho_i+\rho_j}B
\end{equation*}
(untuk $i\neq j$ dan $k\neq 0$), maka ruang baris keduanya sama.
Jadi, matriks-matriks yang ekuivalen baris memiliki ruang baris yang sama dan,
karena itu, rank baris yang sama.
%</lm:RowSpUnchByGR>
\end{lemma}

\begin{proof}
%<*pf:RowSpUnchByGR0>
Korolari Satu.III.\ref{cor:RowsOfEqMatsLinCombos}
menunjukkan bahwa jika \mbox{$A\longrightarrow B$}, maka setiap baris $B$
merupakan kombinasi linear dari baris-baris \( A \).
Dalam istilah di atas, ini berarti bahwa setiap baris \( B \) merupakan unsur
dari ruang baris $A$.
Selanjutnya, $\rowspace{B}\subseteq\rowspace{A}$ karena suatu anggota
$\rowspace{B}$ merupakan kombinasi linear dari baris-baris $B$, sehingga
merupakan kombinasi dari kombinasi-kombinasi baris $A$ dan, menurut Lema
Kombinasi Linear, juga merupakan anggota $\rowspace{A}$.
%</pf:RowSpUnchByGR0>

%<*pf:RowSpUnchByGR1>
Untuk inklusi himpunan yang lain, ingat Lema
Satu.III.\ref{le:RowOpsRev}, bahwa operasi baris dapat dibalik, sehingga
\mbox{$A\longrightarrow B$} jika dan hanya jika
\mbox{$B\longrightarrow A$}.
Maka $\rowspace{A}\subseteq\rowspace{B}$ diperoleh seperti pada paragraf
sebelumnya.
%</pf:RowSpUnchByGR1>
% Thus the two sets are equal.
\end{proof}

Tentu saja, Metode Gauss menjalankan operasi baris secara sistematis dengan
tujuan mencapai bentuk eselon.

\begin{lemma}  \label{le:RowsEchMatLI}
%<*lm:RowsEchMatLI>
Baris-baris tak nol dari suatu matriks bentuk eselon membentuk himpunan bebas
linear.
%</lm:RowsEchMatLI>
\end{lemma}

\begin{proof}
%<*pf:RowsEchMatLI>
Lema~Satu.III.\ref{le:EchFormNoLinCombo}
menyatakan bahwa tidak ada baris tak nol dari suatu matriks bentuk eselon yang
merupakan kombinasi linear dari baris-baris lainnya.
Hasil ini menyatakan ulang fakta tersebut dengan istilah bab ini.
%</pf:RowsEchMatLI>
\end{proof}

Jadi, dalam bahasa bab ini, reduksi Gauss bekerja dengan menghilangkan
kebergantungan linear di antara baris-baris tanpa mengubah rentangnya, sampai
tidak ada lagi hubungan linear tak trivial di antara baris-baris tak nol.
Singkatnya, Metode Gauss menghasilkan basis bagi ruang baris.

\begin{example}
Dari sebarang matriks, kita dapat menghasilkan basis bagi ruang baris dengan
menjalankan Metode Gauss lalu mengambil baris-baris tak nol dari matriks bentuk
eselon yang dihasilkan.
Sebagai contoh,
\begin{equation*}
  \begin{mat}[r]
    1  &3  &1  \\
    1  &4  &1  \\
    2  &0  &5
  \end{mat}
  \grstep[-2\rho_1+\rho_3]{-\rho_1+\rho_2}
  \repeatedgrstep{6\rho_2+\rho_3}
  \begin{mat}[r]
    1  &3  &1  \\
    0  &1  &0  \\
    0  &0  &3
  \end{mat}
\end{equation*}
menghasilkan basis $\sequence{\rowvec{1 &3 &1},
            \rowvec{0 &1 &0},
            \rowvec{0 &0 &3} }$ bagi ruang baris.
Ini merupakan basis bagi ruang baris matriks awal maupun matriks akhir karena
kedua ruang baris tersebut sama.
\end{example}

Dengan teknik ini, kita juga dapat menemukan basis bagi rentang yang tidak
secara langsung melibatkan vektor baris.

\begin{definition} \label{df:ColumnSpace}
%<*df:ColumnSpace>
\definend{Ruang kolom}\index{kolom!ruang}\index{matriks!ruang kolom}
suatu matriks adalah rentang dari himpunan kolom-kolomnya.
\definend{Rank kolom}\index{kolom!rank}\index{rank!kolom}
adalah dimensi ruang kolom, yakni jumlah kolom yang bebas linear.
%</df:ColumnSpace>
\end{definition}

Minat kita pada ruang kolom berasal dari kajian kita tentang sistem linear.
Sebagai contoh, sistem berikut
\begin{equation*}
  \begin{linsys}{3}
    c_1  &+  &3c_2  &+  &7c_3  &=  &d_1  \\
   2c_1  &+  &3c_2  &+  &8c_3  &=  &d_2  \\
         &   &c_2   &+  &2c_3  &=  &d_3  \\
   4c_1  &   &      &+  &4c_3  &=  &d_4   
  \end{linsys}
\end{equation*}
memiliki solusi jika dan hanya jika vektor yang tersusun dari semua \( d \)
merupakan kombinasi linear dari vektor-vektor kolom lainnya,
\begin{equation*}
  c_1\colvec[r]{1 \\ 2 \\ 0 \\ 4}
  +c_2\colvec[r]{3 \\ 3 \\ 1 \\ 0}
  +c_3\colvec[r]{7 \\ 8 \\ 2 \\ 4}
  =\colvec{d_1 \\ d_2 \\ d_3 \\ d_4}
\end{equation*}
yang berarti bahwa vektor yang tersusun dari semua \( d \) berada dalam ruang
kolom matriks koefisien.

\begin{example}   \label{ex:BasisForColSpace}
Diberikan matriks berikut,
\begin{equation*}
  \begin{mat}[r]
    1  &3  &7  \\
    2  &3  &8  \\
    0  &1  &2  \\
    4  &0  &4
  \end{mat}
\end{equation*}
untuk memperoleh basis bagi ruang kolom, ubah sementara kolom-kolomnya menjadi
baris lalu lakukan reduksi.
\begin{equation*}
    \begin{mat}[r]
       1  &2  &0  &4  \\
       3  &3  &1  &0  \\
       7  &8  &2  &4
    \end{mat}
  \grstep[-7\rho_1+\rho_3]{-3\rho_1+\rho_2}
  \repeatedgrstep{-2\rho_2+\rho_3}
  \begin{mat}[r]
     1  &2  &0  &4  \\
     0  &-3 &1  &-12\\
     0  &0  &0  &0
  \end{mat}
\end{equation*}
Sekarang ubah kembali baris-baris itu menjadi kolom.
\begin{equation*}
  \sequence{
     \colvec[r]{1 \\ 2 \\ 0 \\ 4},
     \colvec[r]{0 \\ -3 \\ 1 \\ -12} }
\end{equation*}
Hasilnya merupakan basis bagi ruang kolom matriks yang diberikan.
\end{example}

\begin{definition} \label{df:Transpose}
%<*df:Transpose>
\definend{Transpos}\index{transpos}\index{matriks!transpos}
suatu matriks adalah hasil dari mempertukarkan baris-baris dan kolom-kolomnya,
sehingga kolom~\( j \) dari matriks \( A \) menjadi baris~\( j \) dari
\( \trans{A} \), dan sebaliknya.
%</df:Transpose>
\end{definition}
\noindent Jadi, contoh sebelumnya dapat diringkas sebagai ``transposkan,
reduksi, lalu transposkan kembali.''

Dengan menerima untuk sementara gagasan yang masih samar bahwa ruang-ruang
vektor dapat dianggap ``sama'', kita bahkan dapat menggunakan Metode Gauss
untuk menemukan basis bagi rentang dalam jenis ruang vektor lain.

\begin{example}
Untuk memperoleh basis bagi rentang dari
\( \set{x^2+x^4,2x^2+3x^4,-x^2-3x^4} \) dalam ruang
\( \polyspace_4 \), anggap ketiga polinomial ini ``sama'' dengan vektor baris
\( \rowvec{0 &0 &1 &0 &1} \),
\( \rowvec{0 &0 &2 &0 &3} \), dan
\( \rowvec{0 &0 &-1 &0 &-3} \), lalu terapkan Metode Gauss
\begin{equation*}
  \begin{mat}[r]
    0  &0  &1  &0  &1  \\
    0  &0  &2  &0  &3  \\
    0  &0  &-1 &0  &-3
  \end{mat}
  \grstep[\rho_1+\rho_3]{-2\rho_1+\rho_2}
  \repeatedgrstep{2\rho_2+\rho_3}
  \begin{mat}[r]
    0  &0  &1  &0  &1  \\
    0  &0  &0  &0  &1  \\
    0  &0  &0  &0  &0
  \end{mat}
\end{equation*}
dan terjemahkan kembali untuk memperoleh basis
\( \sequence{x^2+x^4,x^4} \).
(Seperti disebutkan sebelumnya, kita akan membuat frasa ``sama'' ini tepat pada
awal bab berikutnya.)
\end{example}

Jadi, pokok pertama subbagian ini adalah bahwa perangkat dalam bab ini memberi
kita pemahaman yang lebih konseptual tentang reduksi Gauss.

Untuk pokok kedua, perhatikan bahwa operasi baris pada suatu matriks dapat
mengubah ruang kolomnya.
\begin{equation*}
  \begin{mat}[r]
    1  &2  \\
    2  &4
  \end{mat}
  \grstep{-2\rho_1+\rho_2}
  \begin{mat}[r]
    1  &2  \\
    0  &0
  \end{mat}
\end{equation*}
Ruang kolom matriks di sebelah kiri memuat vektor-vektor yang komponen keduanya
tak nol, sedangkan ruang kolom matriks di sebelah kanan hanya memuat
vektor-vektor yang komponen keduanya nol; jadi, kedua ruang itu berbeda.
Pengamatan ini membuat hasil berikut mengejutkan.

\begin{lemma} \label{le:RowOpsNoChngColRnk}
%<*lm:RowOpsNoChngColRnk>
Operasi baris tidak mengubah rank kolom.
%</lm:RowOpsNoChngColRnk>
\end{lemma}

\begin{proof}
%<*pf:RowOpsNoChngColRnk>
Dengan kata lain, jika $A$ tereduksi menjadi $B$, maka rank kolom $B$ sama
dengan rank kolom $A$.

Bukti ini selesai jika kita menunjukkan bahwa operasi baris tidak memengaruhi
hubungan linear di antara kolom-kolom, sebab rank kolom adalah ukuran himpunan
kolom-kolom tak berhubungan yang terbesar.
Dengan kata lain, kita akan menunjukkan bahwa suatu hubungan di antara
kolom-kolom (misalnya, kolom kelima sama dengan dua kali kolom kedua ditambah
kolom keempat) berlaku jika dan hanya jika hubungan itu tetap berlaku setelah
operasi baris.
Namun, inilah tepatnya teorema pertama buku ini, Teorema
Satu.I.\ref{th:GaussMethod}:~dalam suatu hubungan di antara kolom-kolom,
\begin{equation*}
  c_1\cdot\colvec{a_{1,1} \\ a_{2,1} \\ \vdots \\ a_{m,1}}
   +\dots+
  c_n\cdot \colvec{a_{1,n} \\ a_{2,n} \\ \vdots \\ a_{m,n}}
   =\colvec[r]{0 \\ 0 \\ \vdotswithin{0} \\ 0}
\end{equation*}
operasi baris tidak mengubah himpunan solusi \( (c_1,\ldots,c_n) \).
%</pf:RowOpsNoChngColRnk>
\end{proof}

Cara lain untuk menunjukkan bahwa Metode Gauss memberi informasi tentang ruang
kolom maupun ruang baris ialah melalui reduksi Gauss--Jordan.
Reduksi ini berakhir pada bentuk eselon tereduksi suatu matriks, seperti di
bawah ini.
\begin{equation*}
  \begin{mat}[r]
    1  &3  &1  &6  \\
    2  &6  &3  &16 \\
    1  &3  &1  &6
  \end{mat}
  \grstep{}\cdots\grstep{}
  \begin{mat}[r]
    1  &3  &0  &2  \\
    0  &0  &1  &4  \\
    0  &0  &0  &0
  \end{mat}
\end{equation*}
Perhatikan ruang baris dan ruang kolom hasil ini.

Pokok pertama yang dibahas sebelumnya dalam subbagian ini menyatakan bahwa
untuk memperoleh basis bagi ruang baris, kita cukup mengumpulkan baris-baris
yang memiliki entri utama.
Namun, karena matriks ini berbentuk eselon tereduksi, basis bagi ruang kolom
dapat diperoleh dengan sama mudahnya:~kumpulkan kolom-kolom yang memuat entri
utama,
\( \sequence{\vec{e}_1,\vec{e}_2} \).
% (Linear independence is obvious.
% As to span, the other columns are in the span 
% since they all have a third component of zero.)
Jadi, untuk matriks berbentuk eselon tereduksi, kita dapat menemukan basis bagi
ruang baris dan ruang kolom dengan cara yang pada dasarnya sama, yaitu mengambil
bagian-bagian matriks, baik baris maupun kolom, yang memuat entri utama.

\begin{theorem} \label{th:RowRankEqualsColumnRank}
%<*th:RowRankEqualsColumnRank>
Untuk sebarang matriks, rank baris dan rank kolomnya sama.
%</th:RowRankEqualsColumnRank>
\end{theorem}

\begin{proof}
%<*pf:RowRankEqualsColumnRank0>
Bawa matriks tersebut ke bentuk eselon tereduksi.
Rank barisnya sama dengan jumlah entri utama karena jumlah itu sama dengan
jumlah baris tak nol.
Jumlah entri utama juga sama dengan rank kolom, sebab himpunan kolom yang
memuat entri utama terdiri atas beberapa vektor \( \vec{e}_i \) dari suatu
basis standar, dan himpunan itu bebas linear serta merentang himpunan kolom.
Jadi, pada matriks berbentuk eselon tereduksi, rank baris sama dengan rank
kolom karena keduanya sama dengan jumlah entri utama.
%</pf:RowRankEqualsColumnRank0>

%<*pf:RowRankEqualsColumnRank1>
Namun, \nearbylemma{le:RowSpUnchByGR} dan
\nearbylemma{le:RowOpsNoChngColRnk} menunjukkan bahwa penggunaan operasi baris
untuk mencapai bentuk eselon tereduksi tidak mengubah rank baris maupun rank
kolom.
Jadi, rank baris dan rank kolom matriks semula juga sama.
%</pf:RowRankEqualsColumnRank1>
\end{proof}

\begin{definition} \label{df:Rank}
%<*df:Rank>
\definend{Rank\/}\index{rank} suatu matriks adalah rank baris atau rank kolomnya.
%</df:Rank>
\end{definition}

Jadi, pokok kedua yang telah kita bahas dalam subbagian ini adalah bahwa ruang
kolom dan ruang baris suatu matriks memiliki dimensi yang sama.

Pokok terakhir kita adalah bahwa konsep-konsep yang muncul secara alami dalam
kajian ruang vektor tepat sama dengan konsep-konsep yang telah kita pelajari
dalam sistem linear.

\begin{theorem}  \label{th:RankVsSoltnSp}
%<*tm:RankVsSoltnSp>
Untuk sistem linear dengan \( n \) variabel tak diketahui dan matriks koefisien
\( A \), pernyataan-pernyataan
\begin{tfae}
   \item rank \( A \) adalah \( r \)
   \item ruang vektor solusi sistem homogen terkait berdimensi \( n-r \)
\end{tfae}
saling ekuivalen.
%</tm:RankVsSoltnSp>
\end{theorem}

\par\noindent Jadi, jika sistem memiliki sekurang-kurangnya satu solusi
partikular, jumlah parameter bagi himpunan solusinya sama dengan $n-r$, yakni
jumlah variabel dikurangi rank matriks koefisien.

\begin{proof}
%<*pf:RankVsSoltnSp>
Rank \( A \) adalah \( r \) jika dan hanya jika reduksi Gauss pada \( A \)
berakhir dengan \( r \) baris tak nol.
Hal itu berlaku jika dan hanya jika matriks-matriks bentuk eselon yang ekuivalen
baris dengan \( A \) memiliki \( r \) variabel utama.
Selanjutnya, hal itu berlaku jika dan hanya jika terdapat \( n-r \) variabel
bebas.
%</pf:RankVsSoltnSp>
\end{proof}

\begin{corollary} \label{co:EquivToNonsingular}
%<*co:EquivToNonsingular>
Jika matriks $A$ berukuran \( \nbyn{n} \), pernyataan-pernyataan berikut
\begin{tfae}
  \item rank \( A \) adalah \( n \)
  \item \( A \) nonsingular
  \item baris-baris \( A \) membentuk himpunan bebas linear
  \item kolom-kolom \( A \) membentuk himpunan bebas linear
  \item setiap sistem linear yang matriks koefisiennya \( A \) memiliki tepat
    satu solusi
\end{tfae}
saling ekuivalen.
%</co:EquivToNonsingular>
\end{corollary}

\begin{proof}
%<*pf:EquivToNonsingular>
Jelas bahwa \( \text{(1)}\iff\text{(2)}\iff\text{(3)}\iff\text{(4)} \).
Ekuivalensi terakhir, \( \text{(4)}\iff\text{(5)} \), berlaku karena suatu
himpunan yang terdiri atas \( n \) vektor kolom bebas linear jika dan hanya
jika himpunan itu merupakan basis bagi \( \Re^n \), sedangkan sistem
\begin{equation*}
  c_1\colvec{a_{1,1} \\ a_{2,1} \\ \vdots \\ a_{n,1}}
   +\dots+
  c_n\colvec{a_{1,n} \\ a_{2,n} \\ \vdots \\ a_{n,n}}
   =\colvec{d_1 \\ d_2 \\ \vdots \\ d_n}
\end{equation*}
memiliki solusi tunggal untuk setiap pilihan \( d_1,\dots,d_n\in\Re \) jika dan
hanya jika vektor-vektor yang tersusun dari semua \( a \) di ruas kiri
membentuk basis.
%</pf:EquivToNonsingular>
\end{proof}

\begin{remark}\cite{Munkres}  \label{rem:MunkresCommment}
Kadang-kadang hasil subbagian ini keliru diingat sebagai pernyataan bahwa
solusi umum sistem dengan \( m \)~persamaan dan \( n \)~variabel tak diketahui
menggunakan \( n-m \) parameter.
Jumlah persamaan bukanlah bilangan yang relevan; yang penting ialah jumlah
persamaan bebas, yakni jumlah persamaan dalam suatu himpunan bebas maksimal.
Jika terdapat \( r \) persamaan bebas, solusi umumnya melibatkan \( n-r \)
parameter.
\end{remark}


\begin{exercises}
  \item 
    Transposkan masing-masing matriks atau vektor berikut.
    \begin{exparts*}
      \partsitem \( \begin{mat}[r]
                 2  &1  \\
                 3  &1
               \end{mat}  \)
      \partsitem \( \begin{mat}[r]
                 2  &1  \\
                 1  &3
               \end{mat}  \)
      \partsitem \( \begin{mat}[r]
                 1  &4  &3 \\
                 6  &7  &8
               \end{mat}  \)
      \partsitem \( \colvec[r]{0 \\ 0 \\ 0} \)
      \partsitem \( \rowvec{-1 &-2} \)
    \end{exparts*}
    \begin{answer}  
      \begin{exparts*}
        \partsitem \( \begin{mat}[r]
                   2  &3  \\
                   1  &1
                 \end{mat}  \)
        \partsitem \( \begin{mat}[r]
                   2  &1  \\
                   1  &3
                 \end{mat}  \)
        \partsitem \( \begin{mat}[r]
                   1  &6  \\
                   4  &7  \\
                   3  &8
                 \end{mat}  \)
        \partsitem \( \rowvec{0 &0 &0} \)
        \partsitem \( \colvec[r]{-1 \\ -2} \)
      \end{exparts*}  
    \end{answer}
  \recommended \item 
    Tentukan apakah vektor yang diberikan berada dalam ruang baris matriks.
    \begin{exparts*}
       \partsitem \( \begin{mat}[r]
                  2  &1  \\
                  3  &1
                \end{mat}  \),
             \( \rowvec{1 &0} \)
       \partsitem \( \begin{mat}[r]
                  0  &1  &3  \\
                 -1  &0  &1  \\
                 -1  &2  &7
                \end{mat}  \),
             \( \rowvec{1 &1 &1} \)
    \end{exparts*}
    \begin{answer}  
       \begin{exparts}
         \partsitem Ya.
           Untuk menentukan apakah terdapat $c_1$ dan $c_2$ sedemikian sehingga
           \( c_1\cdot\rowvec{2 &1}+c_2\cdot\rowvec{3 &1}=\rowvec{1 &0} \),
           kita selesaikan
           \begin{equation*}
              \begin{linsys}{2}
                 2c_1  &+  &3c_2  &=  &1  \\
                  c_1  &+  &c_2   &=  &0  
              \end{linsys}
           \end{equation*}
           dan memperoleh \( c_1=-1 \) serta \( c_2=1 \).
           Jadi, vektor tersebut berada dalam ruang baris.
         \partsitem Tidak.
           Persamaan
           \( c_1\rowvec{0 &1 &3}
              +c_2\rowvec{-1 &0 &1}
              +c_3\rowvec{-1 &2 &7}
              =\rowvec{1 &1 &1} \)
           tidak memiliki solusi.
           \begin{equation*}
             \begin{amat}[r]{3}
               0  &-1  &-1  &1  \\
               1  &0   &2   &1  \\
               3  &1   &7   &1
             \end{amat}
             \grstep{\rho_1\leftrightarrow\rho_2}
             \grstep{-3\rho_1+\rho_3}
             \repeatedgrstep{\rho_2+\rho_3}
             \begin{amat}[r]{3}
               1  &0   &2   &1  \\
               0  &-1  &-1  &1  \\
               0  &0   &0   &-1
             \end{amat}
           \end{equation*}
           Jadi, vektor tersebut tidak berada dalam ruang baris.
      \end{exparts}  
    \end{answer}
  \recommended \item 
    Tentukan apakah vektor yang diberikan berada dalam ruang kolom matriks.
    \begin{exparts*}
      \partsitem \( \begin{mat}[r]
                 1  &1  \\
                 1  &1
               \end{mat}  \),
            \( \colvec[r]{1 \\ 3} \)
      \partsitem \( \begin{mat}[r]
                 1  &3  &1 \\
                 2  &0  &4 \\
                 1  &-3 &-3
               \end{mat}  \),
            \( \colvec[r]{1 \\ 0 \\ 0} \)
    \end{exparts*}
    \begin{answer}
       \begin{exparts}
          \partsitem Tidak.
            Untuk menentukan apakah terdapat \( c_1,c_2\in\Re \) sedemikian
            sehingga
            \begin{equation*}
              c_1\colvec[r]{1 \\ 1}
              +c_2\colvec[r]{1 \\ 1}
              =\colvec[r]{1 \\ 3}
            \end{equation*}
            kita dapat menggunakan Metode Gauss pada sistem linear yang
            dihasilkan.
            \begin{equation*}
              \begin{linsys}{2}
                 c_1  &+  &c_2  &=  &1  \\
                 c_1  &+  &c_2  &=  &3  
              \end{linsys}
              \grstep{-\rho_1+\rho_2}
              \begin{linsys}{2}
                 c_1  &+  &c_2  &=  &1  \\
                      &   &0    &=  &2  
              \end{linsys}
            \end{equation*}
            Tidak ada solusi, sehingga vektor tersebut tidak berada dalam ruang
            kolom.
          \partsitem Ya.
            Dari hubungan berikut
            \begin{equation*}
              c_1\colvec[r]{1 \\ 2 \\ 1}
              +c_2\colvec[r]{3 \\ 0 \\ -3}
              +c_3\colvec[r]{1 \\ 4 \\ -3}
              =\colvec[r]{1 \\ 0 \\ 0}
            \end{equation*}
            kita memperoleh sistem linear yang, ketika Metode Gauss diterapkan,
            \begin{equation*}
              \begin{amat}[r]{3}
                1  &3  &1  &1  \\
                2  &0  &4  &0  \\
                1  &-3 &-3 &0
              \end{amat}
              \grstep[-\rho_1+\rho_3]{-2\rho_1+\rho_2}
              \repeatedgrstep{-\rho_2+\rho_3}
              \begin{amat}[r]{3}
                1  &3  &1  &1  \\
                0  &-6 &2  &-2 \\
                0  &0  &-6 &1
              \end{amat}
            \end{equation*}
            menghasilkan suatu solusi.
            Jadi, vektor tersebut berada dalam ruang kolom.
% sage: load('gauss_method.sage')
% sage: M = matrix(QQ, [[1,3, 1, 1], [2, 0, 4, 0], [1, -3, -3, 0]])
% sage: M
% [ 1  3  1  1]
% [ 2  0  4  0]
% [ 1 -3 -3  0]
% sage: gauss_method(M)
% [ 1  3  1  1]
% [ 2  0  4  0]
% [ 1 -3 -3  0]
% (' take', -2, 'times row', 1, 'plus row', 2)
% (' take', -1, 'times row', 1, 'plus row', 3)
% [ 1  3  1  1]
% [ 0 -6  2 -2]
% [ 0 -6 -4 -1]
% (' take', -1, 'times row', 2, 'plus row', 3)
% [ 1  3  1  1]
% [ 0 -6  2 -2]
% [ 0  0 -6  1]
      \end{exparts}  
    \end{answer}
  \recommended \item  
    Tentukan apakah vektor yang diberikan berada dalam ruang kolom matriks.
    \begin{exparts*}
      \partsitem \( \begin{mat}[r]
                   2  &1  \\
                   2  &5
                 \end{mat} \),~\( \colvec[r]{1 \\ -3}  \)
      \partsitem \( \begin{mat}[r]
                   4  &-8 \\
                   2  &-4
                 \end{mat} \),~\( \colvec[r]{0 \\ 1}  \)
      \partsitem \( \begin{mat}[r]
                   1  &-1  &1  \\
                   1  &1   &-1 \\
                  -1  &-1  &1
                \end{mat} \),~\( \colvec[r]{2 \\ 0 \\ 0}  \)
    \end{exparts*}
    \begin{answer}
      \begin{exparts}
        \partsitem Ya;
          kita menanyakan apakah terdapat skalar \( c_1 \) dan \( c_2 \)
          sedemikian sehingga
          \begin{equation*}
            c_1\colvec[r]{2 \\ 2}+c_2\colvec[r]{1 \\ 5}=\colvec[r]{1 \\ -3}
          \end{equation*}
          yang menghasilkan sistem linear
          \begin{equation*}
            \begin{linsys}{2}
              2c_1  &+  &c_2  &=  &1  \\
              2c_1  &+  &5c_2 &=  &-3
            \end{linsys}
            \grstep{-\rho_1+\rho_2}
            \begin{linsys}{2}
              2c_1  &+  &c_2  &=  &1  \\
                    &   &4c_2 &=  &-4
            \end{linsys}
          \end{equation*}
          dan Metode Gauss menghasilkan \( c_2=-1 \) dan \( c_1=1 \).
          Jadi, pasangan skalar tersebut memang ada dan vektor itu memang
          berada dalam ruang kolom matriks.
        \partsitem Tidak;
          kita menanyakan apakah terdapat skalar $c_1$ dan $c_2$ sedemikian
          sehingga
          \begin{equation*}
            c_1\colvec[r]{4 \\ 2}+c_2\colvec[r]{-8 \\ -4}=\colvec[r]{0 \\ 1}
          \end{equation*}
          dan salah satu cara untuk melanjutkan ialah memperhatikan sistem
          linear yang dihasilkan
          \begin{equation*}
            \begin{linsys}{2}
              4c_1  &-  &8c_2  &=  &0  \\
              2c_1  &-  &4c_2  &=  &1
            \end{linsys}
          \end{equation*}
          yang dengan mudah terlihat tidak memiliki solusi.
          Cara lain ialah memperhatikan bahwa setiap kombinasi linear kolom di
          ruas kiri memiliki komponen kedua yang nilainya setengah komponen
          pertama, sedangkan vektor di ruas kanan tidak memenuhi kriteria itu.
        \partsitem Ya; kita cukup mengamati bahwa vektor tersebut sama dengan
          kolom pertama dikurangi kolom kedua.
          Jika pengamatan itu terlewat, susun hubungan di antara kolom-kolom
          \begin{equation*}
            c_1\colvec[r]{1 \\ 1 \\ -1}
             +c_2\colvec[r]{-1 \\ 1 \\ -1}
             +c_3\colvec[r]{1 \\ -1 \\ 1}
             =\colvec[r]{2 \\ 0 \\ 0}
          \end{equation*}
          lalu perhatikan sistem linear yang dihasilkan
          \begin{equation*}
            \begin{linsys}{3}
              c_1  &-  &c_2  &+  &c_3  &=  &2  \\
              c_1  &+  &c_2  &-  &c_3  &=  &0  \\
             -c_1  &-  &c_2  &+  &c_3  &=  &0    
            \end{linsys}
            \grstep[\rho_1+\rho_3]{-\rho_1+\rho_2}
            \begin{linsys}{3}
              c_1  &-  &c_2  &+  &c_3  &=  &2  \\
                   &   &2c_2 &-  &2c_3 &=  &-2  \\
                   &   &-2c_2&+  &2c_3 &=  &2    
            \end{linsys}
            \grstep{\rho_2+\rho_3}
            \begin{linsys}{3}
              c_1  &-  &c_2  &+  &c_3  &=  &2  \\
                   &   &2c_2 &-  &2c_3 &=  &-2  \\
                   &   &     &   &0    &=  &0    
            \end{linsys}
          \end{equation*}
          yang memberikan informasi tambahan (selain adanya sekurang-kurangnya
          satu solusi), yaitu bahwa terdapat tak hingga banyak solusi.
          Parametrisasi memberikan $c_2=-1+c_3$ dan $c_1=1$, sehingga dengan
          mengambil $c_3=0$ diperoleh solusi partikular $c_1=1$, $c_2=-1$, dan
          $c_3=0$ (yang tentu saja sama dengan pengamatan di awal).
      \end{exparts}  
    \end{answer}
  \recommended \item  
    Temukan basis bagi ruang baris matriks berikut.
    \begin{equation*}
      \begin{mat}[r]
         2  &0  &3  &4  \\
         0  &1  &1  &-1 \\
         3  &1  &0  &2  \\
         1  &0  &-4 &-1
       \end{mat}
    \end{equation*}
    \begin{answer}
     % Using Sage:
     % M = matrix(QQ, [[2,0,3,4], [0,1,1,-1], [3,1,0,2], [1,0,-4,-1]])
     % M1 = M.with_added_multiple_of_row(2,0,-3/2)
     % M2 = M1.with_added_multiple_of_row(3,0,-1/2)
     % M3 = M2.with_added_multiple_of_row(2,1,-1)
     % M4 = M3.with_added_multiple_of_row(3,2,-1)
     Reduksi Gauss yang rutin
     \begin{equation*}
       \begin{mat}[r]
         2  &0  &3 &4   \\
         0  &1  &1  &-1 \\
         3  &1  &0  &2  \\
         1  &0  &-4  &-1
       \end{mat}
       \grstep[-(1/2)\rho_1+\rho_4]{-(3/2)\rho_1+\rho_3}
       \grstep{-\rho_2+\rho_3}
       \repeatedgrstep{-\rho_3+\rho_4}
       \begin{mat}[r]
         2  &0  &3      &4   \\
         0  &1  &1      &-1 \\
         0  &0  &-11/2  &-3  \\
         0  &0  &0      &0
       \end{mat}
     \end{equation*}
     menyarankan basis berikut
     $\sequence{\rowvec{2 &0 &3 &4},
         \rowvec{0 &1 &1 &-1},\rowvec{0 &0 &-11/2 &-3}}$.

      Prosedur lain yang mungkin lebih mudah ialah mempertukarkan baris lebih
      dahulu,
      % From Sage
      % M = matrix(QQ, [[2,0,3,4], [0,1,1,-1], [3,1,0,2], [1,0,-4,-1]])
      % M1 = M.with_swapped_rows(0,3)
      % M2 = M1.with_added_multiple_of_row(2,0,-3)
      % M3 = M2.with_added_multiple_of_row(3,0,-2)
      % M4 = M3.with_added_multiple_of_row(2,1,-1)
      % M5 = M4.with_added_multiple_of_row(3,2,-1)
      \begin{equation*}
        \grstep{\rho_1\swap\rho_4}
        \repeatedgrstep[-2\rho_1+\rho_4]{-3\rho_1+\rho_3}
        \repeatedgrstep{-\rho_2+\rho_3}
        \repeatedgrstep{-\rho_3+\rho_4}
        \begin{mat}[r]
           1  &0  &-4 &-1 \\
           0  &1  &1  &-1 \\
           0  &0  &11 &6  \\
           0  &0  &0  &0
         \end{mat}
      \end{equation*}
      yang menghasilkan basis
      \( \sequence{\rowvec{1 &0 &-4 &-1},\rowvec{0 &1 &1 &-1},
                   \rowvec{0 &0 &11 &6}} \).  
    \end{answer}
  \recommended \item 
    Temukan rank masing-masing matriks.
    \begin{exparts*}
      \partsitem \(
        \begin{mat}[r]
          2  &1  &3  \\
          1  &-1 &2  \\
          1  &0  &3
        \end{mat}  \)
      \partsitem \(
        \begin{mat}[r]
          1  &-1 &2  \\
          3  &-3 &6  \\
         -2  &2  &-4
        \end{mat}  \)
      \partsitem \(
        \begin{mat}[r]
          1  &3  &2  \\
          5  &1  &1  \\
          6  &4  &3
        \end{mat}  \)
      \partsitem \(
        \begin{mat}[r]
          0  &0  &0  \\
          0  &0  &0  \\
          0  &0  &0
        \end{mat}  \)
    \end{exparts*}
    \begin{answer}
      \begin{exparts}
         \partsitem Reduksi berikut
           \begin{equation*}
             \mbox{}
             \grstep[-(1/2)\rho_1+\rho_3]{-(1/2)\rho_1+\rho_2}
             \repeatedgrstep{-(1/3)\rho_2+\rho_3}
             \begin{mat}[r]
               2  &1     &3     \\
               0  &-3/2  &1/2   \\
               0  &0     &4/3
             \end{mat}
           \end{equation*}
           menunjukkan bahwa rank baris, dan dengan demikian rank matriks,
           adalah tiga.
         \partsitem Pemeriksaan kolom-kolom menunjukkan bahwa semua kolom
           lainnya merupakan kelipatan kolom pertama (pemeriksaan baris-baris
           menunjukkan hal yang sama).
           Jadi, ranknya adalah satu.
            
           Sebagai cara lain, reduksi
           \begin{equation*}
             \begin{mat}[r]
               1  &-1  &2  \\
               3  &-3  &6  \\   
               -2 &2   &-4
             \end{mat}
             \grstep[2\rho_1+\rho_3]{-3\rho_1+\rho_2}
             \begin{mat}[r]
               1  &-1  &2  \\
               0  &0   &0  \\   
               0  &0   &0 
             \end{mat}
           \end{equation*}
           menunjukkan hal yang sama.
         \partsitem Perhitungan berikut
           \begin{equation*}
             \begin{mat}[r]
               1  &3  &2  \\
               5  &1  &1  \\
               6  &4  &3  
             \end{mat}
             \grstep[-6\rho_1+\rho_3]{-5\rho_1+\rho_2}
             \repeatedgrstep{-\rho_2+\rho_3}
             \begin{mat}[r]
               1  &3   &2  \\
               0  &-14 &-9 \\
               0  &0   &0  
             \end{mat}
           \end{equation*}
           menunjukkan bahwa ranknya adalah dua.
         \partsitem Ranknya adalah nol.
       \end{exparts}  
     \end{answer}
  \item Berikan basis bagi ruang kolom matriks berikut.
    Berikan pula rank matriksnya.
    \begin{equation*}
      \begin{mat}
        1 &3 &-1 &2 \\
        2 &1 &1  &0 \\
        0 &1 &1  &4
      \end{mat}
    \end{equation*}
  \begin{answer}
    Kita mencari basis bagi rentang berikut.
    \begin{equation*}
      \spanof{\colvec{1 \\ 2 \\ 0},
              \colvec{3 \\ 1 \\ 1},
              \colvec{-1 \\ 1 \\ 1},
              \colvec{2 \\ 0 \\ 4}}\subseteq\Re^3
    \end{equation*}
    Pendekatan yang paling langsung ialah mentransposkan kolom-kolom tersebut
    menjadi baris, menggunakan Metode Gauss untuk menemukan basis bagi rentang
    baris-baris itu, lalu mentransposkannya kembali menjadi kolom.
    \begin{equation*}
      \begin{mat}
        1 &2 &0 \\
        3 &1 &1 \\
       -1 &1 &1 \\
        2 &0 &4
      \end{mat}
      \grstep[\rho_1+\rho_3 \\ -2\rho_1+\rho_4]{-3\rho_1+\rho_2}
      \grstep[-(4/5)\rho_2+\rho_4]{(3/5)\rho_2+\rho_3}
      \grstep{-2\rho_3+\rho_4}
      \begin{mat}
        1 &2  &0 \\
        0 &-5 &1 \\
        0 &0  &8/5 \\
        0 &0  &0
      \end{mat}
    \end{equation*}
    Buang vektor nol, yang menunjukkan adanya redundansi di antara vektor-vektor
    awal, untuk memperoleh basis ruang kolom berikut.
    \begin{equation*}
      \sequence{
        \colvec{1 \\ 2 \\ 0},
        \colvec{0 \\ -5 \\ 1},
        \colvec{0 \\ 0 \\ 8/5}
        }
    \end{equation*}
    Rank matriks sama dengan dimensi ruang kolomnya, jadi nilainya tiga.
    (Nilai itu juga sama dengan dimensi ruang barisnya.)
  \end{answer}
  \recommended \item 
    Temukan basis bagi rentang masing-masing himpunan.
    \begin{exparts}
      \partsitem \(
        \set{\rowvec{1 &3},
          \rowvec{-1 &3},
          \rowvec{1 &4},
          \rowvec{2 &1}  }\subseteq\matspace_{\nbym{1}{2}} \)
      \partsitem \(
         \set{\colvec[r]{1 \\2 \\1},
           \colvec[r]{3 \\ 1 \\ -1},
           \colvec[r]{1 \\ -3 \\ -3}  }\subseteq\Re^3  \)
      \partsitem \(  \set{1+x,1-x^2,3+2x-x^2}\subseteq\polyspace_3  \)
      \partsitem \( \set{
          \begin{mat}[r]
            1  &0  &1  \\
            3  &1  &-1
          \end{mat},
          \begin{mat}[r]
            1  &0  &3  \\
            2  &1  &4
          \end{mat},
          \begin{mat}[r]
           -1  &0  &-5 \\
           -1  &-1 &-9
          \end{mat}  }  \subseteq\matspace_{\nbym{2}{3}}  \)
    \end{exparts}
    \begin{answer}
      \begin{exparts}
        \partsitem Reduksi berikut
          \begin{equation*}
             \begin{mat}[r]
               1  &3  \\
              -1  &3  \\
               1  &4  \\
               2  &1
             \end{mat}
             \grstep[-\rho_1+\rho_3 \\ -2\rho_1+\rho_4]{\rho_1+\rho_2}
             \grstep[(5/6)\rho_2+\rho_4]{-(1/6)\rho_2+\rho_3}
             \begin{mat}[r]
               1  &3  \\
               0  &6  \\
               0  &0  \\
               0  &0
             \end{mat}
          \end{equation*}
          memberikan \( \sequence{\rowvec{1 &3},\rowvec{0 &6}} \).
        \partsitem Mentransposkan dan mereduksi
          \begin{equation*}
             \begin{mat}[r]
               1  &2  &1  \\
               3  &1  &-1 \\
               1  &-3 &-3
             \end{mat}
             \grstep[-\rho_1+\rho_3]{-3\rho_1+\rho_2}
             \begin{mat}[r]
               1  &2  &1  \\
               0  &-5 &-4 \\
               0  &-5 &-4
             \end{mat}
             \grstep{-\rho_2+\rho_3}
             \begin{mat}[r]
               1  &2  &1  \\
               0  &-5 &-4 \\
               0  &0  &0
             \end{mat}
          \end{equation*}
          lalu mentransposkan kembali memberikan basis berikut.
          \begin{equation*}
            \sequence{\colvec[r]{1 \\ 2 \\ 1},
              \colvec[r]{0 \\ -5 \\ -4}   }
          \end{equation*}
        \partsitem Perhatikan terlebih dahulu bahwa ruang induknya adalah
          $\polyspace_3$, bukan $\polyspace_2$.
          Selanjutnya, anggap polinomial pertama
          $1+1\cdot x+0\cdot x^2+0\cdot x^3$ ``sama'' dengan vektor baris
          $\rowvec{1 &1 &0 &0}$, dan seterusnya, sehingga diperoleh
          \begin{equation*}
            \begin{mat}[r]
              1  &1  &0  &0 \\
              1  &0  &-1 &0 \\
              3  &2  &-1 &0
            \end{mat}
            \grstep[-3\rho_1+\rho_3]{-\rho_1+\rho_2}
            \repeatedgrstep{-\rho_2+\rho_3}
            \begin{mat}[r]
              1  &1  &0  &0 \\
              0  &-1 &-1 &0 \\
              0  &0  &0  &0
            \end{mat}
          \end{equation*}
          yang menghasilkan basis \( \sequence{1+x,-x-x^2} \).
        \partsitem Di sini, anggapan ``sama'' memberikan
          \begin{equation*}
            \begin{mat}[r]
              1  &0  &1  &3  &1  &-1  \\
              1  &0  &3  &2  &1  &4   \\
             -1  &0  &-5  &-1 &-1 &-9
            \end{mat}
            \grstep[\rho_1+\rho_3]{-\rho_1+\rho_2}
            \repeatedgrstep{2\rho_2+\rho_3}
            \begin{mat}[r]
              1  &0  &1  &3  &1  &-1  \\
              0  &0  &2  &-1 &0  &5   \\
              0  &0  &0  &0  &0  &0
            \end{mat}
          \end{equation*}
          yang menghasilkan basis berikut.
          \begin{equation*}
            \sequence{
            \begin{mat}[r]
              1  &0  &1  \\
              3  &1  &-1
            \end{mat},
            \begin{mat}[r]
              0  &0  &2  \\
             -1  &0  &5
            \end{mat}  }
          \end{equation*}
      \end{exparts}   
     \end{answer}
  \item Berikan basis bagi rentang masing-masing himpunan dalam ruang vektor
    alaminya.
    \begin{exparts}
      \partsitem
         $\set{\colvec{1 \\ 1 \\ 3},
              \colvec{-1 \\ 2 \\ 0},
              \colvec{0  \\ 12 \\ 6}}$
      \partsitem $\set{x+x^2, 2-2x, 7, 4+3x+2x^2}$
    \end{exparts}
    \begin{answer}
      \begin{exparts}
        \partsitem
          Transposkan kolom-kolom menjadi baris, bawa ke bentuk eselon
          (lalu buang setiap baris nol), dan transposkan kembali menjadi kolom.
          \begin{equation*}
            \begin{mat}
              1 &1  &3 \\
             -1 &2  &0 \\
              0 &12 &6
            \end{mat}
            \grstep{\rho_1+\rho_2}
            \grstep{-4\rho_2+\rho_3}
            \begin{mat}
              1 &1  &3 \\
              0 &3  &3 \\
              0 &0  &-6
            \end{mat}
          \end{equation*}
          Salah satu basis bagi rentang tersebut adalah
          \begin{equation*}
            \sequence{
              \colvec{1 \\ 1 \\ 3},
              \colvec{0 \\ 3 \\ 3},
              \colvec{0 \\ 0 \\ -6}
          }
          \end{equation*}
       \partsitem
          Seperti pada bagian sebelumnya, kita menganggapnya sebagai baris agar
          dapat memanfaatkan pekerjaan yang telah kita lakukan dengan Metode
          Gauss.
          \begin{multline*}
            \begin{mat}
              0 &1  &1 \\
              2 &-2 &0 \\
              7 &0 &0 \\
              4 &3  &2
            \end{mat}
            \grstep{\rho_1\leftrightarrow\rho_2}
            \grstep[2\rho_1+\rho_4]{-(7/2)\rho_1+\rho_3}
            \grstep[-7\rho_2+\rho_4]{-7\rho_2+\rho_3}          \\
            \grstep{-(5/7)\rho_3+\rho_4}
            \begin{mat}
              2 &-2  &0 \\
              0 &1 &1 \\
              0 &0  &-7 \\
              0 &0 &0
            \end{mat}
          \end{multline*}
          Salah satu basis bagi rentang himpunan itu adalah
          $\sequence{2-2x, x+x^2, -5x^2}$.
      \end{exparts}
    \end{answer}
  \item 
    Matriks manakah yang berank nol?
    Berank satu?
    \begin{answer}
      Hanya matriks nol yang berank nol.
      Matriks berank satu tepatlah matriks yang berbentuk
      \begin{equation*}
        \begin{mat}
           k_1\cdot \rho  \\
            \vdots  \\
           k_m\cdot \rho
        \end{mat}
      \end{equation*}
      dengan \( \rho \) suatu vektor baris tak nol dan tidak semua
      \( k_i \) bernilai nol.
      (\textit{Catatan.}
      Kita tidak dapat sekadar mengatakan bahwa semua baris merupakan kelipatan
      baris pertama, sebab baris pertama mungkin merupakan baris nol.
      \textit{Catatan lain.}
      Pernyataan di atas juga berlaku jika `baris' diganti dengan `kolom'.)
    \end{answer}
  \recommended \item
    Diberikan \( a,b,c\in\Re \), pilihan \( d \) manakah yang membuat matriks
    berikut berank satu?
    \begin{equation*}
      \begin{mat}
         a  &b  \\
         c  &d
      \end{mat}
    \end{equation*}
    \begin{answer}
      Jika \( a\neq 0 \), pilihan \( d=(c/a)b \) membuat baris kedua menjadi
      kelipatan baris pertama, tepatnya \( c/a \) kali baris pertama.
      Jika \( a=0 \) dan \( b=0 \), sebarang pilihan \( d \) yang bukan
      \( 0 \) memastikan bahwa baris kedua tak nol.
      Jika \( a=0 \), \( b\neq 0 \), dan \( c=0 \), sebarang pilihan \( d \)
      dapat digunakan, sebab matriks otomatis berank satu (bahkan untuk pilihan
      $d=0$).
      Terakhir, jika \( a=0 \), \( b\neq 0 \), dan \( c\neq 0 \), tidak ada
      pilihan \( d \) yang memadai karena matriks pasti berank dua.
    \end{answer}
  \item 
    Temukan rank kolom matriks berikut.
    \begin{equation*}
      \begin{mat}[r]
        1  &3  &-1  &5  &0  &4  \\
        2  &0  &1   &0  &4  &1
      \end{mat}
    \end{equation*}
    \begin{answer}
      Rank kolomnya adalah dua.
      Salah satu cara melihatnya ialah melalui pemeriksaan\Dash ruang kolom
      terdiri atas kolom-kolom setinggi dua, sehingga dimensinya paling tinggi
      dua, dan kita dapat dengan mudah menemukan dua kolom yang bersama-sama
      membentuk himpunan bebas linear (misalnya, kolom keempat dan kelima).
      Cara lain ialah mengingat bahwa rank kolom sama dengan rank baris, lalu
      menjalankan Metode Gauss, yang menyisakan dua baris tak nol.
    \end{answer}
  \item 
    Tunjukkan bahwa suatu sistem linear yang memiliki sekurang-kurangnya satu
    solusi memiliki paling banyak satu solusi jika dan hanya jika rank matriks
    koefisiennya sama dengan jumlah kolom matriks tersebut.
    \begin{answer}
      Kita terapkan \nearbytheorem{th:RankVsSoltnSp}.
      Jumlah kolom matriks koefisien $A$ dari suatu sistem linear sama dengan
      jumlah $n$ variabel tak diketahui.
      Sistem linear yang memiliki sekurang-kurangnya satu solusi memiliki
      paling banyak satu solusi jika dan hanya jika ruang solusi sistem homogen
      terkait berdimensi nol (ingat:~dalam persamaan
      `$\text{Umum}=\text{Partikular}+\text{Homogen}$'
      $\vec{v}=\vec{p}+\vec{h}$, asalkan $\vec{p}$ semacam itu ada, solusi
      $\vec{v}$ tunggal jika dan hanya jika vektor $\vec{h}$ tunggal, yakni
      $\vec{h}=\zero$).
      Namun, menurut teorema tersebut, ini berarti bahwa $n=r$.
    \end{answer}
  \recommended \item
    Jika suatu matriks berukuran \( \nbym{5}{9} \), himpunan manakah yang pasti
    bergantung: himpunan barisnya atau himpunan kolomnya?
    \begin{answer}
      Himpunan kolomnya pasti bergantung karena rank matriks paling tinggi lima,
      sedangkan terdapat sembilan kolom.
    \end{answer}
  \item 
    Berikan contoh yang menunjukkan bahwa, meskipun berdimensi sama, ruang
    baris dan ruang kolom suatu matriks tidak harus sama.
    Apakah keduanya pernah sama?
    \begin{answer}
      Keduanya hampir tidak mungkin sama karena ruang baris merupakan himpunan
      vektor baris, sedangkan ruang kolom merupakan himpunan vektor kolom
      (kecuali jika matriks berukuran \( \nbyn{1} \), yang membuat kedua ruang
      pasti sama).

      \textit{Catatan.}
      Perhatikan
      \begin{equation*}
        A=\begin{mat}[r]
            1  &3  \\
            2  &6
          \end{mat}
      \end{equation*}
      dan perhatikan bahwa ruang baris merupakan himpunan semua kelipatan
      \( \rowvec{1 &3} \), sedangkan ruang kolom terdiri atas kelipatan dari
      \begin{equation*}
        \colvec[r]{1 \\ 2}
      \end{equation*}
      sehingga kita juga tidak dapat berargumen bahwa kedua ruang itu cukup
      merupakan transpos satu sama lain.
    \end{answer}
  \item 
    Tunjukkan bahwa himpunan
    \( \set{(1,-1,2,-3),(1,1,2,0),(3,-1,6,-6)} \) tidak memiliki rentang yang
    sama dengan \( \set{(1,0,1,0),(0,2,0,3)} \).
    Omong-omong, apakah ruang vektornya?
    \begin{answer}
      Pertama, ruang vektornya adalah himpunan kuadruplet bilangan real dengan
      operasi alami.
      Meskipun ini bukan himpunan vektor baris selebar empat, perbedaannya
      kecil\Dash himpunan itu ``sama'' dengan himpunan vektor baris tersebut.
      Jadi, kita akan memperlakukan kuadruplet itu seperti vektor selebar empat.

      Dengan pemahaman itu, salah satu cara melihat bahwa \( (1,0,1,0) \) tidak
      berada dalam rentang himpunan pertama ialah memperhatikan bahwa reduksi
      berikut
      \begin{equation*}
        \begin{mat}[r]
          1  &-1  &2  &-3  \\
          1  &1   &2  &0   \\
          3  &-1  &6  &-6
        \end{mat}
        \grstep[-3\rho_1+\rho_3]{-\rho_1+\rho_2}
        \repeatedgrstep{-\rho_2+\rho_3}
        \begin{mat}[r]
          1  &-1  &2  &-3  \\
          0  &2   &0  &3   \\
          0  &0   &0  &0
        \end{mat}
      \end{equation*}
      dan reduksi berikut
      \begin{equation*}
        \begin{mat}[r]
          1  &-1  &2  &-3  \\
          1  &1   &2  &0   \\
          3  &-1  &6  &-6  \\
          1  &0   &1  &0
        \end{mat}
        \grstep[-3\rho_1+\rho_3 \\ -\rho_1+\rho_4]{-\rho_1+\rho_2}
        \repeatedgrstep[-(1/2)\rho_2+\rho_4]{-\rho_2+\rho_3}
        \repeatedgrstep{\rho_3\leftrightarrow\rho_4}
        \begin{mat}[r]
          1  &-1  &2  &-3  \\
          0  &2   &0  &3   \\
          0  &0   &-1 &3/2 \\
          0  &0   &0  &0   
        \end{mat}
      \end{equation*}
      menghasilkan matriks dengan rank yang berbeda.
      Ini berarti bahwa penambahan $(1,0,1,0)$ ke himpunan tiga kuadruplet
      pertama meningkatkan rank, dan dengan demikian memperbesar rentang,
      himpunan tersebut.
      Jadi, $(1,0,1,0)$ belum berada dalam rentangnya.
    \end{answer}
  \recommended \item 
    Tunjukkan bahwa himpunan vektor kolom berikut
    \begin{equation*}
      \set{\colvec{d_1 \\ d_2 \\ d_3}
           \suchthat
           \text{terdapat $x$, $y$, dan $z$ sedemikian sehingga:\ }
           \begin{linsys}{3}
             3x  &+  &2y  &+  &4z  &=   &d_1   \\
              x  &   &    &-  &z   &=   &d_2   \\
             2x  &+  &2y  &+  &5z  &=   &d_3   
           \end{linsys}
           }
    \end{equation*}
    merupakan subruang dari \( \Re^3 \).
    Temukan sebuah basis.
    \begin{answer}
      Himpunan itu merupakan subruang karena merupakan ruang kolom dari matriks
      \begin{equation*}
        \begin{mat}[r]
          3  &2  &4  \\
          1  &0  &-1 \\
          2  &2  &5
        \end{mat}
      \end{equation*}
      koefisien.
      Untuk menemukan basis bagi ruang kolom,
      \begin{equation*}
        \set{c_1\colvec[r]{3 \\ 1 \\ 2}
             +c_2\colvec[r]{2 \\ 0 \\ 2}
             +c_3\colvec[r]{4 \\ -1 \\ 5}
             \suchthat c_1,c_2,c_3\in\Re}
      \end{equation*}
      kita hilangkan hubungan linear di antara ketiga vektor kolom dari
      himpunan perentang dengan mentransposkan, mereduksi,
      \begin{equation*}
         \begin{mat}[r]
           3  &1  &2  \\
           2  &0  &2  \\
           4  &-1 &5
         \end{mat}
         \grstep[-(4/3)\rho_1+\rho_3]{-(2/3)\rho_1+\rho_2}
         \repeatedgrstep{-(7/2)\rho_2+\rho_3}
         \begin{mat}[r]
           3  &1     &2  \\
           0  &-2/3  &2/3  \\
           0  &0     &0
         \end{mat}
      \end{equation*}
      membuang baris nol, lalu mentransposkan kembali.
      \begin{equation*}
        \sequence{ \colvec[r]{3 \\ 1 \\ 2},
                   \colvec[r]{0 \\ -2/3 \\ 2/3}  }
      \end{equation*}  
    \end{answer}
  \item 
    Tunjukkan bahwa operasi transpos bersifat
    linear:\index{transpos!bersifat linear}
    \begin{equation*}
      \trans{(rA+sB)}  = r\trans{A}+s\trans{B}
    \end{equation*}
    untuk \( r,s\in\Re \) dan \( A,B\in\matspace_{\nbym{m}{n}} \).
    \begin{answer}
      Kita dapat menunjukkannya melalui perhitungan langsung.
      \begin{align*}
        \trans{(rA+sB)}
        &=\trans{\begin{mat}
             ra_{1,1}+sb_{1,1}  &\ldots &ra_{1,n}+sb_{1,n} \\
                          &\vdots            \\
             ra_{m,1}+sb_{m,1}  &\ldots &ra_{m,n}+sb_{m,n}
                \end{mat}  }  \\
        &=\begin{mat}
           ra_{1,1}+sb_{1,1}  &\ldots &ra_{m,1}+sb_{m,1}  \\
                            &\vdots                           \\
           ra_{1,n}+sb_{1,n}  &\ldots &ra_{m,n}+sb_{m,n}
         \end{mat}           \\
        &=\begin{mat}
           ra_{1,1}  &\ldots &ra_{m,1}  \\
                    &\vdots                       \\
           ra_{1,n}  &\ldots &ra_{m,n}
         \end{mat}           
        +\begin{mat}
           sb_{1,1}  &\ldots &sb_{m,1}  \\
                    &\vdots                       \\
           sb_{1,n}  &\ldots &sb_{m,n}
         \end{mat}           \\
        &=r\trans{A}+s\trans{B}
      \end{align*}  
    \end{answer}
  \recommended \item  
    Dalam subbagian ini, kita telah menunjukkan bahwa reduksi Gauss menemukan
    basis bagi ruang baris.
    \begin{exparts}
      \partsitem Tunjukkan bahwa basis ini tidak tunggal\Dash reduksi yang
        berbeda dapat menghasilkan basis yang berbeda.
      \partsitem Berikan matriks-matriks yang ruang barisnya sama tetapi jumlah
        barisnya berbeda.
      \partsitem Buktikan bahwa dua matriks memiliki ruang baris yang sama jika
        dan hanya jika, setelah reduksi Gauss--Jordan, keduanya memiliki
        baris-baris tak nol yang sama.
    \end{exparts}
    \begin{answer}
      \begin{exparts}
        \partsitem Reduksi-reduksi berikut memberikan basis yang berbeda.
          \begin{equation*}
            \begin{mat}[r]
              1  &2  &0  \\
              1  &2  &1
            \end{mat}
            \grstep{-\rho_1+\rho_2}
            \begin{mat}[r]
              1  &2  &0  \\
              0  &0  &1
            \end{mat}
            \qquad
            \begin{mat}[r]
              1  &2  &0  \\
              1  &2  &1
            \end{mat}
            \grstep{-\rho_1+\rho_2}
            \repeatedgrstep{2\rho_2}
            \begin{mat}[r]
              1  &2  &0  \\
              0  &0  &2
            \end{mat}
          \end{equation*}
        \partsitem Berikut sebuah contoh sederhana.
          \begin{equation*}
            \begin{mat}[r]
              1  &2  &1  \\
              3  &1  &4
            \end{mat}
            \qquad
            \begin{mat}[r]
              1  &2  &1  \\
              3  &1  &4  \\
              0  &0  &0
            \end{mat}
          \end{equation*}
          Berikut contoh yang tidak terlalu sederhana.
          \begin{equation*}
            \begin{mat}[r]
              1  &2  &1  \\
              3  &1  &4
            \end{mat}
            \qquad
            \begin{mat}[r]
              1  &2  &1  \\
              3  &1  &4  \\
              2  &4  &2  \\
              4  &3  &5
            \end{mat}
          \end{equation*}
        \partsitem 
          % We give two proofs.
          % The first is a bit lower-level but more constructive.
          % The second is higher level.

          % \textit{First proof.}
          % Assume that \( A \) and \( B \) are matrices with equal row spaces.
          % Construct a matrix \( C \) with the rows of \( A \) above the rows
          % of \( B \), and another matrix \( D \) with the rows of \( B \)
          % above the rows of \( A \).
          % \begin{equation*}
          %   C=\begin{mat}
          %        A \\ B
          %     \end{mat}
          %   \qquad
          %   D=\begin{mat}
          %        B \\ A
          %     \end{mat}
          % \end{equation*}
          % Observe that \( C \) and \( D \) are row-equivalent (via a sequence
          % of row-swaps) and so Gauss-Jordan reduce to the same reduced
          % echelon form matrix.

          % Because the row spaces are equal, the rows of \( B \)
          % are linear combinations of the rows of \( A \) so Gauss-Jordan
          % reduction on \( C \) simply turns the rows of \( B \) to zero rows
          % and thus the nonzero rows of $C$ are just the nonzero rows obtained
          % by Gauss-Jordan reducing \( A \).
          % The same can be said for the matrix \( D \)\Dash Gauss-Jordan
          % reduction on \( D \) gives the same non-zero rows as are produced
          % by reduction on \( B \) alone.
          % Therefore, 
          % \( A \) yields the same nonzero rows as \( C  \), which yields
          % the same nonzero rows as \( D \), which yields the same nonzero
          % rows as \( B \).

          % \textit{Second proof (N~Schenker).}
          Karena ruang baris \( A \) dan \( B \) sama, kedua matriks itu
          ekuivalen baris.
          Karena setiap kelas ekuivalensi baris memuat tepat satu bentuk eselon
          tereduksi (Teorema Satu.III.\ref{th:ReducedEchelonFormIsUnique}),
          bentuk eselon tereduksi \( A \) pasti sama dengan bentuk eselon
          tereduksi \( B \).
      \end{exparts}  
    \end{answer}
  \item 
    Mengapa \nearbyremark{rem:MunkresCommment} tidak menimbulkan masalah pada
    kasus ketika \( r \) lebih besar daripada \( n \)?
    \begin{answer}
      \( r \) tidak mungkin lebih besar.
    \end{answer}
  \item  
    Tunjukkan bahwa rank baris suatu matriks \( \nbym{m}{n} \) paling tinggi
    \( m \).
    Adakah batas yang lebih baik?
    \begin{answer}
      Jumlah baris dalam suatu himpunan bebas linear maksimal tidak dapat
      melebihi jumlah seluruh baris.
      Batas yang lebih baik (dan secara umum merupakan batas terbaik yang
      mungkin) adalah nilai minimum dari \( m \) dan \( n \), karena rank baris
      sama dengan rank kolom.
     \end{answer}
  \item
    Tunjukkan bahwa rank suatu matriks sama dengan rank transposnya.
    \begin{answer}
      Karena baris-baris matriks \( A \) merupakan kolom-kolom \( \trans{A} \),
      dimensi ruang baris \( A \) sama dengan dimensi ruang kolom
      \( \trans{A} \).
      Namun, dimensi ruang baris \( A \) adalah rank \( A \), dan dimensi ruang
      kolom \( \trans{A} \) adalah rank \( \trans{A} \).
      Jadi, kedua rank tersebut sama.
    \end{answer}
  \item 
    Benar atau salah:~ruang kolom suatu matriks sama dengan ruang baris
    transposnya.
    \begin{answer}
      Salah.
      Yang pertama merupakan himpunan kolom, sedangkan yang kedua merupakan
      himpunan baris.

      Namun, contoh berikut
      \begin{equation*}
         A=\begin{mat}[r]
             1  &2  &3  \\
             4  &5  &6
           \end{mat},
         \qquad
         \trans{A}=\begin{mat}[r]
             1  &4  \\
             2  &5  \\
             3  &6
           \end{mat}
      \end{equation*}
      menunjukkan bahwa begitu kita memiliki makna formal bagi ``sama'', kita
      dapat menerapkannya di sini:
      \begin{equation*}
        \colspace{A}=\spanof{\set{
                        \colvec[r]{1 \\ 4},
                        \colvec[r]{2 \\ 5},
                        \colvec[r]{3 \\ 6} }  }
      \end{equation*}
      sedangkan
      \begin{equation*}
         \rowspace{\trans{A}}=\spanof{\set{
                                 \rowvec{1 &4},
                                 \rowvec{2 &5},
                                 \rowvec{3 &6} }  }
      \end{equation*}  
      keduanya ``sama'' satu sama lain.
    \end{answer}
  \recommended \item  
    Kita telah melihat bahwa operasi baris dapat mengubah ruang kolom.
    Haruskah operasi baris selalu mengubahnya?
    \begin{answer}
      Tidak.
      Dalam contoh berikut, Metode Gauss tidak mengubah ruang kolom.
      \begin{equation*}
        \begin{mat}[r]
          1  &0  \\
          3  &1
        \end{mat}
        \grstep{-3\rho_1+\rho_2}
        \begin{mat}[r]
          1  &0  \\
          0  &1
        \end{mat}
      \end{equation*} 
    \end{answer}
  \item
    Buktikan bahwa suatu sistem linear memiliki solusi jika dan hanya jika
    matriks koefisien sistem itu memiliki rank yang sama dengan matriks
    diperbesarnya.
    \begin{answer}
      Suatu sistem linear
      \begin{equation*}
        c_1\vec{a}_1+\dots+c_n\vec{a}_n=\vec{d}
      \end{equation*}
      memiliki solusi jika dan hanya jika \( \vec{d} \) berada dalam rentang
      himpunan \( \set{\vec{a}_1,\dots,\vec{a}_n} \).
      Hal itu berlaku jika dan hanya jika rank kolom matriks diperbesar sama
      dengan rank kolom matriks koefisien.
      Karena rank sama dengan rank kolom, sistem tersebut memiliki solusi jika
      dan hanya jika rank matriks diperbesarnya sama dengan rank matriks
      koefisiennya.
    \end{answer}
  \item 
    Suatu matriks \( \nbym{m}{n} \) memiliki
    \definend{rank baris penuh}\index{rank baris penuh}\index{rank baris!penuh}
    jika rank barisnya \( m \), dan memiliki
    \definend{rank kolom penuh}\index{rank kolom penuh}\index{kolom!rank!penuh}
    jika rank kolomnya \( n \).
    \begin{exparts}
      \partsitem Tunjukkan bahwa suatu matriks hanya dapat memiliki rank baris
        penuh sekaligus rank kolom penuh jika matriks itu persegi.
      \partsitem Buktikan bahwa sistem linear dengan matriks koefisien \( A \)
        memiliki solusi untuk sebarang \( d_1 \), \ldots, \( d_m \) di ruas
        kanan jika dan hanya jika \( A \) memiliki rank baris penuh.
      \partsitem Buktikan bahwa suatu sistem homogen memiliki solusi tunggal
        jika dan hanya jika matriks koefisiennya $A$ memiliki rank kolom penuh.
      \partsitem Buktikan bahwa pernyataan ``jika suatu sistem dengan matriks
        koefisien \( A \) memiliki solusi, maka solusinya tunggal'' berlaku jika
        dan hanya jika \( A \) memiliki rank kolom penuh.
    \end{exparts}
    \begin{answer}
      \begin{exparts}
        \partsitem Rank baris sama dengan rank kolom, sehingga masing-masing
          paling tinggi sama dengan nilai minimum dari jumlah baris dan jumlah
          kolom.
          Jadi, keduanya hanya dapat penuh jika jumlah baris sama dengan jumlah
          kolom.
          (Tentu saja, arah sebaliknya tidak berlaku:~matriks persegi tidak
          harus memiliki rank baris penuh atau rank kolom penuh.)
        \partsitem Jika \( A \) memiliki rank baris penuh, maka apa pun ruas
          kanannya, Metode Gauss pada matriks diperbesar berakhir dengan satu
          entri utama pada setiap baris dan tidak satu pun entri utama itu
          berada dalam kolom paling kanan (kolom ``augmentasi'').
          Substitusi mundur kemudian memberikan suatu solusi.

          Di sisi lain, jika sistem linear tidak memiliki solusi untuk suatu
          ruas kanan, satu-satunya penyebabnya ialah Metode Gauss menyisakan
          suatu baris yang seluruh entrinya di sebelah kiri garis
          ``augmentasi'' bernilai nol, tetapi entrinya di sebelah kanan tak nol.
          Jadi, jika sistem dengan matriks koefisien $A$ tidak memiliki solusi
          untuk suatu ruas kanan, maka \( A \) tidak memiliki rank baris penuh
          karena beberapa barisnya telah tereliminasi.
        \partsitem Matriks \( A \) memiliki rank kolom penuh jika dan hanya jika
          kolom-kolomnya membentuk himpunan bebas linear.
          Ini ekuivalen dengan hanya adanya hubungan linear trivial di antara
          kolom-kolom tersebut, sehingga satu-satunya solusi sistem adalah saat
          setiap variabel bernilai $0$.
        \partsitem Matriks \( A \) memiliki rank kolom penuh jika dan hanya jika
          himpunan kolomnya bebas linear, sehingga membentuk basis bagi
          rentangnya.
          Ini ekuivalen dengan adanya representasi linear yang tunggal bagi
          setiap vektor dalam rentang tersebut.
          Itulah hasil yang diminta, sebab setiap representasi linear suatu
          vektor dalam rentang merupakan solusi sistem linear.
      \end{exparts}  
    \end{answer}
  \item 
    Bagaimana kesimpulan \nearbylemma{le:RowSpUnchByGR} berubah jika Metode
    Gauss diubah agar mengizinkan perkalian suatu baris dengan nol?
    \begin{answer}
      Alih-alih ruang baris keduanya sama, ruang baris $B$ akan menjadi subruang
      (yang mungkin sama dengan) ruang baris $A$.
    \end{answer}
  \item
    Apakah hubungan antara \( \rank(A) \) dan \( \rank(-A) \)?
    Antara \( \rank(A) \) dan \( \rank(kA) \)?
    Hubungan apakah, jika ada, antara \( \rank(A) \), \( \rank(B) \), dan
    \( \rank(A+B) \)?
    \begin{answer}
      Jelas bahwa \( \rank(A)=\rank(-A) \), sebab Metode Gauss mengizinkan kita
      mengalikan semua baris suatu matriks dengan \( -1 \).
      Dengan cara yang sama, ketika \( k\neq 0 \), kita memperoleh
      \( \rank(A)=\rank(kA) \).

      Penjumlahan lebih menarik.
      Rank suatu jumlah dapat lebih kecil daripada rank suku-sukunya.
      \begin{equation*}
        \begin{mat}[r]
          1  &2  \\
          3  &4
        \end{mat}
        +
        \begin{mat}[r]
         -1  &-2  \\
         -3  &-4
        \end{mat}
        =
        \begin{mat}[r]
         0   &0   \\
         0   &0
        \end{mat}
      \end{equation*}
      Rank suatu jumlah dapat lebih besar daripada rank suku-sukunya.
      \begin{equation*}
        \begin{mat}[r]
          1  &2  \\
          0  &0
        \end{mat}
        +
        \begin{mat}[r]
          0  &0   \\
          3  &4
        \end{mat}
        =
        \begin{mat}[r]
         1   &2   \\
         3   &4
        \end{mat}
      \end{equation*}
      Namun, terdapat batas atas (selain ukuran matriksnya).
      Secara umum, \( \rank(A+B)\leq\rank(A)+\rank(B) \).

      Untuk membuktikannya, perhatikan bahwa kita dapat melakukan eliminasi
      Gauss pada \( A+B \) melalui salah satu dari dua cara:
      pertama-tama kita dapat menjumlahkan \( A \) dan \( B \), lalu menerapkan
      barisan langkah reduksi yang sesuai
      \begin{equation*}
        (A+B)
        \grstep{\text{langkah}_1}
        \cdots
        \grstep{\text{langkah}_k}
        \text{bentuk eselon}
      \end{equation*}
      atau kita dapat memperoleh hasil yang sama dengan menjalankan
      \( \text{langkah}_1 \) sampai \( \text{langkah}_k \) secara terpisah pada
      \( A \) dan \( B \), lalu menjumlahkannya.
      Rank terbesar yang dapat dihasilkan dalam kasus kedua jelaslah jumlah
      kedua rank tersebut.
      (Matriks-matriks di atas memberi contoh terjadinya kedua kemungkinan,
      yakni $\rank(A+B)<\rank(A)+\rank(B)$ dan
      $\rank(A+B)=\rank(A)+\rank(B)$.)
    \end{answer}
    % Relationship between rank(A), rank(B) and rank(cA+dB)?
\end{exercises}























\subsectionoptional{Mengombinasikan Subruang}
\index{jumlah langsung|(}
\textit{Subbagian ini bersifat opsional.
Subbagian ini hanya diperlukan untuk bagian terakhir Bab Tiga dan Bab Lima,
serta untuk beberapa latihan.
Anda dapat melewatinya tanpa kehilangan kesinambungan.}

% This chapter opened with the definition of a vector space, and the middle
% consisted of a first analysis of the idea.
% This subsection closes the chapter by finishing the analysis, in the sense
% that `analysis' means ``method of determining the \ldots{} essential
% features of something by separating it into parts'' \cite{MacmillanDict}.

Salah satu cara memahami sesuatu ialah melihat bagaimana sesuatu itu dibangun
dari bagian-bagian penyusunnya.
Sebagai contoh, kita kadang-kadang membayangkan \( \Re^3 \) tersusun dari
sumbu-\( x \), sumbu-\( y \), dan sumbu-\( z \).
Dalam subbagian ini, kita akan menjelaskan cara menguraikan suatu ruang vektor
menjadi kombinasi beberapa subruangnya.
Saat mengembangkan gagasan kombinasi subruang ini, kita akan terus menggunakan
contoh $\Re^3$ sebagai prototipe.

Subruang merupakan subhimpunan, dan himpunan dapat digabungkan melalui operasi
gabungan.
Namun, operasi gabungan himpunan biasa bukanlah operasi kombinasi subruang yang
kita inginkan.
Sebagai contoh, gabungan sumbu-\( x \), sumbu-\( y \), dan sumbu-\( z \) bukan
seluruh $\Re^3$.
Bahkan, gabungan itu bukan subruang karena tidak tertutup terhadap penjumlahan:
vektor berikut
\begin{equation*}
  \colvec[r]{1 \\ 0 \\ 0}
  +
  \colvec[r]{0 \\ 1 \\ 0}
  +
  \colvec[r]{0 \\ 0 \\ 1}
  =
  \colvec[r]{1 \\ 1 \\ 1}
\end{equation*}
tidak berada pada satu pun dari ketiga sumbu tersebut, sehingga tidak berada
dalam gabungannya.
Karena itu, untuk mengombinasikan subruang, selain anggota-anggota subruang
tersebut, kita sekurang-kurangnya juga harus menyertakan semua kombinasi
linearnya.

\begin{definition}
Jika \( W_1,\dots, W_k \) merupakan subruang dari suatu ruang vektor,
\definend{jumlah}\index{jumlah!subruang}\index{subruang!jumlah}
subruang-subruang itu adalah rentang dari gabungannya,
$W_1+W_2+\dots +W_k=\spanof{W_1\union W_2\union \cdots W_k}$.
\end{definition}

\noindent Penulisan `\( + \)' sesuai dengan kebiasaan menggunakan lambang ini
untuk operasi penggabungan yang alami.

\begin{example} \label{ex:RThreeSumAxes}
Prototipe $\Re^3$ kita sesuai dengan definisi ini.
Setiap vektor $\vec{w}\in\Re^3$ merupakan kombinasi linear
$c_1\vec{v}_1+c_2\vec{v}_2+c_3\vec{v}_3$, dengan $\vec{v}_1$ anggota
sumbu-$x$, dan seterusnya, sebagai berikut:
\begin{equation*}
    \colvec{w_1 \\ w_2 \\ w_3}
    =1\cdot\colvec{w_1 \\ 0 \\ 0}
    +1\cdot\colvec{0 \\ w_2 \\ 0}
    +1\cdot\colvec{0 \\ 0 \\ w_3}
\end{equation*}
sehingga $\text{sumbu-$x$}+\text{sumbu-$y$}+\text{sumbu-$z$}=\Re^3$.
\end{example}

\begin{example}
Jumlah subruang dapat lebih kecil daripada seluruh ruang.
Di dalam \( \polyspace_4 \), misalkan \( L \) merupakan subruang polinomial
linear
$\set{a+bx\suchthat a,b\in\Re}$
dan misalkan \( C \) merupakan subruang polinomial kubik murni
\( \set{cx^3\suchthat c\in\Re}\).
Maka \mbox{$L+C$} bukan seluruh $\polyspace_4$.
Sebaliknya,
\( \mbox{$L+C$}=\set{a+bx+cx^3\suchthat a,b,c\in\Re} \).
\end{example}

\begin{example} \label{exam:RThreeIsSumXYAndYZ}
Suatu ruang dapat dinyatakan sebagai kombinasi subruang dengan lebih dari satu
cara.
Selain penguraian
$\Re^3=\text{sumbu-$x$}+\text{sumbu-$y$}+\text{sumbu-$z$}$,
kita juga dapat menulis
$\Re^3=\text{bidang-$xy$}+\text{bidang-$yz$}$.
Untuk memeriksanya, perhatikan bahwa setiap $\vec{w}\in\Re^3$ dapat ditulis
sebagai kombinasi linear dari anggota bidang-$xy$ dan anggota bidang-$yz$;
berikut dua kombinasi semacam itu.
\begin{equation*}
  \colvec{w_1 \\ w_2 \\ w_3}
  =1\cdot\colvec{w_1 \\ w_2 \\ 0}
   +1\cdot\colvec{0 \\ 0 \\ w_3}
  \qquad
  \colvec{w_1 \\ w_2 \\ w_3}
  =1\cdot\colvec{w_1 \\ w_2/2 \\ 0}
   +1\cdot\colvec{0 \\ w_2/2 \\ w_3}
\end{equation*}
\end{example}

Definisi di atas memberikan salah satu cara untuk memandang suatu ruang sebagai
kombinasi beberapa bagiannya.
Namun, contoh sebelumnya menunjukkan bahwa sekurang-kurangnya ada satu sifat
menarik dari model acuan kita yang tidak tercakup dalam definisi jumlah
subruang.
Dalam penguraian $\Re^3$ yang lazim, kita sering berbicara tentang
`bagian~$x$', `bagian~$y$', atau `bagian~$z$' dari suatu vektor.
Dengan kata lain, dalam prototipe kita, setiap vektor memiliki penguraian
tunggal menjadi komponen-komponen dari bagian yang menyusun seluruh ruang.
Namun, dalam penguraian yang digunakan pada
\nearbyexample{exam:RThreeIsSumXYAndYZ}, kita tidak dapat merujuk pada
``bagian~$xy$'' suatu vektor\Dash ketiga penjumlahan berikut
\begin{equation*}
  \colvec[r]{1 \\ 2 \\ 3}
  =\colvec[r]{1 \\ 2 \\ 0}
  +\colvec[r]{0 \\ 0 \\ 3}
  =\colvec[r]{1 \\ 0 \\ 0}
  +\colvec[r]{0 \\ 2 \\ 3} 
  =\colvec[r]{1 \\ 1 \\ 0}
  +\colvec[r]{0 \\ 1 \\ 3} 
\end{equation*}
semuanya menyatakan vektor tersebut sebagai penjumlahan sesuatu dari bidang
pertama dengan sesuatu dari bidang kedua, tetapi ``bagian~$xy$'' pada
masing-masing penjumlahan berbeda.

%What's happening here, what allows the vector's second entry 
%$2$ to be taken from the $xy$-plane or
%from the $yz$-plane or from a combination, is that the two planes overlap.
%Their intersection is the $y$-axis and so the middle component 
%can shift around.
%In terms of our `analysis', the problem with the sum of subspaces operation is
%that, while the expression
%$\Re^3=\text{$xy$-plane}+\text{$yz$-plane}$ describes the space as a
%combination of its parts, those parts are not separate.

%\begin{lemma}
%The intersection of subspaces is a subspace.\index{subspace!intersection}
%\end{lemma}

%\begin{proof}
%To show that it is a subspace, we need only show that it is nonempty and
%closed under linear combinations.
%It is nonempty because it contains the zero vector.
%For closure, consider a linear combination of vectors from the
%intersection. 
%We can view that as a linear combination of vectors from the first subspace,
%because each vector in the intersection is in the first subspace.
%As all of the vectors are in the first subspace, 
%their combination sums to a vector that is also in the first subspace.
%In the same way, the combination sums to a vector that is in each subspace, 
%and so is in the intersection.
%\end{proof}

%Therefore, an intersection of subspaces is never empty; the smallest that it
%can be is the trivial subspace.
%So the above reference to subspaces that ``overlap'' must be taken to mean
%that their intersection is larger than the trivial subspace, 
%not just that their intersection is nonempty.

Jadi, ketika mempertimbangkan bagaimana $\Re^3$ disusun dari ketiga sumbu, yang
kita maksud mungkin ``sedemikian sehingga setiap vektor memiliki
sekurang-kurangnya satu penguraian'', yang menghasilkan definisi di atas.
Namun, jika yang kita maksud adalah ``sedemikian sehingga setiap vektor
memiliki tepat satu penguraian'', kita memerlukan syarat lain pada kombinasi
tersebut.
Untuk melihat syarat itu, ingat bahwa vektor direpresentasikan secara tunggal
terhadap suatu basis.
Kita dapat menggunakan fakta ini untuk menguraikan suatu ruang menjadi jumlah
subruang sedemikian sehingga setiap vektor dalam ruang terurai secara tunggal
menjadi jumlah anggota-anggota subruang tersebut.

\begin{example} \label{exam:BenchDirSum}
Perhatikan $\Re^3$ dengan basis standarnya
$\stdbasis_3=\sequence{\vec{e}_1,\vec{e}_2,\vec{e}_3}$.
Subruang dengan basis $B_1=\sequence{\vec{e}_1}$ adalah sumbu-$x$, subruang
dengan basis $B_2=\sequence{\vec{e}_2}$ adalah sumbu-$y$, dan subruang dengan
basis $B_3=\sequence{\vec{e}_3}$ adalah sumbu-$z$.
Fakta bahwa setiap anggota $\Re^3$ dapat dinyatakan sebagai jumlah vektor dari
subruang-subruang ini
\begin{equation*}
  \colvec{x \\ y \\ z}
   =\colvec{x \\ 0 \\ 0}
    +\colvec{0 \\ y \\ 0}
    +\colvec{0 \\ 0 \\ z}
\end{equation*}
mencerminkan fakta bahwa $\stdbasis_3$ merentang ruang tersebut\Dash persamaan
berikut
\begin{equation*}
  \colvec{x \\ y \\ z}
   =c_1\colvec[r]{1 \\ 0 \\ 0}
    +c_2\colvec[r]{0 \\ 1 \\ 0}
    +c_3\colvec[r]{0 \\ 0 \\ 1}
\end{equation*}
memiliki solusi untuk setiap $x,y,z\in\Re$.
Fakta bahwa setiap pernyataan semacam itu tunggal mencerminkan kebebasan linear
$\stdbasis_3$, sehingga setiap persamaan seperti di atas memiliki solusi
tunggal.
\end{example}

\begin{example}
Kita tidak harus mengambil vektor basis satu per satu; kita dapat
menggabungkannya menjadi barisan-barisan yang lebih besar.
Perhatikan kembali ruang $\Re^3$ dan vektor-vektor dari basis standar
$\stdbasis_3$.
Subruang dengan basis $B_1=\sequence{\vec{e}_1,\vec{e}_3}$ adalah bidang-$xz$.
Subruang dengan basis $B_2=\sequence{\vec{e}_2}$ adalah sumbu-$y$.
Seperti pada contoh sebelumnya, fakta bahwa setiap anggota ruang merupakan
jumlah anggota kedua subruang dalam tepat satu cara
\begin{equation*}
  \colvec{x \\ y \\ z}
   =\colvec{x \\ 0 \\ z}
    +\colvec{0 \\ y \\ 0}
\end{equation*}
mencerminkan fakta bahwa vektor-vektor ini membentuk basis\Dash persamaan
berikut
\begin{equation*}
  \colvec{x \\ y \\ z}
   =(c_1\colvec[r]{1 \\ 0 \\ 0}
    +c_3\colvec[r]{0 \\ 0 \\ 1})
    +c_2\colvec[r]{0 \\ 1 \\ 0}
\end{equation*}
memiliki tepat satu solusi untuk setiap $x,y,z\in\Re$.
\end{example}

% These examples illustrate 
% the natural way to decompose a space into a sum of subspaces so that
% each vector decomposes uniquely into a sum of vectors from the parts.

\begin{definition}
\definend{Konkatenasi}\index{konkatenasi barisan}%
\index{barisan!konkatenasi}
dari barisan-barisan
$
   B_1=\sequence{\vec{\beta}_{1,1},\dots,\vec{\beta}_{1,n_1}}$, 
  \ldots,
  $B_k=\sequence{\vec{\beta}_{k,1},\dots,\vec{\beta}_{k,n_k}}$ 
menggabungkannya menjadi satu barisan.
\begin{equation*}
 \cat{\cat{\cat{B_1}{B_2}}{\cdots}}{B_k}=
   \sequence{\vec{\beta}_{1,1},\dots,\vec{\beta}_{1,n_1},
            \vec{\beta}_{2,1},\dots,\vec{\beta}_{k,n_k} } 
\end{equation*}
\end{definition}

\begin{lemma} \label{le:UniqDecIffBasisDec}
Misalkan $V$ merupakan ruang vektor yang menjadi jumlah beberapa subruangnya,
$V=W_1+\dots+W_k$.
Misalkan $B_1$, \ldots, $B_k$ merupakan basis bagi subruang-subruang tersebut.
Pernyataan-pernyataan berikut saling ekuivalen.
\begin{tfae}
  \item Pernyataan sebarang $\vec{v}\in V$ sebagai kombinasi
     $\vec{v}=\vec{w}_1+\dots+\vec{w}_k$ 
     dengan $\vec{w}_i\in W_i$ bersifat tunggal.
  \item Konkatenasi $\cat{\cat{B_1}{\cdots}}{B_k}$ merupakan basis bagi $V$.
  \item Setiap hubungan linear di antara vektor-vektor tak nol dari
    subruang-subruang $W_i$ yang berbeda bersifat trivial.
% ; that is, 
%     any set of nonzero vectors $\set{\vec{w}_1,\dots,\vec{w}_k}$
%     with $\vec{w}_i\in W_i$ is linearly independent.
\end{tfae}
\end{lemma}

\begin{proof}
Kita akan menunjukkan bahwa $\text{(1)}\implies\text{(2)}$, kemudian
$\text{(2)}\implies\text{(3)}$, dan terakhir
$\text{(3)}\implies\text{(1)}$.
Untuk argumen-argumen ini, perhatikan bahwa kita dapat beralih dari kombinasi
vektor-vektor $\vec{w}$ ke kombinasi vektor-vektor $\vec{\beta}$
\begin{align*}
  d_1\vec{w}_1+\dots+d_k\vec{w}_k                                 
    &= d_1(c_{1,1}\vec{\beta}_{1,1}+\dots+c_{1,n_1}\vec{\beta}_{1,n_1})   \\
    &\qquad\hbox{}+\dots
        +d_k(c_{k,1}\vec{\beta}_{k,1}+\dots+c_{k,n_k}\vec{\beta}_{k,n_k}) \\
    &= d_1c_{1,1}\cdot\vec{\beta}_{1,1}
        +\dots
        +d_kc_{k,n_k}\cdot\vec{\beta}_{k,n_k}   
  \tag{$*$}
\end{align*}
dan sebaliknya (kita dapat bergerak dari baris bawah ke baris atas dengan
mengambil setiap $d_i$ sama dengan $1$).

Untuk $\text{(1)}\implies\text{(2)}$, andaikan semua penguraian bersifat
tunggal.
Kita akan menunjukkan bahwa $\cat{\cat{B_1}{\cdots}}{B_k}$ merentang ruang dan
bebas linear.
Konkatenasi itu merentang ruang karena asumsi $V=W_1+\dots+W_k$ berarti bahwa
setiap $\vec{v}$ dapat dinyatakan sebagai
$\vec{v}=\vec{w}_1+\dots+\vec{w}_k$, yang melalui persamaan~($*$) berubah
menjadi pernyataan $\vec{v}$ sebagai kombinasi linear vektor-vektor
$\vec{\beta}$ dari konkatenasi tersebut.
Untuk kebebasan linear, perhatikan hubungan linear berikut.
\begin{equation*}
  \zero=c_{1,1}\vec{\beta}_{1,1}+\dots+c_{k,n_k}\vec{\beta}_{k,n_k}
\end{equation*}
Kelompokkan kembali seperti dalam ($*$) (yakni, bergerak dari baris bawah ke
baris atas) untuk memperoleh penguraian
$\zero=\vec{w}_1+\dots+\vec{w}_k$.
Karena vektor nol jelas memiliki penguraian
$\zero=\zero+\dots+\zero$, asumsi bahwa penguraian bersifat tunggal menunjukkan
bahwa setiap $\vec{w}_i$ merupakan vektor nol.
Ini berarti bahwa
$c_{i,1}\vec{\beta}_{i,1}+\dots+c_{i,n_i}\vec{\beta}_{i,n_i}=\zero$, dan
karena setiap $B_i$ merupakan basis, kita memperoleh kesimpulan yang diinginkan
bahwa semua $c$ bernilai nol.

Untuk $\text{(2)}\implies\text{(3)}$, andaikan konkatenasi basis-basis tersebut
merupakan basis bagi seluruh ruang.
Perhatikan suatu hubungan linear di antara vektor-vektor tak nol dari
subruang-subruang $W_i$ yang berbeda.
Hubungan ini mungkin melibatkan atau tidak melibatkan suatu vektor dari $W_1$,
dari $W_2$, dan seterusnya, sehingga kita menuliskannya sebagai
$\zero=\mbox{}\cdots+d_i\vec{w}_i+\cdots\mbox{}$.
Seperti dalam persamaan~($*$), uraikan vektornya.
\begin{align*}
  \zero   
   &= \mbox{}\dots
      +d_i(c_{i,1}\vec{\beta}_{i,1}+\dots+c_{i,n_i}\vec{\beta}_{i,n_i})
      +\cdots\mbox{}                                               \\        
   &= \mbox{}\dots
      +d_ic_{i,1}\cdot\vec{\beta}_{i,1}
      +\dots+
      d_ic_{i,n_i}\cdot\vec{\beta}_{i,n_i}
      +\cdots\mbox{}                                                      
\end{align*}
Kebebasan linear $\cat{\cat{B_1}{\cdots}}{B_k}$ memberikan bahwa setiap
koefisien $d_ic_{i,j}$ bernilai nol.
Karena $\vec{w}_i$ merupakan vektor tak nol, sekurang-kurangnya satu
$c_{i,j}$ tak nol, sehingga $d_i$ bernilai nol.
Hal ini berlaku bagi setiap $d_i$, sehingga hubungan linear tersebut trivial.

Terakhir, untuk $\text{(3)}\implies\text{(1)}$, andaikan setiap hubungan linear
di antara vektor-vektor tak nol dari subruang-subruang $W_i$ yang berbeda
bersifat trivial.
Perhatikan dua penguraian suatu vektor,
$\vec{v}=\mbox{}\cdots+\vec{w}_i+\cdots\mbox{}$ 
dan $\vec{v}=\mbox{}\cdots+\vec{u}_j+\cdots\mbox{}$
dengan $\vec{w}_i\in W_i$ dan $\vec{u}_j\in W_j$.
Kurangkan satu penguraian dari yang lain untuk memperoleh hubungan linear
seperti berikut (jika tidak ada $\vec{u}_i$ atau $\vec{w}_j$, abaikan suku
tersebut).
\begin{equation*}
  \zero
   =\mbox{}\cdots+(\vec{w}_i-\vec{u}_i)
     +\cdots
     +(\vec{w}_j-\vec{u}_j)+\cdots\mbox{}
\end{equation*}
Asumsi kasus bahwa pernyataan~(3) berlaku menyiratkan bahwa setiap suku sama
dengan vektor nol,
 $\vec{w}_i-\vec{u}_i=\zero$.
Jadi, penguraiannya tunggal.
\end{proof}

\begin{definition}
Suatu koleksi subruang \( \set{W_1,\ldots, W_k} \) disebut
\definend{bebas}\index{subruang!kebebasan}
jika tidak ada vektor tak nol dari sebarang \( W_i \) yang merupakan kombinasi
linear vektor-vektor dari subruang-subruang lain
\( W_1,\dots, W_{i-1},W_{i+1},\dots, W_k \).
\end{definition}

\begin{definition}
Suatu ruang vektor \( V \) merupakan
\definend{jumlah langsung}\index{subruang!jumlah langsung}%
\index{jumlah langsung!definisi}
(atau \definend{jumlah langsung internal}\index{jumlah langsung internal})
dari subruang-subruangnya \( W_1,\dots, W_k \) jika
\( V=W_1+W_2+\dots +W_k \)
dan koleksi \( \set{W_1,\dots, W_k} \) bebas.
Kita menulis \( V=W_1\directsum W_2\directsum \cdots\directsum W_k \).
\end{definition}

\begin{example}
Prototipe kita berlaku:
$\Re^3=\text{sumbu-$x$}\directsum\text{sumbu-$y$}\directsum\text{sumbu-$z$}$.
\end{example}

\begin{example}
Ruang matriks \( \nbyn{2} \) merupakan jumlah langsung berikut.
\begin{equation*}
  \set{\begin{mat}
         a  &0  \\
         0  &d
       \end{mat}  \suchthat a,d\in\Re }
  \,\mbox{}\directsum\mbox{}\,
  \set{\begin{mat}
         0  &b  \\
         0  &0
       \end{mat}  \suchthat b\in\Re }
  \,\mbox{}\directsum\mbox{}\,
  \set{\begin{mat}
         0  &0  \\
         c  &0
       \end{mat}  \suchthat c\in\Re }
\end{equation*}
Ruang itu juga merupakan jumlah langsung subruang melalui banyak cara lain;
penguraian jumlah langsung tidak tunggal.
\end{example}

\begin{corollary} \label{cor:DirSumDimsAdd}
Dimensi suatu jumlah langsung adalah jumlah dimensi suku-sukunya.
\end{corollary}

\begin{proof}
Dalam \nearbylemma{le:UniqDecIffBasisDec}, jumlah vektor basis dalam konkatenasi
sama dengan jumlah banyaknya vektor dalam subbasis-subbasisnya.
\end{proof}

Kasus khusus dua subruang layak dibahas secara tersendiri.

\begin{definition}
Jika suatu ruang vektor merupakan jumlah langsung dari dua subruangnya, maka
kedua subruang itu disebut
\definend{komplementer}.\index{subruang!komplementer}%
\index{subruang komplementer}
\end{definition}

\begin{lemma}
\label{le:DirectSumTwoSp}
Suatu ruang vektor \( V \) merupakan
jumlah langsung\index{jumlah langsung!dua subruang}
dari dua subruangnya \( W_1 \) dan \( W_2 \) jika dan hanya jika \( V \)
merupakan jumlah keduanya, \( V=W_1+W_2 \), dan irisannya trivial,
\( W_1\intersection W_2=\set{\zero\,} \).
\end{lemma}

\begin{proof}
Pertama, andaikan $V=W_1\directsum W_2$.
Menurut definisi, $V$ merupakan jumlah keduanya, $V=W_1+ W_2$.
Untuk menunjukkan bahwa irisan keduanya trivial, misalkan $\vec{v}$ merupakan
vektor dari $W_1\intersection W_2$ dan perhatikan persamaan
$\vec{v}=\vec{v}$.
Ruas kiri persamaan itu merupakan anggota $W_1$, sedangkan ruas kanannya
merupakan anggota $W_2$, yang dapat kita pandang sebagai kombinasi linear
anggota-anggota $W_2$.
Namun, kedua ruang tersebut bebas, sehingga satu-satunya cara suatu anggota
$W_1$ dapat menjadi kombinasi linear vektor-vektor dari $W_2$ ialah jika
anggota itu merupakan vektor nol, $\vec{v}=\zero$.

Untuk arah lainnya, andaikan $V$ merupakan jumlah dua ruang yang irisannya
trivial.
Untuk menunjukkan bahwa $V$ merupakan jumlah langsung keduanya, kita hanya
perlu menunjukkan bahwa kedua ruang itu bebas\Dash bahwa tidak ada anggota tak
nol dari ruang pertama yang dapat dinyatakan sebagai kombinasi linear
anggota-anggota ruang kedua, dan sebaliknya.
Hal ini berlaku karena setiap hubungan
$\vec{w}_1=c_1\vec{w}_{2,1}+\dots+c_k\vec{w}_{2,k}$ (dengan $\vec{w}_1\in W_1$
dan $\vec{w}_{2,j}\in W_2$ untuk semua $j$) menunjukkan bahwa vektor di ruas
kiri juga berada dalam $W_2$, sebab ruas kanan merupakan kombinasi
anggota-anggota $W_2$.
Irisan kedua ruang ini trivial, sehingga $\vec{w}_1=\zero$.
Argumen yang sama berlaku bagi sebarang $\vec{w}_2$.
\end{proof}

%\begin{corollary}
%\label{cor:CompsIffDimsAndIntTriv}
%Two subspaces \( W_1,W_2 \) are complements in \( V \) if and only if
%\( \dim(W_1)+\dim(W_2)=\dim(V) \) and
%\( W_1\intersection W_2=\set{\zero\,} \).
%\end{corollary}

%\begin{proof}
%This follows immediately from 
%\nearbycorollary{cor:DimSumEqDimSummandsMinDimInter}.
%\end{proof}

\begin{example}
Dalam $\Re^2$, sumbu-\( x \) dan sumbu-\( y \) saling komplementer, yakni
$\Re^2=\text{sumbu-$x$}\directsum\text{sumbu-$y$}$.
Ini menunjukkan bahwa komplemen subruang sedikit berbeda dari komplemen
himpunan;~sumbu $x$ dan~$y$ bukan komplemen himpunan karena irisannya bukan
himpunan kosong.

Suatu ruang dapat memiliki lebih dari satu pasangan subruang komplementer;
pasangan lain bagi~$\Re^2$ adalah subruang-subruang berupa garis \( y=x \) dan
\( y=2x \).
\end{example}

\begin{example}
Dalam ruang \( F=\set{a\cos\theta+b\sin\theta\suchthat a,b\in\Re} \),
subruang \( W_1=\set{a\cos\theta\suchthat a\in\Re} \) dan
\( W_2=\set{b\sin\theta\suchthat b\in\Re} \) saling komplementer.
Contoh sebelumnya mencatat bahwa suatu ruang dapat diuraikan menjadi lebih dari
satu pasangan komplemen.
Selain itu, perhatikan bahwa $F$ dapat memiliki lebih dari satu pasangan
subruang komplementer dengan $W_1$ sebagai anggota pertama\Dash komplemen lain
dari \( W_1 \) adalah
\( W_3=\set{b\sin\theta+b\cos\theta \suchthat b\in\Re} \).
\end{example}

\begin{example}
Dalam \( \Re^3 \), bidang-\( xy \) dan bidang-\( yz \) tidak saling
komplementer; inilah pokok pembahasan setelah
\nearbyexample{exam:RThreeIsSumXYAndYZ}.
Salah satu komplemen bidang-\( xy \) adalah sumbu-\( z \).
% A complement of the \( yz \)-plane is the line through \( (1,1,1) \).
\end{example}

Berikut pertanyaan alami yang muncul dari \nearbylemma{le:DirectSumTwoSp}:~untuk
$k>2$, apakah jumlah biasa $V=W_1+\dots+W_k$ juga merupakan jumlah langsung
jika dan hanya jika irisan subruang-subruangnya trivial?

\begin{example}    \label{exam:DirSumThree}
Jika terdapat lebih dari dua subruang, irisan yang trivial tidak cukup untuk
menjamin penguraian tunggal (yakni, tidak cukup untuk memastikan bahwa
ruang-ruang tersebut bebas).
Dalam \( \Re^3 \), misalkan \( W_1 \) adalah sumbu-\( x \), \( W_2 \) adalah
sumbu-\( y \), dan $W_3$ adalah subruang berikut.
\begin{equation*}
  W_3=\set{\colvec{q \\ q \\ r} \suchthat q,r\in\Re}
\end{equation*}
Pemeriksaan bahwa \( \Re^3=W_1+W_2+W_3 \) mudah.
Irisan
\( W_1\intersection W_2\intersection W_3 \) 
trivial, tetapi penguraiannya tidak tunggal.
\begin{equation*}
  \colvec{x \\ y \\ z}
  =\colvec{0 \\ 0 \\ 0}
  +\colvec{0 \\ y-x \\ 0}
  +\colvec{x \\ x \\ z}
  =\colvec{x-y \\ 0 \\ 0}
  +\colvec{0 \\ 0 \\ 0}
  +\colvec{y \\ y \\ z}
\end{equation*}
(Contoh ini juga menunjukkan bahwa syarat berikut pun tidak cukup:~semua irisan
berpasangan dari subruang-subruang tersebut trivial.
Lihat \nearbyexercise{exer:ThreeSubsPairwseNonTriv}.)
\end{example}

Dalam subbagian ini, kita telah melihat dua cara memandang suatu ruang sebagai
sesuatu yang dibangun dari bagian-bagian penyusunnya.
Keduanya berguna; secara khusus, kita akan menggunakan definisi jumlah langsung
pada akhir Bab Lima.

%In looking at how this overlap of subspaces interacts with the 
%`$+$' operation, we first consider the simplest case, the two-subspace case.
%To understand the relationship between the two subspaces $U$ and $W$, their
%sum $U+W$, and their intersection $U\intersection W$, we can ask about the
%bases.  
%Note that a basis for the intersection $U\intersection W$ is a linearly 
%independent subset of $U$ and also of $W$.
%Thus, we can think to expand it to make bases for the larger sets.

%\begin{lemma} \label{le:BasesSumTwoSubs}
%Let \( U \) and \( W \) be subspaces of a vector space.
%If the sequence \( \sequence{\vec{\beta}_1,\dots,\vec{\beta}_k} \) 
%is a basis for
%\( U\intersection W \), and
%the sequence \( \sequence{\vec{\mu}_1,\dots,\vec{\mu}_j,
%   \vec{\beta}_1,\dots,\vec{\beta}_k} \)
%is a basis for \( U \), and
%\( \sequence{\vec{\beta}_1,\dots,\vec{\beta}_k,
%   \vec{\omega}_1,\dots,\vec{\omega}_p} \)
%is a basis for \( W \), then
%\begin{equation*}
%   \sequence{\vec{\mu}_1,\dots,\vec{\mu}_j,
%             \vec{\beta}_1,\dots,\vec{\beta}_k,
%             \vec{\omega}_1,\dots,\vec{\omega}_p}
%\end{equation*}
%is a basis for \( U+W \).
%\end{lemma}

%\begin{proof}
%We write $B_{U\intersection W}$ for the basis for $U\intersection W$,
%we write $B_U$ for the basis for $U$, we write $B_W$ for the basis for $W$,
%and we write $B_{U+W}$ for the basis under consideration.

%To see that that $B_{U+W}$ spans $U+W$, observe that
%any vector $c\vec{u}+d\vec{w}$ from $U+W$ can be written as a linear
%combination of the vectors in $B_{U+W}$, simply by expressing $\vec{u}$ in 
%terms of $B_U$ and expressing $\vec{w}$ in terms of $B_W$.

%We finish by showing that $B_{U+W}$ is linearly independent. 
%Consider
%\begin{equation*}
%  c_1\vec{\mu}_1+\dots+c_{j+1}\vec{\beta}_1+\dots+c_{j+k+p}\vec{\omega}_p=
%    \zero
%\end{equation*}
%which can be rewritten in this way.
%\begin{equation*}
%  c_1\vec{\mu}_1+\dots+c_j\vec{\mu}_j=
%       -c_{j+1}\vec{\beta}_1-\dots-c_{j+k+p}\vec{\omega}_p
%\end{equation*}
%Note that the left side sums to a vector in \( U \) while 
%right side sums to a vector in \( W \), and thus both sides sum to
%a member of $U\intersection W$.
%Since the left side is a member of $U\intersection W$, 
%it is expressible in terms of the members of $B_{U\intersection W}$, 
%which gives the combination of $\vec{\mu}$'s from the left side above
%as equal to a combination of $\vec{\beta}$'s. 
%But, the fact that 
%the basis $B_U$ is linearly independent shows that any such combination
%is trivial, and in particular, the coefficients $c_1$, \ldots,
%$c_j$ fro the left side above are all zero.
%Similarly, the coefficients of the \( \vec{\omega} \)'s are all zero.
%This leaves the above equation as a linear relationship among the 
%\( \vec{\beta} \)'s, 
%but $B_{U\intersection W}$ is linearly independent, 
%and therefore all of the coefficients of the $\vec{\beta}$'s are also zero.
%\end{proof}

%\begin{example}
%Label the axes in \( \Re^4 \) with \( x \), \( y \), \( z \), and \( w \).
%If \( U \) is \( xyz \)-space
%(that is, the subspace spanned by the \( x \), \( y \), and
%\( z \) axes), and \( W \) is \( xyw \)-space, then this
%\begin{equation*}
%  B_{U\intersection W}
%     =\sequence{\colvec{1 \\ 0 \\ 0 \\ 0},
%                \colvec{0 \\ 2 \\ 0 \\ 0} }
%\end{equation*}
%is a basis for \( U\intersection W \).
%It can be extended to these bases
%\begin{equation*}
%  B_U=\sequence{\colvec{1 \\ 1 \\ 1 \\ 0},
%            \colvec{1 \\ 0 \\ 0 \\ 0},
%            \colvec{0 \\ 2 \\ 0 \\ 0} }
%  \quad\text{and}\quad
%  B_W=\sequence{\colvec{1 \\ 0 \\ 0 \\ 0},
%            \colvec{0 \\ 2 \\ 0 \\ 0},
%            \colvec{1 \\ 0 \\ 0 \\ -1} }
%\end{equation*}
%for \( U \) and \( W \).
%Then this 
%\begin{equation*}
%  B_{U+W}
%   =\sequence{\colvec{1 \\ 1 \\ 1 \\ 0},
%            \colvec{1 \\ 0 \\ 0 \\ 0},
%            \colvec{0 \\ 2 \\ 0 \\ 0},
%            \colvec{1 \\ 0 \\ 0 \\ -1} }
%\end{equation*}
%is a basis for \( U+W=\Re^4 \).
%\end{example}

%\begin{corollary}
%\label{cor:DimSumEqDimSummandsMinDimInter}
%For any subspaces \( U \) and \( W \) of a vector space,
%\begin{equation*}
%  \dim(U+W)=\dim(U)+\dim(W)-\dim(U\intersection W).
%\end{equation*}
%\end{corollary}

%\begin{proof}
%Using the notation from the lemma, $\dim (U\intersection W)=k$,
%and $\dim(U)=j+k$, and $\dim(W)=k+p$, and $\dim (U+W)=j+k+p$. 
%\end{proof}

%\begin{example}
%In the prior example, $\dim (U\intersection W)=2$,
%and $\dim(U)=3$, and $\dim (W)=3$, and $\dim (U+W)=\dim(\Re^4)=4$,
%and so the equation becomes $4=3+3-2$.     
%\end{example}

%\begin{example}
%In \nearbyexample{exam:RThreeIsSumXYAndYZ} the equation gives
%$\dim(\Re^3)=\dim(\text{$xy$-plane})+\dim(\text{$yz$-plane})
%    -\dim(\text{$y$-axis})$ which is, of course, simply $3=2+2-1$.
%\end{example}

%The lemma gives us a way to ensure that every vector $\vec{v}\in U+W$ has a 
%unique decomposition as the sum of a vector from $U$ and a vector from $W$.
%Recall that vectors are uniquely represented in terms of a basis.
%If $B_{U\intersection W}$ is the empty basis (that is, if the intersection of
%$U$ and $W$ is trivial) then in the representation
%$\vec{v}=c_1\vec{\mu}_1+\dots+c_{j+p}\vec{\omega}_{j+p}$ we can take
%$\vec{u}$ to be $c_1\vec{\mu}_1+\dots+c_{j}\vec{\mu}_{j}$
%and take $\vec{w}$ to be 
%$c_{j+1}\vec{\omega}_{j+1}+\dots+c_{j+p}\vec{\omega}_{j+p}$
%to have a unique $\vec{v}=\vec{u}+\vec{w}$.

%\begin{definition}
%\label{def:DirectSumTwoSp}
%A vector space \( V \) is the 
%\definend{direct sum}\index{subspace!direct sum}%
%\index{direct sum!of two subspaces}
%of two of its subspaces \( W_1 \) and \( W_2 \), written 
%$V=W_1\directsum W_2$, if
%\( W_1+W_2=V \) and~\( W_1\intersection W_2=\set{\zero\,} \).
%The two subspaces are
%\definend{complements} %
%\index{subspace!complementary}\index{complementary subspaces}
%in \( V \).
%\end{definition}

%\begin{corollary}
%\label{cor:CompsIffDimsAndIntTriv}
%Two subspaces \( W_1,W_2 \) are complements in \( V \) if and only if
%\( \dim(W_1)+\dim(W_2)=\dim(V) \) and
%\( W_1\intersection W_2=\set{\zero\,} \).
%\end{corollary}

%\begin{proof}
%This follows immediately from 
%\nearbycorollary{cor:DimSumEqDimSummandsMinDimInter}.
%\end{proof}

%\begin{example}
%The \( x \)-axis and the \( y \)-axis are complements in \( \Re^2 \),
%that is, $\Re^2=\text{$x$-axis}\directsum\text{$y$-axis}$.
%A space can have more than one pair of complementary subspaces;
%another pair here are the subspaces consisting of the 
%lines \( y=x \) and \( y=2x \).
%\end{example}

%\begin{example}
%In the space \( F=\set{a\cos\theta+b\sin\theta\suchthat a,b\in\Re} \),
%the subspaces \( W_1=\set{a\cos\theta\suchthat a\in\Re} \) and
%\( W_2=\set{b\sin\theta\suchthat b\in\Re} \) are complements.
%In addition to the fact that a space like $F$ can have more than one pair of
%complementary subspaces, inside of the space a single subspace like $W_1$ can
%have more than one subspace that is a complement to it---another 
%complement to \( W_1 \) is
%\( W_3=\set{b\sin\theta+c\cos\theta \suchthat b,c\in\Re} \).
%\end{example}

%\begin{example}
%In \( \Re^3 \), the \( xy \)-plane and the \( yz \)-planes are not
%complements, because they have a nontrivial intersection.
%One complement of the \( xy \)-plane is the \( z \)-axis.
%A complement of the \( yz \)-plane is the line through \( (1,1,1) \).
%\end{example}

%\begin{lemma}
%For a space $V$ and two subspaces $W_1$, $W_2$, every vector $\vec{v}$ in the
%space has a decomposition $\vec{v}=\vec{u}+\vec{w}$ that is unique (where
%$\vec{u}\in U$ and $\vec{w}\in W$) if and only if the space is the direct sum
%of the subspaces.
%\end{lemma}

%\begin{proof}
%The `if' direction is handled as discussed before the definition.
%If $V=W_1\directsum W_2$ then the intersection of the two spaces is trivial
%and so \nearbylemma{le:BasesSumTwoSubs} shows that there is a basis for
%the sum $W_1+W_2=V$ consisting of a basis for $W_1$ with a basis for $W_2$
%adjoined.
%The representation of any $\vec{v}\in V$ with respect to this basis is unique,
%and we can regroup inside of this representation to have a unique expression
%$\vec{v}=\vec{w}_1+\vec{w}_2$.

%For the `only if' direction, we assume that decompositions exist and are
%unique.
%The fact that for every $\vec{v}\in V$ a decomposition exists shows that
%the space is the sum of the subspaces.
%To show that the intersection of the spaces is trivial,
%consider a member of the intersection \ldots
  
%\end{proof}

%We finish this subsection by going to the case of more than two subspaces.
%This is a bit trickier that the two-subspace case; the right notion of
%``do not overlap'' is perhaps not evident.
%That is, we are looking for a definition of a space as a direct sum of its
%subspaces where the first condition is that the space be the sum of those 
%subspaces, and the second condition is strong enough to give that each vector
%in the space has a unique decomposition into parts, one from each subspace.

%A natural first guess at the second condition, based on the material above,
%is to require that the spaces have a trivial intersection.
%The next example shows that this guess is, however, not enough to guarantee
%unique decomposition. 

%\begin{example}    \label{exam:DirSumThree}
%In \( \Re^3 \), let
%\( W_1 \) be the \( x \)-axis, let \( W_2 \) be the \( y \)-axis, and let
%\begin{equation*}
%  W_3=\set{\colvec{q \\ q \\ r} \suchthat q,r\in\Re}.
%\end{equation*}
%The check that \( \Re^3=W_1+W_2+W_3 \) is easy.
%The intersection of all three is trivial
%\( W_1\intersection W_2\intersection W_3=\set{\zero} \) but
%decompositions aren't unique.
%\begin{equation*}
%  \colvec{x \\ y \\ z}
%  =\colvec{0 \\ 0 \\ 0}
%  +\colvec{0 \\ y-x \\ 0}
%  +\colvec{x \\ x \\ z}
%  =\colvec{x-y \\ 0 \\ 0}
%  +\colvec{0 \\ 0 \\ 0}
%  +\colvec{y \\ y \\ z}
%\end{equation*}
%(This example also shows that this requirement is also not 
%enough:~that all pairwise intersections of the subspaces be trivial.
%See \nearbyexercise{exer:ThreeSubsPairwseNonTriv}.)
%\end{example}

%To ensure that each vector in the sum has a unique decomposition,
%we need a still stronger condition.

%\begin{definition}
%A collection of subspaces \( \set{W_1,\ldots, W_k} \) is
%{\em independent\/}\index{subspace!independence}
%if no nonzero vector from any \( W_i \) is a linear combination of
%vectors from the other subspaces \( W_1,\dots, W_{i-1},W_{i+1},\dots, W_k \).
%\end{definition}

%\begin{definition}
%A vector space \( V \) is the
%{\em direct sum\/}\index{subspace!direct sum}\index{direct sum!definition}
%(or {\em internal direct sum\/}) of its subspaces \( W_1,\dots, W_k \) if
%\( W_1+W_2+\dots +W_k=V \)
%and the collection \( \set{W_1,\dots, W_k} \) is independent.
%We write \( V=W_1\directsum W_2\directsum \dots\directsum W_k \).
%\end{definition}

%\begin{example}
%The benchmark example fits under this definition:
%$\Re^3=\text{$x$-axis}\directsum\text{$y$-axis}\directsum\text{$z$-axis}$ 
%\end{example}

%\begin{example}
%The space of \( \nbyn{2} \) matrices is this direct sum
%\begin{equation*}
%  \set{\begin{mat}
%         a  &0  \\
%         0  &d
%       \end{mat}  \suchthat a,d\in\Re }
%  \,\mbox{}\directsum\mbox{}\,
%  \set{\begin{mat}
%         0  &b  \\
%         0  &0
%       \end{mat}  \suchthat b\in\Re }
%  \,\mbox{}\directsum\mbox{}\,
%  \set{\begin{mat}
%         0  &0  \\
%         c  &0
%       \end{mat}  \suchthat c\in\Re }
%\end{equation*}
%(it is the direct sum of subspaces in many other ways also; as with the case
%of only two subspaces, direct sum decompositions into three or more 
%subspaces need not be unique).
%\end{example}

%Of course, following the work in this subsection, that definition will only be
%justified if we support it with a result saying that under those conditions,
%each vector in the summation has a unique decomposition into a sum of vectors,
%one from each summand.
%For that, 
%we will relate the decomposion a vector space into a direct sum
%to a decomposition of some basis for that space.
%We already know that uniqueness of decomposition relates to bases.

%\begin{definition}
%The sequence {\em concatenation\/}\index{concatenation} of
%\( B_1=\sequence{\vec{\beta}_{1,1},\dots,\vec{\beta}_{1,n_1}} \),
%\dots, \( B_k=\sequence{\vec{\beta}_{k,1},\dots,\vec{\beta}_{k,n_k}} \) is
%\(
% \cat{\cat{B_1}{\cdots}}{B_k}=
%  \sequence{\vec{\beta}_{1,1},\dots,\vec{\beta}_{1,n_1},
%            \vec{\beta}_{2,1},\dots,\vec{\beta}_{k,n_k} }.
%\)
%\end{definition}

%\begin{corollary}
%\label{cor:DirSumIffBasesCat}
%If \( V \) has subspaces \( W_1,\dots, W_k \)
%with bases \( B_1,\dots, B_k \)
%then \( W_1\directsum\dots\directsum W_k=V \) if and only if
%\( \cat{\cat{B_1}{\cdots}}{B_k} \) is a basis for \( V \).
%\end{corollary}

%\begin{proof}
%We will show that 
%the sum \( W_1+\dots+W_k \) equals \( V \) 
%if and only if \( \cat{\cat{B_1}{\cdots}}{B_k} \) spans \( V \), and that
%the set \( \set{W_1,\dots, W_k} \) is independent if and only if
%\( \cat{\cat{B_1}{\cdots}}{B_k} \) is linearly independent.
%These two halves prove the `if and only if' equivalence.

%First~(1).
%Left to right:
%assume the subspaces sum to \( V \) so any vector in \( V \) can be represented
%as \( \vec{v}=\vec{w}_1+\dots+\vec{w}_k \),
%rewrite each \( \vec{w}_i \) as a sum of basis elements and
%conclude \( \cat{\cat{B_1}{\cdots}}{B_k} \) spans \( V \).
%Right to left: if \( \cat{\cat{B_1}{\cdots}}{B_k} \) spans \( V \)
%then any \( \vec{v}\in V \) is a linear combination of vectors from
%\( \cat{\cat{B_1}{\cdots}}{B_k} \).
%Inside that combination commute to put vectors from \( B_1 \) first,
%followed from vectors from \( B_2 \), etc.
%Sum to first part to some \( \vec{w}_1\in W_1\), sum the second part to
%a \( \vec{w}_2\in W_2 \), etc.
%The result has the desired form.

%Now we prove item~(2).
%`If' is easy.
%For `only if' suppose the set of subspaces is independent and consider a linear
%relationship among vectors in
%\( \cat{\cat{B_1}{\cdots}}{B_k} \).
%Rewrite that relationship so all the vectors from \( B_1 \) are on one side and
%the other vectors are on the other side.
%This expresses some member of \( W_1 \) as a linear combination of members of
%the other subspaces, and independence of the set of subspaces then implies
%this linear combination of members of \( B_1 \) is \( \zero \).
%Thus, in this relationship, all the scalars associated with members of
%\( B_1 \) are \( 0 \).
%Similarly every scalar in the combination is \( 0 \).
%\end{proof}

%Hence, as mentioned above, each direct sum decomposition of a vector space is
%associated with a decomposition of one of that space's bases.

%\begin{corollary}
%\label{cor:DirectSumIffUniqueSum}
%Where \( W_1,\dots, W_k \) are subspaces,
%\( V=W_1\directsum\dots\directsum W_k \) if and
%only if every vector \( \vec{v}\in V \) has one and only one representation
%\( \vec{v}=\vec{w}_1+\dots+\vec{w}_k \) where \( \vec{w}_i\in W_i \).
%\end{corollary}

%\begin{proof}
%A chain of double implications works:
%\( V=W_1\directsum\dots\directsum W_k \)
%if and only if
%\( \cat{\cat{B_1}{\cdots}}{B_k} \) is a basis for \( V \),
%which is true if and only if each vector in \( V \)
%has a unique representation as a sum of vectors from
%\( \cat{\cat{B_1}{\cdots}}{B_k} \).
%\end{proof}

%\begin{corollary}
%The dimension of a direct sum is the sum of the dimensions of its summands.
%\end{corollary}


\begin{exercises}
  \recommended \item
    Tentukan apakah \( \Re^2 \) merupakan jumlah langsung dari masing-masing
    pasangan \( W_1 \) dan \( W_2 \) berikut.
    \begin{exparts}
       \partsitem \( W_1=\set{\colvec{x \\ 0}\suchthat x\in\Re} \),
             \( W_2=\set{\colvec{x \\ x}\suchthat x\in\Re} \)
       \partsitem \( W_1=\set{\colvec{s \\ s}\suchthat s\in\Re} \),
             \( W_2=\set{\colvec{s \\ 1.1s}\suchthat s\in\Re} \)
       \partsitem \( W_1=\Re^2 \), \( W_2=\set{\zero} \)
       \partsitem \( W_1=W_2=\set{\colvec{t \\ t}\suchthat t\in\Re} \)
       \partsitem \( W_1=\set{\colvec[r]{1 \\ 0}+\colvec{x \\ 0}
                                \suchthat x\in\Re} \),
             \( W_2=\set{\colvec[r]{-1 \\ 0}+\colvec[r]{0 \\ y}\suchthat y\in\Re} \)
    \end{exparts}
    \begin{answer}
       Pada setiap kasus, kita dapat menerapkan
       \nearbylemma{le:DirectSumTwoSp}.
       \begin{exparts}
         \partsitem Ya.
           Bidang tersebut merupakan jumlah $W_1$ dan $W_2$ ini karena untuk
           sebarang skalar $a$ dan $b$
           \begin{equation*}
             \colvec{a \\ b}
             =\colvec{a-b \\ 0}
             +\colvec{b \\ b}
           \end{equation*}
           menunjukkan bahwa vektor umum merupakan jumlah vektor dari kedua
           bagian tersebut.
           Selain itu, kedua subruang ini merupakan garis (yang berbeda) melalui
           titik asal, sehingga irisannya trivial.
         \partsitem Ya.
           Untuk melihat bahwa setiap vektor dalam bidang merupakan kombinasi
           vektor dari bagian-bagian ini, perhatikan hubungan berikut.
           \begin{equation*}
             \colvec{a \\ b}=c_1\colvec[r]{1 \\ 1}+c_2\colvec[r]{1 \\ 1.1}
           \end{equation*}
           Kita dapat langsung memperhatikan bahwa himpunan
           \begin{equation*}
             \set{\colvec[r]{1 \\ 1},\colvec[r]{1 \\ 1.1}}
           \end{equation*}
           merupakan basis bagi ruang tersebut (karena jelas bebas linear dan
           berukuran dua dalam $\Re^2$), sehingga persamaan di atas memiliki
           tepat satu solusi; ini menyiratkan bahwa semua penguraian tunggal.
           Sebagai cara lain, kita dapat menyelesaikan
           \begin{equation*}
             \begin{linsys}{2}
               c_1  &+  &c_2    &=  &a  \\
               c_1  &+  &1.1c_2 &=  &b
             \end{linsys}
             \grstep{-\rho_1+\rho_2}
             \begin{linsys}{2}
               c_1  &+  &c_2    &=  &a  \\
                    &   &0.1c_2 &=  &-a+b
             \end{linsys}
           \end{equation*}
           untuk memperoleh $c_2=10(-a+b)$ dan $c_1=11a-10b$, sehingga
           \begin{equation*}
             \colvec{a \\ b}=\colvec{11a-10b \\ 11a-10b}
                             +\colvec{-10a+10b \\ 1.1\cdot(-10a+10b)}
           \end{equation*}
           sebagaimana diminta.
           Seperti pada jawaban sebelumnya, masing-masing dari kedua subruang
           merupakan garis melalui titik asal, dan irisannya trivial.
         \partsitem Ya.
           Setiap vektor dalam bidang merupakan jumlah dengan cara berikut:
           \begin{equation*}
             \colvec{x \\ y}=\colvec{x \\ y}+\colvec[r]{0 \\ 0}
           \end{equation*}
           dan irisan kedua subruang tersebut trivial.
         \partsitem Tidak.
           Irisannya tidak trivial.
         \partsitem Tidak.
           Keduanya bukan subruang.
      \end{exparts}  
    \end{answer}
  \recommended \item
    Tunjukkan bahwa \( \Re^3 \) merupakan jumlah langsung bidang-\( xy \)
    dengan masing-masing subruang berikut.
    \begin{exparts}
      \partsitem sumbu-\( z \)
      \partsitem garis
        \begin{equation*}
          \set{\colvec{z \\ z \\ z} \suchthat z\in\Re }
        \end{equation*}
    \end{exparts}
    \begin{answer}
      Pada setiap kasus, kita dapat menggunakan
      \nearbylemma{le:DirectSumTwoSp}.
      \begin{exparts}
        \partsitem Setiap vektor dalam \( \Re^3 \) dapat diuraikan sebagai
          jumlah berikut.
          \begin{equation*}
            \colvec{x \\ y \\ z}
            =\colvec{x \\ y \\ 0}
            +\colvec{0 \\ 0 \\ z}
          \end{equation*}
          Selain itu, irisan bidang-$xy$ dan sumbu-$z$ merupakan subruang
          trivial.
        \partsitem Setiap vektor dalam \( \Re^3 \) dapat diuraikan sebagai
          \begin{equation*}
            \colvec{x \\ y \\ z}
            =\colvec{x-z \\ y-z \\ 0}
            +\colvec{z \\ z \\ z}
          \end{equation*}
          dan irisan kedua ruang tersebut trivial.
      \end{exparts}  
    \end{answer}
  \item 
    Apakah \( \polyspace_2 \) merupakan jumlah langsung dari
    \( \set{a+bx^2\suchthat a,b\in\Re} \) dan
    \( \set{cx\suchthat c\in\Re} \)?
    \begin{answer}
      Ya.
      Menunjukkan bahwa keduanya merupakan subruang bersifat rutin.
      Untuk melihat bahwa ruang tersebut merupakan jumlah langsung keduanya,
      cukup perhatikan bahwa setiap anggota \( \polyspace_2 \) memiliki
      penguraian tunggal
      \( m+nx+px^2=(m+px^2)+(nx) \).  
     \end{answer}
  \recommended \item 
    Dalam \( \polyspace_n \), polinomial
    \definend{genap}\index{fungsi genap} merupakan anggota himpunan berikut:
    \begin{equation*}
      \mathcal{E}=
      \set{p\in\polyspace_n \suchthat \text{\( p(-x)=p(x) \) untuk semua \( x \)}}
    \end{equation*}
    dan polinomial \definend{ganjil}\index{fungsi ganjil}
    merupakan anggota himpunan berikut.
    \begin{equation*}
      \mathcal{O}=
      \set{p\in\polyspace_n \suchthat \text{\( p(-x)=-p(x) \)
                                              untuk semua \( x \)}}
    \end{equation*}
    Tunjukkan bahwa keduanya merupakan subruang komplementer.
    \begin{answer}
      Menunjukkan bahwa keduanya merupakan subruang bersifat rutin.
      Kita akan membuktikan bahwa keduanya komplementer dengan
      \nearbylemma{le:DirectSumTwoSp}.
      Irisan \( \mathcal{E}\intersection\mathcal{O} \) trivial karena satu-satunya
      polinomial yang memenuhi kedua syarat $p(-x)=p(x)$ dan $p(-x)=-p(x)$
      adalah polinomial nol.
      Untuk melihat bahwa seluruh ruang merupakan jumlah kedua subruang,
      \( \mathcal{E}+\mathcal{O}=\polyspace_n \), 
      perhatikan bahwa polinomial \( p_0(x)=1 \), \( p_2(x)=x^2 \),
      \( p_4(x)=x^4 \), dan seterusnya berada dalam \( \mathcal{E} \), serta
      polinomial \( p_1(x)=x \), \( p_3(x)=x^3 \), dan seterusnya berada dalam
      \( \mathcal{O} \).
      Jadi, setiap anggota \( \polyspace_n \) merupakan kombinasi anggota
      \( \mathcal{E} \) dan \( \mathcal{O} \).
    \end{answer}
  \item 
    Manakah di antara subruang-subruang \( \Re^3 \) berikut
    \begin{center}
     \begin{tabular}{l}
      $W_1$: sumbu-\( x \),\quad
      $W_2$:~sumbu-\( y \),\quad
      $W_3$:~sumbu-\( z \),             \\
      $W_4$:~bidang \( x+y+z=0 \),\quad
      $W_5$:~bidang-\( yz \)
     \end{tabular}
    \end{center}
    yang dapat dikombinasikan untuk
    \begin{exparts*}
      \partsitem berjumlah \( \Re^3 \)?
      \partsitem berjumlah langsung \( \Re^3 \)?
    \end{exparts*}
    \begin{answer}
     Masing-masing hasil berikut adalah $\Re^3$.
     \begin{exparts}
      \partsitem 
        Daftar ini dipecah menjadi beberapa baris agar mudah dibaca.
        \begin{center}
          \begin{tabular}{l} 
            $W_1+W_2+W_3$, 
              $W_1+W_2+W_3+W_4$, 
              $W_1+W_2+W_3+W_5$,  \\ \quad %line too long
              $W_1+W_2+W_3+W_4+W_5$, 
              $W_1+W_2+W_4$, 
              $W_1+W_2+W_4+W_5$,     \\ \quad %line too long
              $W_1+W_2+W_5$,         
              $W_1+W_3+W_4$, 
              $W_1+W_3+W_5$, 
              $W_1+W_3+W_4+W_5$,     \\ \quad 
              $W_1+W_4$, 
              $W_1+W_4+W_5$,         
              $W_1+W_5$,             \\ 
              $W_2+W_3+W_4$, 
              $W_2+W_3+W_4+W_5$,     
              $W_2+W_4$, 
              $W_2+W_4+W_5$,         \\ 
              $W_3+W_4$, 
              $W_3+W_4+W_5$,         \\ 
              $W_4+W_5$ 
          \end{tabular}
         \end{center}
     \partsitem
       $W_1\directsum W_2\directsum W_3$,    
       $W_1\directsum W_4$,    
       $W_1\directsum W_5$,    
       $W_2\directsum W_4$,    
       $W_3\directsum W_4$    
     \end{exparts}
    \end{answer}
  \recommended \item
    Tunjukkan bahwa \( \polyspace_n=\set{a_0 \suchthat a_0\in\Re}\directsum\dots
                              \directsum\set{a_nx^n\suchthat a_n\in\Re} \).
    \begin{answer}
      Jelas bahwa masing-masing merupakan subruang.
      Ketika dikonkatenasikan, basis \( B_i=\sequence{x^i} \) bagi
      subruang-subruang tersebut membentuk basis bagi seluruh ruang.
    \end{answer}
  \item 
    Apakah \( W_1+W_2 \) jika \( W_1\subseteq W_2 \)?
    \begin{answer}
       Hasilnya adalah \( W_2 \).
    \end{answer}
  \item 
    Apakah \nearbyexample{exam:BenchDirSum} dapat digeneralisasi?
    Dengan kata lain, benar atau salah:~jika suatu ruang vektor \( V \) memiliki
    basis
    \( \sequence{\vec{\beta}_1,\dots ,\vec{\beta}_n} \), maka ruang itu
    merupakan jumlah langsung rentang subruang-subruang
    berdimensi satu,
    \( V=\spanof{\set{\vec{\beta}_1}}\directsum\dots
           \directsum\spanof{\set{\vec{\beta}_n}} \)?
    \begin{answer}
      Benar menurut \nearbylemma{le:UniqDecIffBasisDec}.
    \end{answer}
  \item 
    Dapatkah \( \Re^4 \) diuraikan sebagai jumlah langsung melalui dua cara
    yang berbeda?
    Dapatkah \( \Re^1 \)?
    \begin{answer}
      Dua penguraian jumlah langsung \( \Re^4 \) yang berbeda mudah ditemukan.
      Dua di antaranya adalah
      \( W_1=\spanof{\set{\vec{e}_1,\vec{e}_2}} \) dan
      \( W_2=\spanof{\set{\vec{e}_3,\vec{e}_4}} \), serta
      \( U_1=\spanof{\set{\vec{e}_1}} \) dan
      \( U_2=\spanof{\set{\vec{e}_2,\vec{e}_3,\vec{e}_4}} \).
      (Masih banyak kemungkinan lain, misalnya \( \Re^4 \) dan subruang
      trivialnya.)

      Sebaliknya, setiap partisi basis vektor tunggal \( \Re^1 \) menghasilkan
      satu basis tanpa anggota dan satu basis dengan satu anggota.
      Jadi, setiap penguraian melibatkan \( \Re^1 \) dan subruang trivialnya.
    \end{answer}
  \item 
     Latihan ini membuat notasi `$+$' di antara himpunan terasa lebih alami.
     Buktikan bahwa jika \( W_1,\dots, W_k \) merupakan subruang dari suatu
     ruang vektor, maka
     \begin{equation*}
       W_1+\dots+W_k
       =\set{\vec{w}_1+\vec{w}_2+\dots+\vec{w}_k
             \suchthat \vec{w}_1\in W_1,\dots,\vec{w}_k\in W_k},
     \end{equation*}
     sehingga jumlah subruang merupakan subruang yang terdiri atas semua
     jumlah tersebut.
     \begin{answer}
       Inklusi himpunan ke satu arah mudah:
       \( \set{\vec{w}_1+\dots+\vec{w}_k\suchthat \vec{w}_i\in W_i} \)
       merupakan subhimpunan dari \( \spanof{W_1\union \dots \union W_k}  \)
       karena setiap \( \vec{w}_1+\dots+\vec{w}_k \) merupakan jumlah vektor
       dari gabungan tersebut.

       Untuk inklusi lainnya, pada sebarang kombinasi linear vektor-vektor dari
       gabungan, terapkan sifat komutatif penjumlahan vektor untuk menempatkan
       vektor-vektor dari \( W_1 \) terlebih dahulu, diikuti vektor-vektor dari
       \( W_2 \), dan seterusnya.
       Jumlahkan vektor-vektor dari \( W_1 \) untuk memperoleh
       \( \vec{w}_1\in W_1 \), jumlahkan vektor-vektor dari \( W_2 \) untuk
       memperoleh \( \vec{w}_2\in W_2 \), dan seterusnya.
       Hasilnya memiliki bentuk yang diinginkan.
     \end{answer}
  \item \label{exer:ThreeSubsPairwseNonTriv}
    (Lihat \nearbyexample{exam:DirSumThree}.
    Latihan ini menunjukkan bahwa syarat irisan berpasangan trivial memang lebih
    kuat daripada syarat yang hanya menyatakan bahwa irisan semua subruang
    trivial.)
    Berikan suatu ruang vektor dan tiga subruang $W_1$, $W_2$, dan $W_3$
    sedemikian sehingga ruang tersebut merupakan jumlah subruang-subruang itu,
    irisan ketiganya $W_1\intersection W_2\intersection W_3$ trivial, tetapi
    irisan berpasangannya
    $W_1\intersection W_2$, $W_1\intersection W_3$, dan $W_2\intersection W_3$
    tidak trivial.
    \begin{answer}
      Salah satu contoh diperoleh dengan mengambil ruang $\Re^3$ dan memilih
      bidang-$xy$, bidang-$xz$, dan bidang-$yz$ sebagai subruang-subruangnya.
    \end{answer}
  \item 
    Buktikan bahwa jika \( V=W_1\directsum\dots\directsum W_k \), maka
    \( W_i\intersection W_j \) trivial setiap kali \( i\neq j \).
    Ini menunjukkan bahwa separuh pertama bukti
    \nearbylemma{le:DirectSumTwoSp} meluas ke kasus lebih dari dua subruang.
    (\nearbyexample{exam:DirSumThree} menunjukkan bahwa implikasi ini tidak
    berlaku dalam arah sebaliknya; separuh lainnya tidak meluas.)
    \begin{answer}
      Tentu saja, vektor nol berada dalam semua subruang, sehingga irisannya
      sekurang-kurangnya memuat vektor tersebut.
      Menurut definisi jumlah langsung, himpunan \( \set{W_1,\dots,W_k} \)
      bebas, sehingga tidak ada vektor tak nol dari \( W_i \) yang merupakan
      kelipatan suatu anggota \( W_j \) ketika \( i\neq j \).
      Secara khusus, tidak ada vektor tak nol dari \( W_i \) yang sama dengan
      suatu anggota \( W_j \).
    \end{answer}
  \item 
    Ingat bahwa tidak ada himpunan bebas linear yang memuat vektor nol.
    Dapatkah suatu himpunan subruang yang bebas memuat subruang trivial?
    \begin{answer}
      Ya, himpunan itu dapat memuat subruang trivial;
      himpunan subruang dari \( \Re^3 \) berikut bebas:
      \( \set{\set{\zero},\text{sumbu-\( x \)}} \).
      Tidak ada vektor tak nol dari ruang trivial \( \set{\zero} \) yang
      merupakan kelipatan vektor dari sumbu-\( x \), sebab ruang trivial tidak
      memiliki vektor tak nol yang dapat menjadi calon kelipatan semacam itu
      (dan tidak ada pula vektor tak nol dari sumbu-\( x \) yang merupakan
      kelipatan vektor nol dari subruang trivial).
     \end{answer}
  \recommended \item 
    Apakah setiap subruang memiliki komplemen?
    \begin{answer}
      Ya.
      Untuk sebarang subruang dari suatu ruang vektor, kita dapat mengambil
      sebarang basis \( \sequence{\vec{\omega}_1,\dots,\vec{\omega}_k} \) bagi
      subruang itu dan memperluasnya menjadi basis
      \( \sequence{\vec{\omega}_1,\dots,\vec{\omega}_k,
                     \vec{\beta}_{k+1},\dots,\vec{\beta}_n} \) bagi seluruh
      ruang.
      Maka komplemen subruang semula memiliki basis berikut:
      \( \sequence{\vec{\beta}_{k+1},\dots,\vec{\beta}_n} \).  
     \end{answer}
  \recommended \item 
    Misalkan \( W_1, W_2 \) merupakan subruang dari suatu ruang vektor.
    \begin{exparts}
      \partsitem Andaikan himpunan \( S_1 \) merentang \( W_1 \), dan himpunan
        \( S_2 \) merentang \( W_2 \).
        Dapatkah \( S_1\union S_2 \) merentang \( W_1+W_2 \)?
        Haruskah demikian?
      \partsitem Andaikan \( S_1 \) merupakan subhimpunan bebas linear dari
        \( W_1 \), dan \( S_2 \) merupakan subhimpunan bebas linear dari
        \( W_2 \).
        Dapatkah \( S_1\union S_2 \) menjadi subhimpunan bebas linear dari
        \( W_1+W_2 \)?
        Haruskah demikian?
    \end{exparts}
    \begin{answer}
      \begin{exparts}
        \partsitem Harus.
          Setiap anggota \( W_1+W_2 \) dapat kita tulis sebagai
          \( \vec{w}_1+\vec{w}_2 \), dengan \( \vec{w}_1\in W_1 \) dan
          \( \vec{w}_2\in W_2 \).
          Karena \( S_1 \) merentang \( W_1 \), vektor \( \vec{w}_1 \)
          merupakan kombinasi anggota-anggota \( S_1 \).
          Demikian pula, \( \vec{w}_2 \) merupakan kombinasi anggota-anggota
          \( S_2 \).
        \partsitem Cara mudah untuk melihat bahwa gabungannya dapat bebas linear
          ialah mengambil keduanya sebagai himpunan kosong.
          Di sisi lain, dalam ruang $\Re^1$, jika \( W_1=\Re^1 \),
          \( W_2=\Re^1 \), $S_1=\set{1}$, dan $S_2=\set{2}$, maka gabungannya
          $S_1\union S_2$ tidak bebas.
      \end{exparts}  
     \end{answer}
  \item \label{exer:BasesSumTwoSubs}
    Ketika kita menguraikan suatu ruang vektor sebagai jumlah langsung,
    dimensi subruang-subruangnya berjumlah sama dengan dimensi ruang tersebut.
    Keadaan pada suatu ruang yang diberikan sebagai jumlah subruangnya tidak
    sesederhana itu.
    Latihan ini membahas kasus khusus dua subruang.
    \begin{exparts}
      \partsitem Untuk subruang-subruang \( \matspace_{\nbyn{2}} \) berikut,
        temukan
        \( W_1\intersection W_2 \), \( \dim(W_1\intersection W_2) \),
        \( W_1+W_2 \), dan \( \dim(W_1+W_2) \).
        \begin{equation*}
          W_1=\set{\begin{mat}
                    0  &0  \\
                    c  &d
                  \end{mat} \suchthat c,d\in\Re  }
          \qquad
         W_2=\set{\begin{mat}
                    0  &b  \\
                    c  &0
                  \end{mat} \suchthat b,c\in\Re  }
       \end{equation*}
     \partsitem Andaikan \( U \) dan \( W \) merupakan subruang dari suatu
       ruang vektor.
       Andaikan barisan
       \( \sequence{\vec{\beta}_1,\dots,\vec{\beta}_k} \) 
       merupakan basis bagi
       \( U\intersection W \).
       Terakhir, andaikan barisan sebelumnya telah diperluas sehingga
       menghasilkan barisan \( \sequence{\vec{\mu}_1,\dots,\vec{\mu}_j,
       \vec{\beta}_1,\dots,\vec{\beta}_k} \)
       yang merupakan basis bagi \( U \), serta barisan
       \( \sequence{\vec{\beta}_1,\dots,\vec{\beta}_k,
          \vec{\omega}_1,\dots,\vec{\omega}_p} \)
       yang merupakan basis bagi \( W \).
       Buktikan bahwa barisan berikut
       \begin{equation*}
         \sequence{\vec{\mu}_1,\dots,\vec{\mu}_j,
              \vec{\beta}_1,\dots,\vec{\beta}_k,
              \vec{\omega}_1,\dots,\vec{\omega}_p}
       \end{equation*}
       merupakan basis bagi jumlah \( U+W \).
      \partsitem Simpulkan bahwa
          $\dim (U+W)=\dim(U)+\dim(W)-\dim(U\intersection W)$.
      \partsitem Misalkan \( W_1 \) dan \( W_2 \) merupakan subruang
        berdimensi delapan dari suatu ruang berdimensi sepuluh.
        Daftarkan semua nilai yang mungkin bagi
        \( \dim(W_1\intersection W_2) \).
    \end{exparts}
    \begin{answer}
      \begin{exparts}
        \partsitem Irisan dan jumlahnya adalah
          \begin{equation*}
             \set{\begin{mat}
                    0  &0  \\
                    c  &0
                  \end{mat} \suchthat c\in\Re  }
             \qquad
             \set{\begin{mat}
                    0  &b  \\
                    c  &d
                   \end{mat} \suchthat b,c,d\in\Re  }
          \end{equation*}
        yang masing-masing berdimensi satu dan tiga.
      \partsitem Kita menulis $B_{U\intersection W}$ untuk basis bagi
        $U\intersection W$, $B_U$ untuk basis bagi $U$, $B_W$ untuk basis bagi
        $W$, dan $B_{U+W}$ untuk basis yang sedang dibahas.

        Untuk melihat bahwa $B_{U+W}$ merentang $U+W$, perhatikan bahwa setiap
        vektor $c\vec{u}+d\vec{w}$ dari $U+W$ dapat kita tulis sebagai kombinasi
        linear vektor-vektor dalam $B_{U+W}$, cukup dengan menyatakan
        $\vec{u}$ terhadap $B_U$ dan $\vec{w}$ terhadap $B_W$.

         Kita selesaikan bukti dengan menunjukkan bahwa $B_{U+W}$ bebas linear.
         Perhatikan
         \begin{equation*}
           c_1\vec{\mu}_1+\dots+c_{j+1}\vec{\beta}_1+\dots
              +c_{j+k+p}\vec{\omega}_p=\zero
         \end{equation*}
         yang dapat ditulis ulang sebagai berikut.
         \begin{equation*}
           c_1\vec{\mu}_1+\dots+c_j\vec{\mu}_j=
             -c_{j+1}\vec{\beta}_1-\dots-c_{j+k+p}\vec{\omega}_p
         \end{equation*}
         Perhatikan bahwa ruas kiri berjumlah menjadi vektor dalam \( U \),
         sedangkan ruas kanan berjumlah menjadi vektor dalam \( W \), sehingga
         kedua ruas merupakan anggota $U\intersection W$.
         Karena ruas kiri merupakan anggota $U\intersection W$, ruas itu dapat
         dinyatakan terhadap anggota-anggota $B_{U\intersection W}$; dengan
         demikian, kombinasi vektor-vektor $\vec{\mu}$ dari ruas kiri di atas
         sama dengan suatu kombinasi vektor-vektor $\vec{\beta}$.
         Namun, karena basis $B_U$ bebas linear, setiap kombinasi semacam itu
         trivial; khususnya, semua koefisien $c_1$, \ldots, $c_j$ dari ruas kiri
         di atas bernilai nol.
         Demikian pula, semua koefisien vektor-vektor \( \vec{\omega} \)
         bernilai nol.
         Persamaan di atas kini tinggal menjadi hubungan linear di antara
         vektor-vektor \( \vec{\beta} \), tetapi
         $B_{U\intersection W}$ bebas linear, sehingga semua koefisien
         vektor-vektor $\vec{\beta}$ juga bernilai nol.
        \partsitem Cukup hitung vektor basis dalam butir sebelumnya:
          $\dim(U+W)=j+k+p$, $\dim(U)=j+k$, $\dim(W)=k+p$, dan
          $\dim(U\intersection W)=k$.
        \partsitem Kita tahu bahwa
          \( \dim(W_1+W_2)=\dim(W_1)+\dim(W_2)-\dim(W_1\intersection W_2) \).
          Karena \( W_1\subseteq W_1+W_2 \), kita tahu bahwa dimensi
          \( W_1+W_2 \) sekurang-kurangnya sebesar dimensi $W_1$, sehingga
          dimensinya harus delapan, sembilan, atau sepuluh.
          Substitusi memberi tiga kemungkinan,
          $8=8+8-\dim(W_1\intersection W_2)$ atau
          $9=8+8-\dim(W_1\intersection W_2)$ atau
          $10=8+8-\dim(W_1\intersection W_2)$.
          Jadi, \( \dim(W_1\intersection W_2) \) harus bernilai delapan, tujuh,
          atau enam.
          (Mudah memberi contoh yang menunjukkan bahwa ketiga kasus tersebut
          mungkin, misalnya dalam \( \Re^{10} \).)
      \end{exparts}
    \end{answer}
  \item 
    Misalkan \( V=W_1\directsum\cdots\directsum W_k \), dan untuk setiap indeks
    \( i \), andaikan \( S_i \) merupakan subhimpunan bebas linear dari
    \( W_i \).
    Buktikan bahwa gabungan semua \( S_i \) bebas linear.
    \begin{answer}
      Perluas setiap \( S_i \) menjadi basis $B_i$ bagi \( W_i \).
      Konkatenasi basis-basis tersebut $\cat{\cat{B_1}{\cdots}}{B_k}$ merupakan
      basis bagi \( V \), sehingga anggota-anggotanya membentuk himpunan bebas
      linear.
      Namun, gabungan $S_1\union\cdots\union S_k$ merupakan subhimpunan dari
      himpunan bebas linear tersebut, sehingga gabungan itu sendiri bebas
      linear.
     \end{answer}
  \item 
    Suatu matriks disebut \definend{simetris}\index{matriks!simetris}%
    \index{matriks simetris}
    jika untuk setiap pasangan indeks \( i \) dan \( j \), entri \( i,j \)
    sama dengan entri \( j,i \).
    Suatu matriks disebut \definend{antisimetris}\index{matriks!antisimetris}%
    \index{matriks antisimetris}
    jika setiap entri \( i,j \) merupakan negatif dari entri \( j,i \).
    \begin{exparts}
      \partsitem Berikan matriks simetris $\nbyn{2}$ dan matriks antisimetris
        $\nbyn{2}$.
        (\textit{Catatan.}
        Untuk matriks kedua, perhatikan entri-entri pada diagonal.)
      \partsitem Apakah hubungan antara suatu matriks simetris persegi dan
        transposnya?
        Antara suatu matriks antisimetris persegi dan transposnya?
      \partsitem Tunjukkan bahwa \( \matspace_{\nbyn{n}} \) merupakan jumlah
        langsung dari ruang matriks simetris dan ruang matriks antisimetris.
    \end{exparts}
    \begin{answer}
      \begin{exparts}
        \partsitem Dua contoh yang diminta adalah
          \begin{equation*}
            \begin{mat}[r]
              1  &2  \\
              2  &3
            \end{mat}
            \qquad
            \begin{mat}[r]
              0  &1  \\
              -1 &0
            \end{mat}
          \end{equation*}
          Untuk matriks antisimetris, entri-entri diagonal harus nol.
        \partsitem Matriks simetris persegi sama dengan transposnya.
          Matriks antisimetris persegi sama dengan negatif transposnya.
        \partsitem Menunjukkan bahwa kedua himpunan merupakan subruang bersifat
          mudah.
         Andaikan \( A\in\matspace_{\nbyn{n}} \).
         Untuk menyatakan $A$ sebagai jumlah matriks simetris dan matriks
         antisimetris, perhatikan bahwa
         \begin{equation*}
            A=(1/2)(A+\trans{A}) + (1/2)(A-\trans{A})
         \end{equation*}
         dan suku pertama simetris, sedangkan suku kedua antisimetris.
         Jadi, \( \matspace_{\nbyn{n}} \) merupakan jumlah kedua subruang.
         Untuk menunjukkan bahwa jumlah itu langsung, andaikan suatu matriks
         \( A \) sekaligus simetris \( A=\trans{A} \) dan antisimetris
         \( A=-\trans{A} \).
         Maka \( A=-A \), sehingga semua entri \( A \) bernilai nol.
      \end{exparts}
     \end{answer}
  \item 
    Misalkan \( W_1,W_2,W_3 \) merupakan subruang dari suatu ruang vektor.
    Buktikan bahwa \( (W_1\intersection W_2)+(W_1\intersection W_3)\subseteq
              W_1\intersection (W_2+W_3) \).
    Apakah inklusi sebaliknya berlaku?
    \begin{answer}
      Andaikan
      \( \vec{v}\in (W_1\intersection W_2)+(W_1\intersection W_3) \).
      Maka \( \vec{v}=\vec{w}_2+\vec{w}_3 \), dengan
      \( \vec{w}_2\in W_1\intersection W_2 \) dan
      \( \vec{w}_3\in W_1\intersection W_3 \).
      Perhatikan bahwa \( \vec{w}_2,\vec{w}_3\in W_1 \), dan karena subruang
      tertutup terhadap penjumlahan,
      \( \vec{w}_2+\vec{w}_3\in W_1 \).
      Jadi, \( \vec{v}=\vec{w}_2+\vec{w}_3\in W_1\intersection(W_2+W_3) \).

      Contoh berikut membuktikan bahwa inklusinya dapat bersifat sejati:
      dalam \( \Re^2 \), ambil \( W_1 \) sebagai sumbu-\( x \), \( W_2 \)
      sebagai sumbu-\( y \), dan \( W_3 \) sebagai garis \( y=x \).
      Maka \( W_1\intersection W_2 \) dan \( W_1\intersection W_3 \) trivial,
      sehingga jumlah keduanya trivial.
      Namun, \( W_2+W_3 \) adalah seluruh \( \Re^2 \), sehingga
      \( W_1\intersection(W_2+W_3) \) adalah sumbu-\( x \).
     \end{answer}
  \item 
    Contoh sumbu-$x$ dan sumbu-$y$ dalam \( \Re^2 \) menunjukkan bahwa
    \( W_1\directsum W_2=V \) tidak menyiratkan \( W_1\union W_2=V \).
    Dapatkah \( W_1\directsum W_2=V \) dan \( W_1\union W_2=V \) terjadi
    bersamaan?
    \begin{answer}
      Hal itu terjadi ketika sekurang-kurangnya satu di antara \( W_1,W_2 \)
      trivial.
      Namun, itulah satu-satunya cara hal tersebut dapat terjadi.

      Untuk membuktikannya, andaikan keduanya tidak trivial, pilih vektor tak nol
      \( \vec{w}_1, \vec{w}_2 \) dari masing-masing subruang, lalu perhatikan
      \( \vec{w}_1+\vec{w}_2 \).
      Jumlah ini tidak berada dalam \( W_1 \), sebab
      \( \vec{w}_1+\vec{w}_2=\vec{v}\in W_1 \) akan menyiratkan bahwa
      \( \vec{w}_2=\vec{v}-\vec{w}_1 \) berada dalam \( W_1 \), yang
      bertentangan dengan asumsi kebebasan subruang-subruang tersebut.
      Demikian pula, \( \vec{w}_1+\vec{w}_2 \) tidak berada dalam \( W_2 \).
      Jadi, terdapat unsur \( V \) yang tidak berada dalam
      \( W_1\union W_2 \).
     \end{answer}
%   \recommended \item
%     Our model for complementary subspaces, the $x$-axis and the
%     $y$-axis in \( \Re^2 \), has one property not used here.
%     Where \( U \) is a subspace of \( \Re^n \) we define the
%     \definend{orthocomplement}\index{subspace!orthocomplement} 
%     of \( U \) to be
%     \begin{equation*}
%       U^\perp=
%       \set{\vec{v}\in\Re^n \suchthat \text{\( \vec{v}\dotprod\vec{u}=0 \)
%                                             for all \( \vec{u}\in U \)}}
%     \end{equation*}
%     (read ``\( U \) perp'').
%     \begin{exparts}
%       \partsitem Find the orthocomplement of the \( x \)-axis in \( \Re^2 \).
%       \partsitem Find the orthocomplement of the \( x \)-axis in \( \Re^3 \).
%       \partsitem Find the orthocomplement of the \( xy \)-plane in \( \Re^3 \).
%       \partsitem Show that the orthocomplement of a subspace is a subspace.
%       \partsitem Show that if \( W \) is the orthocomplement of \( U \) then
%         \( U \) is the orthocomplement of \( W \).
%       \partsitem Prove that a subspace and its orthocomplement have a trivial
%         intersection.
%       \partsitem Conclude that for any \( n \) and subspace 
%         \( U\subseteq \Re^n \) we have that \( \Re^n=U\directsum U^\perp \).
%       \partsitem Show that 
%         \( \dim(U)+\dim(U^\perp) \) equals the dimension of the
%         enclosing space.
%     \end{exparts}
%     \begin{answer}
%       \begin{exparts}
%         \partsitem The set 
%           \begin{equation*}
%             \set{\colvec{v_1 \\ v_2} 
%                \suchthat \colvec{v_1 \\ v_2}\dotprod\colvec{x \\ 0}=0 
%                               \text{\ for all $x\in\Re$}}
%           \end{equation*}
%           is easily seen to be the $y$-axis.
%         \partsitem The \( yz \)-plane.
%         \partsitem The \( z \)-axis.
%         \partsitem 
%           Assume that \( U \) is a subspace of some \( \Re^n \).
%           Because $U^\perp$ contains the zero vector, since that
%           vector is perpendicular to everything, we need only show that the
%           orthocomplement is closed under linear combinations of two elements.
%           If \( \vec{w}_1, \vec{w}_2\in U^\perp \) then
%           \( \vec{w}_1\dotprod\vec{u}=0 \) and \( \vec{w}_2\dotprod\vec{u}=0 \)
%           for all \( \vec{u}\in U \).
%           Thus \( (c_1\vec{w}_1+c_2\vec{w}_2)\dotprod\vec{u}=
%               c_1(\vec{w}_1\dotprod\vec{u})+c_2(\vec{w}_2\dotprod\vec{u})=0 \)
%           for all \( \vec{u}\in U \) and so \( U^\perp \) is closed under
%           linear combinations.
%         \partsitem The only vector orthogonal to itself is the zero vector.
%         \partsitem This is immediate.
%         \partsitem To prove that the dimensions add, 
%           it suffices by \nearbycorollary{cor:DirSumDimsAdd} and 
%           \nearbylemma{le:DirectSumTwoSp} to show that 
%           \( U\intersection U^\perp \) is the trivial subspace
%           \( \set{\zero} \).
%           But this is one of the prior items in this problem.
%       \end{exparts}  
%      \end{answer}
  \item 
    Perhatikan \nearbycorollary{cor:DirSumDimsAdd}.
    Apakah hasil itu berlaku dalam kedua arah\Dash yakni, dengan mengandaikan
    \( V=W_1+\dots+ W_k \), apakah
    \( V=W_1\directsum\cdots\directsum W_k \) jika dan hanya jika
    \( \dim(V)=\dim(W_1)+\dots+\dim(W_k) \)?
    \begin{answer}
      Ya.
      Implikasi dari kiri ke kanan adalah
      \nearbycorollary{cor:DirSumDimsAdd}.
      Untuk arah lainnya, andaikan
      \( \dim(V)=\dim(W_1)+\dots+\dim(W_k) \).
      Misalkan \( B_1,\dots, B_k \) merupakan basis bagi
      \( W_1,\dots, W_k \).
      Karena $V$ merupakan jumlah subruang-subruang itu, setiap
      \( \vec{v}\in V \) dapat ditulis sebagai
      \( \vec{v}=\vec{w}_1+\cdots+\vec{w}_k \); dengan menyatakan setiap
      \( \vec{w}_i \) sebagai kombinasi vektor-vektor dari basis terkait
      \( B_i \), terlihat bahwa konkatenasi
      \( \cat{B_1}{\cat{\cdots}{B_k}} \) merentang \( V \).
      Konkatenasi itu memiliki \( \dim(W_1)+\dots+\dim(W_k) \) anggota, sehingga
      merupakan himpunan perentang berukuran \( \dim(V) \).
      Oleh karena itu, konkatenasi tersebut merupakan basis bagi \( V \).
      Jadi, \( V \) merupakan jumlah langsung.
     \end{answer}
  \item 
    Kita tahu bahwa jika \( V=W_1\directsum W_2 \), maka terdapat basis bagi
    \( V \) yang terpecah menjadi basis bagi \( W_1 \) dan basis bagi
    \( W_2 \).
    Dapatkah kita membuat pernyataan yang lebih kuat bahwa setiap basis bagi
    \( V \) terpecah menjadi basis bagi \( W_1 \) dan basis bagi \( W_2 \)?
    \begin{answer}
      Tidak.
      Basis standar bagi \( \Re^2 \) tidak terpecah menjadi basis-basis bagi
      subruang komplementer berupa garis \( x=y \) dan garis \( x=-y \).
    \end{answer}
  \item 
     Kita dapat menanyakan sifat aljabar operasi `$+$'.
     \begin{exparts}
       \partsitem Apakah operasi itu komutatif; apakah
         \( W_1+W_2=W_2+W_1 \)?
       \partsitem Apakah operasi itu asosiatif; apakah
         \( (W_1+W_2)+W_3=W_1+(W_2+W_3) \)?
       \partsitem Misalkan \( W \) merupakan subruang dari suatu ruang vektor.
         Tunjukkan bahwa \( W+W=W \).
       \partsitem Haruskah terdapat unsur identitas, yakni subruang \( I \)
         sedemikian sehingga \( I+W=W+I=W \) untuk semua subruang \( W \)?
       \partsitem Apakah pencoretan kiri berlaku:~jika
         \( W_1+W_2=W_1+W_3 \), apakah \( W_2=W_3 \)?
         Bagaimana dengan pencoretan kanan?
    \end{exparts}
    \begin{answer}
      \begin{exparts}
        \partsitem Ya, \( W_1+W_2=W_2+W_1 \) untuk semua subruang \( W_1,W_2 \),
          sebab masing-masing ruas merupakan rentang dari
          \( W_1\union W_2=W_2\union W_1 \).
        \partsitem Hasil ini serupa dengan yang sebelumnya\Dash masing-masing
          ruas persamaan merupakan rentang dari
          \( (W_1\union W_2)\union W_3=W_1\union (W_2\union W_3) \).
        \partsitem Karena ini merupakan kesamaan antara himpunan, kita dapat
          menunjukkannya melalui inklusi timbal balik.
          Jelas bahwa \( W\subseteq W+W \).
          Untuk \( W+W\subseteq W \), ingat bahwa setiap subruang tertutup
          terhadap penjumlahan, sehingga setiap jumlah berbentuk
          \( \vec{w}_1+\vec{w}_2 \) berada dalam \( W \).
        \partsitem Dalam setiap ruang vektor, unsur identitas terhadap
          penjumlahan subruang adalah subruang trivial.
        \partsitem Pencoretan kiri maupun kanan tidak harus berlaku.
          Sebagai contoh, dalam \( \Re^3 \), ambil \( W_1 \) sebagai
          bidang-\( xy \), \( W_2 \) sebagai sumbu-\( x \), dan \( W_3 \)
          sebagai sumbu-\( y \).
      \end{exparts}
     \end{answer}
  % 2014-Dec-19 I commented this out to make the pages fit better; it is perfectly fine
  % \item 
  %   Consider the algebraic properties of the direct sum operation.
  %   \begin{exparts}
  %      \partsitem Does direct sum commute: does \( V=W_1\directsum W_2 \) imply
  %        that \( V=W_2\directsum W_1 \)?
  %     \partsitem Prove that direct sum is associative:
  %       \( (W_1\directsum W_2)\directsum W_3=
  %           W_1\directsum(W_2\directsum W_3) \).
  %      \partsitem Show that \( \Re^3 \) is the direct sum of the three axes
  %        (the relevance here is that by the previous item,
  %         we needn't specify which two of the three axes are combined first).
  %      \partsitem Does the direct sum operation left-cancel:~does
  %        \( W_1\directsum W_2=W_1\directsum W_3 \) imply \( W_2=W_3 \)?
  %        Does it right-cancel?
  %      \partsitem There is an identity element with respect to this operation.
  %        Find it.
  %      \partsitem Do some, or all, subspaces have inverses with respect to this
  %        operation:~is there a subspace \( W \) of some vector space such
  %        that there is a subspace \( U \) with the property that
  %        \( U\directsum W \) equals the identity element from
  %        the prior item?
  %   \end{exparts}
  %   \begin{answer}
  %     \begin{exparts}
  %        \partsitem They are equal because for each, 
  %          \( V \) is the direct sum if
  %          and only if we can write each \( \vec{v}\in V \) in a unique
  %          way as a sum \( \vec{v}=\vec{w}_1+\vec{w}_2 \) and
  %          \( \vec{v}=\vec{w}_2+\vec{w}_1 \).
  %        \partsitem They are equal because for each, 
  %          \( V \) is the direct sum if
  %          and only if we can write each \( \vec{v}\in V \) in a unique
  %          way as a sum of a vector from each
  %          $\vec{v}=(\vec{w}_1+\vec{w}_2)+\vec{w}_3$
  %          and $\vec{v}=\vec{w}_1+(\vec{w}_2+\vec{w}_3)$.
  %        \partsitem We can decompose any vector in \( \Re^3 \) uniquely into
  %          the sum of a vector from each axis.
  %        \partsitem No.
  %          For an example, in \( \Re^2 \) take \( W_1 \) to be the
  %          \( x \)-axis, take \( W_2 \) to be the \( y \)-axis, and
  %          take \( W_3 \) to be the line \( y=x \).
  %        \partsitem In any vector space the trivial subspace acts as 
  %          the identity element with respect to direct sum.
  %        \partsitem In any vector space, only the trivial subspace has
  %          a direct-sum inverse (namely, itself).
  %          One way to see this is that dimensions add, and so increase.
  %     \end{exparts}  
  %    \end{answer}
\end{exercises}
\index{jumlah langsung|)}
